\section{Lift estimates I: heights, cusps and defects}\label{sec:pf:lift1}

\providecommand{\liftonemainT}{\mathrm{mainT}}
\providecommand{\liftonemuc}{\mu_{\mathrm{core}}}
\providecommand{\liftonedfl}{d_{\mathrm{fl}}}
\providecommand{\liftonelogOm}{\log\Omega}
\providecommand{\liftonedefB}{\mathrm{defB}}
\providecommand{\liftoneflat}{\mathrm{fd}}
\providecommand{\liftonerd}{\mathrm{rd}}
\providecommand{\liftonecusp}{\mathcal S}
\providecommand{\liftonelat}{\Lambda}
\providecommand{\liftonehgt}{\mathrm{ht}}
\providecommand{\liftonedg}{\mathrm{dg}}
\providecommand{\liftonedegC}{\mathrm{deg}_{\mathrm{C}}}
\providecommand{\liftonerkC}{\mathrm{rk}_{\mathrm{C}}}
\providecommand{\liftonerk}{\operatorname{rk}}
\providecommand{\liftonevdm}{\mathrm{vdm}}
\providecommand{\liftoneerr}{\mathrm{err}}
\providecommand{\liftonebnd}{\mathrm{bnd}}
\providecommand{\liftonehtc}{\mathrm{htc}}
\providecommand{\liftonevol}{\mathrm{vol}}
\providecommand{\liftoneH}{\mathrm{Ht}}
\providecommand{\liftoneS}{\mathrm{Sh}}
\providecommand{\liftoneG}{\mathrm{Geo}}
\providecommand{\liftoneM}{\mathrm{Mlt}}
\providecommand{\liftoneB}{\mathrm{Bud}}
\providecommand{\liftoneP}{\mathrm{Pin}}
\providecommand{\liftoneR}{\mathrm{Shr}}
\providecommand{\liftoneHS}{\mathrm{HSB}}
\providecommand{\liftoneAnc}{\mathrm{Anc}}
\providecommand{\liftoneNw}{\mathcal N_{\mathrm{wt}}}
\providecommand{\liftoneNh}{\mathcal N_{\mathrm{host}}}
\providecommand{\liftoneDef}{\mathrm{Def}}
\providecommand{\liftoneFl}{\mathcal F}
\providecommand{\liftoneFib}{\mathrm{Fib}}
\providecommand{\liftoneeig}{\Lambda^{\chi}}
\providecommand{\liftoneHC}{\mathrm{HC}}

This section proves five results which enter the proof of the lift estimates
(\cref{thm:covolume:lift}). We use the notation of \cref{sec:covolume}: points $p$ of the admissible
families of \cref{sec:families}, indices $i\in\{0,\dots,m-1\}$ of the chain of lifts, the error unit
$U=\mathrm{unitE}(p)$, the main term $\liftonemainT=\liftonemainT_i(p)$, the core slope $\liftonemuc$,
the closedness quantifier $\mathrm{Closed}$ and deep covers $Y$ of level $i+1$
(\cref{sec:covolume:lift}). The results are the following.
\begin{enumerate}[label=(\Alph*)]
\item \emph{The low defect bound} (\cref{pf:lift1:lowdef}). At the top index $i+1=m$ one has
  $\liftonemuc\ge-\liftonemainT-C\,U$. At this index the unitary datum is definite, the moduli space is
  a finite set of CM points, and the core lattice is an eigenlattice of a finite group acting freely on
  these points. Its slope is bounded below by an explicit quantity $2\,\liftonedefB$ which is computed
  from the dual Weyl modules of the weight, the degree of the Lie algebra at the CM points (a period
  of Chowla--Selberg type) and the degree of the theta line; the anchored part of $2\,\liftonedefB$ is
  the main term, the rest is $O(U)$.
\item \emph{The height bound} (\cref{pf:lift1:height}). On a good deep cover, the heights of the
  connected components with respect to the automorphic line bundle of a weight $u_1$ are bounded below,
  in mean, in terms of the mean-slope bound $\sfname{DeepStR}$ (\cref{def:covolume:deepstr}). The bridge
  is an upper bound for the arithmetic volume of a component by its degree times its height.
\item \emph{The Hilbert--Samuel lower bound} (\cref{pf:lift1:hs}). The arithmetic degree of the
  lattice of cusp forms of the anchored weight $u_M$ on a component is bounded below by the degree and
  height of the component, through a Hilbert--Samuel estimate.
\item \emph{A cusp-regular level determinant} (\cref{pf:lift1:cusp}). On a good deep cover there is a
  cusp form of weight $c_0$ times the Hodge weight whose determinant along a group of level
  automorphisms is a non-zero-divisor in the ring of sections.
\item \emph{The Hodge determinant dominates} (\cref{pf:lift1:hodge}). On a toroidal model of a unitary
  moduli space, every line bundle is dominated at interior points by a power of the Hodge
  determinant $\det\Omega$.
\end{enumerate}
(A), (B) and (C) are three of the conditions of \cref{tab:covolume:conditions} (the low defect, height
and cover height sum conditions); (D) and (E) are geometric statements used elsewhere in the proof of
\cref{thm:covolume:lift}, and (E) is also used in the proof of (D). \Cref{pf:lift1:setting} fixes
notation.

In all five cases the formal proof is long (between $12{,}000$ and $32{,}000$ lines each) and its
weight lies in a few geometric or arithmetic estimates. We give complete proofs of the reductions to
these estimates and of the arguments that are specific to this part of the proof; the remaining
estimates are stated precisely and cited. The intermediate results proved in this way include the
reference weight and carrying section conditions of the height bound, the degree of flat lines at CM
points, an abstract arithmetic Hilbert--Samuel estimate (with its slope bound and a Bertini theorem over
number fields), the geometric inputs of the separation of orbits by cusp forms, and invariant smooth
refinements of cone decompositions.

% ---------------------------------------------------------------------------
\subsection{Setting}\label{pf:lift1:setting}

\subsubsection*{Lattices}
Let $K$ be a totally complex number field. For a finite-dimensional $K$-vector space $V$ with a family
of norms $\nu=(\nu_\sigma)_{\sigma\colon K\to\C}$ having a Gram frame and a finitely generated
$\OO_K$-submodule $L\ne0$ we use $\liftonerk L=\dim_K KL$ and the slope $\mu_\nu(L)$ of
\cref{def:covolume:slopes}; the \emph{degree} is $\deg_\nu L=\tfrac12\liftonerk L\cdot\mu_\nu(L)$. For a
finite set $\iota$ and a family $Y=(Y_\sigma)_\sigma$ of invertible complex $\iota\times\iota$
matrices, $\|x\|_{Y,\sigma}=|Y_\sigma\,\sigma(x)|$ (Euclidean norm of the vector
$(\sum_jY_{\sigma,ij}\sigma(x_j))_i$) is a family of norms on $K^\iota$ with a Gram frame. A
\emph{lattice} in a $K$-vector space is a finitely generated $\OO_K$-submodule spanning it. We use two
facts.

\begin{lemma}\label{pf:lift1:facts}
\begin{enumerate}[label=(\alph*)]
\item Let $L_o\subseteq V_o$ ($o\in O$, $O$ finite) be nonzero finitely generated submodules of spaces
  with norm families $\nu^o$ having Gram frames, and give $\bigoplus_oV_o$ the norms
  $\nu_\sigma(x)^2=\sum_o\nu^o_\sigma(x_o)^2$. Then $\liftonerk\bigoplus_oL_o=\sum_o\liftonerk L_o$ and
  $\deg_\nu\bigoplus_oL_o=\sum_o\deg_{\nu^o}L_o$.
\item Let $\Phi\colon V\to W$ be injective and $K$-linear, $\nu,\nu'$ norm families with Gram frames on
  $V,W$ with $\nu'_\sigma(\Phi x)\le\nu_\sigma(x)$ for all $\sigma,x$, and $A\subseteq V$,
  $B\subseteq W$ nonzero finitely generated with $\Phi(A)\subseteq B$ and $\liftonerk A=\liftonerk B$.
  Then $\mu_\nu(A)\le\mu_{\nu'}(B)$.
\end{enumerate}
\end{lemma}

\begin{proof}
(a) The realisation of $\bigoplus_oL_o$ is the orthogonal sum of the realisations of the $L_o$, so its
covolume is the product of their covolumes and its $\Z$-rank the sum of their $\Z$-ranks; the claim
follows from the definition of $\mu$ (the additive constants in $\mu$ do not depend on the lattice).
(b) This is \cref{lem:covolume:slopes}(a) with $\liftonerk A=\liftonerk B$ and $h(\Phi)\le0$.
\end{proof}

\subsubsection*{Unitary data, weights and forms}
The objects of this subsection are those of the core data and deep covers of \cref{sec:covolume}; we
recall what is used. A \emph{host} for a CM field $E$ consists of a prime $d\equiv7\bmod16$, an element
$\mathrm{rt}\in\OO_E$ with $\mathrm{rt}^2=-d$, a complex number $\mathrm{rt}_0$ with
$\mathrm{rt}_0^2=-d$ and further data; its \emph{CM type} is
$\Psi=\{\tau\colon E\to\C:\tau(\mathrm{rt})=\mathrm{rt}_0\}$. A \emph{datum} $D_w$ is a finitely
generated projective $\OO_E$-module $\Lambda$ with a unimodular alternating $\Z$-bilinear form
$\langle\,,\rangle$ satisfying $\langle ax,y\rangle=\langle x,\bar ay\rangle$, of signature
$(r_+,r_-)$ ($r_++r_-=\liftonerk\Lambda$); the multiplicity of $\tau$ is $r_+$ if $\tau\in\Psi$ and
$r_-$ otherwise. With a level $l$ (an odd modulus $M\ge3$, an ideal $\mathfrak M$ with
$\mathfrak M\bar{\mathfrak M}=(M)$, and an integer $A\ge1$, the ``$A$-part'') and a coefficient datum
(a number field $F$ containing the reflex data) one has the unitary moduli space $X^o$ over $F$, its
set $\mathrm{Comp}$ of connected components and, for each embedding $\sigma\colon F\to\C$, a complex
uniformisation. A \emph{weight} is a function $u\colon\Psi\times\{0,\dots,r-1\}\to\Z$; $u$ defines an
automorphic line bundle on the flag space of $X^o$ and a weight $G=\mathrm{wt}(u)$ of automorphic
forms. To a weight $G$ the development attaches the $F$-space $\liftonecusp(G)$ of cusp forms (formalised
as rules on framed points), its lattice $\liftonelat(G)$ of forms integral at all finite places, and
Petersson-type norms $\mathrm{nrm}_\sigma(G)$; for a component $Z$, the subspace
$\liftonecusp_Z(G)\subseteq\liftonecusp(G)$ of forms supported on $Z$, its lattice
$\liftonelat_Z(G)$ and norms $\mathrm{nrm}_{Z,\sigma}(G)$. The \emph{Hodge weight} $h$ is $0$ on the
first $r_+$ entries and $-1$ on the others; $\liftonedfl=|\Psi|\binom r2$ is the dimension of the flag
variety. For weights $u$ with $\liftonevdm(u)=\prod_\tau\prod_{a<b}(u_{\tau,a}-u_{\tau,b})\ne0$ put
\begin{equation}\label{pf:lift1:eq:rho}
u'=u-3r\,h,\qquad\rho(u)=1-\frac{\liftonevdm(u')}{\liftonevdm(u)};
\end{equation}
so $u'$ adds $3r$ to the last $r_-$ entries of each $u_\tau$. For a component $Z$ and a weight $u$ the
development defines (with $G=\mathrm{wt}(u)$):
\begin{itemize}
\item $\liftonedg(u,Z)\in\N$, the degree of the image of the generic point of $Z$ under the map defined
  by a basis of $\liftonecusp_Z(G)$, a subvariety of dimension $\liftonedfl$ of projective space (the
  degree of its Chow form);
\item $\liftonehgt(u,Z)\in\R$, the height of this image with respect to the lattice $\liftonelat_Z(G)$
  and the norms $\mathrm{nrm}_Z(G)$ (the height of its Chow form; defined when the norms have a Gram
  frame and $\liftonedg(u,Z)>0$);
\item $\liftonerkC(u,Z)=\liftonerk\liftonelat_Z(G)$ and $\liftonedegC(u,Z)=\deg\liftonelat_Z(G)$ for the
  norms $\mathrm{nrm}_Z(G)$;
\item the \emph{top degree} $\mathrm{top}(Z,u)=v_0(Z)\,\liftonevdm(u)$, where $v_0(Z)\ge0$ is the
  normalised growth of $\dim\liftonecusp_Z(t\,\mathrm{wt}(u_{\mathrm{std}}))$ as $t\to\infty$ for the
  standard weight, and $\delta'(Z,u)=\mathrm{top}(Z,u')$; thus
  $\delta'(Z,u)=(1-\rho(u))\,\mathrm{top}(Z,u)$ when $\liftonevdm(u)\ne0$.
\end{itemize}

\subsubsection*{Points, deep covers and lift conditions}
Let $p$ be a point of a family at $(x,\iota)$ (\cref{sec:covolume:statement}) and $i<m$. We write $j=i+1$,
$D_j$ for the datum of the core datum of $p$ at level $j$, and $(r_+,r_-)$ for its signature; by the
numerical conditions below $(r_+,r_-)=(k+j-1,m-j)$ and $r_++r_-=r$. The weight $\kappa_i$ of $p$ gives
exponents $a_\kappa\colon\{\text{embeddings of }E\}\times\N\to\N$ and an even integer $n_\kappa$; we write
$G_\kappa=\mathrm{wt}(a_\kappa,n_\kappa)$ for the corresponding weight of forms and $\lambda$ for the
weight of $D_j$ given by
\[
\lambda_{\tau,a}=a_\kappa(\tau,a)+\tfrac{n_\kappa}2\ (a<r_+),\qquad
\lambda_{\tau,a}=-a_\kappa(\bar\tau,r_--1-(a-r_+))-\tfrac{n_\kappa}2\ (a\ge r_+).
\]
The \emph{anchor} is $\mathrm{anc}(\tau)=\lambda_{\tau,k-2}$ ($\tau\in\Psi$); the anchored weights are
$u_W=\lambda-\mathrm{anc}$ (subtracting $\mathrm{anc}(\tau)$ from every entry of $\lambda_\tau$) and
$u_M=u_W-c_0h$, where $c_0=c_0(\mathcal D,j)$ is the least multiple of $8$ which is at least
$r+\binom{fr_+r_-+1}2+2$ and at least $8fr+4$; one has $0\le c_0\le22f^2r^4$. Further, $|\kappa_i|=\sum_{\tau,a}|\lambda_{\tau,a}|$, $L=x^\eps$ and $b=e^L$ are
size parameters of $p$, $\liftonelogOm=\liftonelogOm(p)$ is the constant of
\cref{def:covolume:mainT}, $\liftonerd(E)=|d_E|^{1/[E:\Q]}$, and $PP$ is a parameter of the family.
For $PP\in\R$ put
\[
T=\Bigl\lceil\frac{PP^2\liftonedfl^2k}{m}\Bigr\rceil,\qquad
B_T=\binom{T-1}{\liftonedfl},\qquad
\zeta(u_1)=7\Bigl(\frac{\liftonedfl^2}{T}+\liftonedfl\,\rho(u_1)\Bigr).
\]
\emph{Weight parameters.} Let $\kappa_{\mathrm{fl}}$ be the weight of the canonical bundle of the flag
variety and $T_1=T+\liftonedfl+1$. Dividing the successive differences of each
$(u_M-\kappa_{\mathrm{fl}})_\tau$ by $T_1$ (quotient and remainder) gives weights $u_1$ and
$u_{\mathrm{rem}}$, both vanishing at the entry $k-2$, the successive differences of $u_{\mathrm{rem}}$
lying in $[0,T+\liftonedfl]$; one puts $u_2=u_1+u_{\mathrm{rem}}$ and
$\mathrm{wt}_t=\kappa_{\mathrm{fl}}+u_2+(\liftonedfl+t)u_1$ for $0\le t\le T$. A weight $u_1$ is called
\emph{admissible} if it satisfies the divisibility relation of the development with respect to
$(u_W,c_0,k-1,T)$ (which expresses that $u_1$ is obtained by the division just described). At the
points considered, for $j<m$, every admissible weight equals the weight $u_1$ just constructed, the
weights $u_1$ and $u_1'$ are strictly decreasing in each $\tau$, and $u_M=\mathrm{wt}_T+w$ for a parallel
weight $w$ with $|w(\tau)|\le3r$ (\lref{GlobalTheory.T.weight_parameter_strict_antitone_alt}, a lift
condition proved from the size condition which states these properties).

The conditions of \cref{tab:covolume:conditions} are \emph{lift conditions}: for a predicate
$Q(p,i,PP,C)$, the lift condition defined by $Q$ is
\begin{multline*}
\mathrm{Closed}\Bigl(\mathrm{FamHyp},(F,C_H)\mapsto\exists C\ge0\ \exists x_0\ \forall x\ge x_0\
\forall\iota\ \forall p:\\
\mathrm{PointConditions}'(p)\wedge\mathrm{RecTied}(p)\Rightarrow\forall i:\
Q(p,i,PP_F(x,\iota),C)\Bigr),
\end{multline*}
where $PP_F(x,\iota)$ is the parameter $PP$ of the family $F$ at $(x,\iota)$ and $p$ runs over the points
of $F$ at $(x,\iota)$ (\cref{sec:covolume:statement}). For a predicate $Q'(p,i,Y,PP,C)$ on deep
covers, the \emph{cover condition} defined by $Q'$ is the lift condition defined by
$Q(p,i,PP,C)\colon i+1<m\Rightarrow\forall Y\,(\mathrm{GoodCov}(Y)\Rightarrow Q'(p,i,Y,PP,C))$.

\begin{lemma}[Combination]\label{pf:lift1:combine}
Let $Q_1,\dots,Q_s,Q$ be predicates as above, and suppose that there is a function
$\Phi\colon\R_{\ge0}^s\to\R_{\ge0}$, which may depend on finitely many reals not depending on $\eps$,
the family and $C_H$, such that for all $p,i,PP$ and $C_1,\dots,C_s\ge0$,
$Q_1(\dots,C_1)\wedge\dots\wedge Q_s(\dots,C_s)$ implies $Q(\dots,\Phi(C_1,\dots,C_s))$. If the lift
conditions defined by the $Q_l$ hold, so does the lift condition defined by $Q$. The same holds for cover
conditions.
\end{lemma}

\begin{proof}
$\mathrm{Closed}$ is stable under finite conjunctions and under implications between the inner
predicates (proof of \cref{thm:covolume:assembly}). Given $\eps$, the family and $C_H$, take the
constants $C_l$ and thresholds $x_l$ of the $Q_l$, the constant $\Phi(C_1,\dots,C_s)$ and the
threshold $\max_lx_l$.
\end{proof}

All proofs below that start from lift or cover conditions are applications of
\cref{pf:lift1:combine}; we only give the pointwise implications and the constants. We also use the
following numerical lift conditions, proved elsewhere.

\begin{definition}[Numerical conditions]\label{pf:lift1:num}
$\liftoneNw(p,i,PP)$ comprises: $(r_+,r_-)=(k+j-1,m-j)$ and $r_++r_-=r$; $|\Psi|=f$;
$\liftonedfl=f\binom r2$; $n_\kappa=n_V$; dominance of $(a_\kappa,n_\kappa)$;
$U=f\,m\,b\,L\log\log x$; $|\liftonelogOm|\le3L$; $|\lambda_{\tau,a}|\le4b$; size relations between
$PP$, $T$, $\liftonedfl$, $L$ and $b$; and, for $j<m$, the hypotheses of the weight construction
defining the admissible weights. $\liftoneNh(PP,A,M)$ comprises: $M\le L$, $10\le L\le x$,
$1\le\log\log x$; $8\le f$; $[E:\Q]=2f$; $\log(fr),\log n_V\le L$; $\log g_k\le9L+1$ for the auxiliary
primes $g_k$; $\log|d_E|\le4fL$; $\log d\le L$; $PP\le x^A$, $fr\,n_VL\le x^A$; $(x^A)^M\le b$; and
further size relations.
\end{definition}

The lift condition defined by $\liftoneNw$ (with any constant) holds by
\lref{GlobalTheory.CovTop.WtShape.wtN_z}, and for every $M$ the lift condition defined by
$\liftoneNh(PP,C,M)$ holds by \lref{GlobalTheory.CovTop.Sizes.hostN_alt}. In particular
$b\ge1$, $U\ge0$ and $L\ge10$ at all points considered.

\subsubsection*{Inputs}
The results of this section are proved from the following statements of the literature, formalised
in the development and listed in \cref{tab:covolume:inputs} (some through packages of named
statements); we use the abbreviations:
(Hdg) the Hodge package of \cref{tab:covolume:inputs};
(GM) Grothendieck--Messing deformation theory \cite{Messing1972};
(Sieg) toroidal compactifications of Siegel moduli spaces \cite{FaltingsChai1990};
(Tor) smooth toroidal models of unitary moduli spaces with the logarithmic Kodaira--Spencer isomorphism
\cite{Lan2013};
(Lie) seminegativity of subbundles of $\Lie G$ for semi-abelian $G$ over curves
\cite{FaltingsChai1990,MoretBailly1985};
(X) the package consisting of (Tor), (Lie), Green's theorem on cuspidal characters, point separation for
pairs \cite{AngehrnSiu1995} and a Fubini--Study/Chow height formula;
(H) the package consisting of (Sieg), toroidal charts, toroidal charts for fans and the compatibilities
of arithmetic degrees \cite{GilletSoule1990};
(Q) the package consisting of Kawamata--Viehweg vanishing, $L^2$ extension for log canonical pairs,
Serre's GAGA for holomorphic sections, the numerics of nef divisors \cite{Lazarsfeld2004}, the existence
of Chow forms (integral ones and those of hyperplane sections) and lengths of proper intersections;
(V) the bound of the arithmetic self-intersection of a hermitian model of the image of a component by
$[K:\Q](\liftonedfl+1)$ times its degree times its (shifted) height, together with the injectivity of small
forms into small sections of such a model.
Further inputs are introduced where they are used.

% ---------------------------------------------------------------------------
\subsection{The low defect bound}\label{pf:lift1:lowdef}

Recall from \cref{def:covolume:predicates} the predicates $\mathsf{GramAt}$ and
$\mathsf{LowAt}(C)\colon-\liftonemainT-C\,U\le\liftonemuc$, and the anchor predicate
$\liftoneAnc(C)$: $|\liftonemainT-\sum_{\tau\in\Psi}(r_+-r_-)\,\mathrm{anc}(\tau)\liftonelogOm|\le
C\,U$, together with two size bounds involving $\mathrm{anc}$ (\cref{tab:covolume:conditions}). The
last component of $\mathsf{GramAt}$ says that the forms of weight $G_\kappa$ of the core datum at level
$j$ are \emph{good}: their norms have a Gram frame and their lattice $\liftonelat(G_\kappa)$ is
finitely generated and spans $\liftonecusp(G_\kappa)$. The core slope is
$\liftonemuc=\mu(\liftonelat_{\mathrm{type}})$, where
$\liftonelat_{\mathrm{type}}\subseteq\liftonelat(G_\kappa)$ is the lattice of forms of the type
prescribed by $p$ and $\mu$ is taken for the norms $\mathrm{nrm}(G_\kappa)$.

\begin{theorem}[Low defect bound]\label{pf:lift1:lowdef-thm}
Assume (Hdg), the Gram condition (the lift condition defined by
$\mathsf{GramAt}$) and the anchor condition (the lift condition defined by $\liftoneAnc(C)$).
Then the low defect condition holds: the lift condition defined by
\[
i+1=m\ \Longrightarrow\ -\liftonemainT_i(p)-C\,U\le\liftonemuc(p,i).
\]
\formallabel{GlobalTheory.T.lowDef_of_gram_of_anchor}
\formallabel{GlobalTheory.CovTop.DefLow.covering_low_defect_alt}
\end{theorem}

\begin{remark}\label{pf:lift1:lowdef-rem}
Hypothesis (Hdg) is not used in the proof: the statements through which it is passed
(\cref{pf:lift1:defect,pf:lift1:defectlift}) do not use this hypothesis. Also, the case distinction in
the proof of \cref{pf:lift1:defectlift} has a vacuous branch.
\end{remark}

\subsubsection*{The definite datum and the defect bound}
At the top index $j=m$ the numerical conditions give $(r_+,r_-)=(k+m-1,0)=(r,0)$: the datum is definite.
For the weight exponents $a=a_\kappa$, $\tau\in\Psi$, let $a_\tau=(a(\tau,0),\dots,a(\tau,r_+-1))$,
which is non-increasing by dominance, $\mathrm{size}(a_\tau)=\sum_ia(\tau,i)$, and $\nabla(a_\tau)$ the
irreducible polynomial representation of $\GL_{r_+}(\C)$ of highest weight $a_\tau$ (the dual Weyl
module). For $c_1\in\R$ and a real number $\ell$ put
\begin{align}
\liftonedefB(a,n,c_1)&=\sum_{\tau\in\Psi}\Bigl(\Bigl(\frac{\mathrm{size}(a_\tau)}{r_+}+\frac n2\Bigr)c_1
-\frac12\log\dim\nabla(a_\tau)\Bigr)+n\,|\Psi|\,r_+\Bigl(\frac{\log2\pi}4-\frac{\log\liftonerd(E)}2\Bigr),
\label{pf:lift1:eq:defB}\\
\liftoneflat(\ell)&=\frac{r_+-r_-}2\,\ell+r_-\bigl(\log d-\log\liftonerd(E)\bigr).\label{pf:lift1:eq:fd}
\end{align}
The number $\liftoneflat(\liftonelogOm)$ is the degree of the flat sections (\cref{pf:lift1:flat}); at the
top index $\liftoneflat(\ell)=\frac r2\ell$.

\begin{lemma}\label{pf:lift1:twodef}
Let $D_w$ be a datum with $r_-=0$ and $r_+\ge1$, $a$ weight exponents, $n$ even, $k_c\in\N$ and
$\ell\in\R$; write $\lambda=\lambda(a,n)$ and $\mathrm{anc}(\tau)=\lambda_{\tau,k_c-1}$ (or $0$ if
$k_c-1\ge r$). Then
\begin{multline*}
-\sum_{\tau\in\Psi}r_+\mathrm{anc}(\tau)\,\ell-|\ell|\sum_{\tau\in\Psi}\sum_a|\lambda_{\tau,a}-\mathrm{anc}(\tau)|
-\sum_{\tau\in\Psi}\log\dim\nabla(a_\tau)\\
-n|\Psi|r_+|\log\liftonerd(E)|\le2\,\liftonedefB\bigl(a,n,-\liftoneflat(\ell)\bigr).
\end{multline*}
\formallabel{GlobalTheory.DefArith.two_defBound_ge}
\end{lemma}

\begin{proof}
Since $r_-=0$, all entries of $\lambda_\tau$ are $a(\tau,a)+n/2$, so
$\sum_a\lambda_{\tau,a}=\mathrm{size}(a_\tau)+r_+n/2$, and $\liftoneflat(\ell)=r_+\ell/2$. Hence the
$\tau$-term of $2\,\liftonedefB(a,n,-\liftoneflat(\ell))$ is
$-\bigl(\sum_a\lambda_{\tau,a}\bigr)\ell-\log\dim\nabla(a_\tau)$. Writing
$\sum_a\lambda_{\tau,a}=r_+\mathrm{anc}(\tau)+\sum_a(\lambda_{\tau,a}-\mathrm{anc}(\tau))$ and
$-\sum_a(\lambda_{\tau,a}-\mathrm{anc}(\tau))\ell\ge-|\ell|\sum_a|\lambda_{\tau,a}-\mathrm{anc}(\tau)|$
bounds the $\tau$-term below by the $\tau$-part of the left side. The last term of
$2\,\liftonedefB$ is $n|\Psi|r_+(\frac12\log2\pi-\log\liftonerd(E))\ge-n|\Psi|r_+|\log\liftonerd(E)|$,
as $\log2\pi>0$. Summing gives the claim.
\end{proof}

\begin{lemma}[Crude errors]\label{pf:lift1:crude}
The lift condition defined by the following predicate holds (the proof gives $C=35$): if $i+1=m$, then
\[
|\liftonelogOm|\sum_{\tau\in\Psi}\sum_a|\lambda_{\tau,a}-\mathrm{anc}(\tau)|
+\sum_{\tau\in\Psi}\log\dim\nabla(a_\tau)+n_\kappa|\Psi|r_+|\log\liftonerd(E)|\le C\,U .
\]
\formallabel{GlobalTheory.CovTop.DefSizes.parameters_scale_rank_gap_bounds}
\end{lemma}

\begin{proof}
By \cref{pf:lift1:combine} it suffices to prove the inequality with $C=35$ at points where
$\liftoneNw(p,i,PP)$, $\liftoneNh(PP,A,1)$ and the size relations of the family hold
(\cref{pf:lift1:num}). There, $|\liftonelogOm|\le3L$, $|\Psi|=f$, and the development proves the
per-$\tau$ bounds $\sum_a|\lambda_{\tau,a}-\mathrm{anc}(\tau)|\le10mb$ and
$\log\dim\nabla(a_\tau)\le3b$ (\lref{GlobalTheory.CovTop.DefSizes.sizes_at}; the second uses Weyl's
bound $\dim\nabla(a)\le(\sum_i|a_i-w|+1)^{\binom N2}$ for every $w\in\Z$,
\lref{GlobalTheory.T.dualWeyl_dim_formula_duality_bound}). Further
$n_\kappa|\Psi|r_+|\log\liftonerd(E)|\le n_Vfr\cdot2L\le2b$ by $|\log\liftonerd(E)|\le2L$ and the
host relations. As $U=fmbL\log\log x\ge fmbL$ and $f,m,L\ge1$, the sum is at most
$30fmbL+3fb+2b\le35U$.
\end{proof}

\subsubsection*{Definite moduli spaces}
Let $D_w$ be a datum, $l$ a level, $X^o$ the unitary moduli space over a coefficient field $F$ with
uniformisations $U_\sigma$, and $\xi>0$. The \emph{definite case} $\liftoneDef(U,\xi)$ means:
$r_-=0$, $r_+\ge1$; the level is good at every prime and has trivial $A$-part; $F$ is algebraically
closed in the function field of every component; the uniformisations at the good embeddings are unitary;
and the base points of the uniformisations have absolute value $\xi$ at every embedding of $E$. In this
case the components of $X^o$ are points, and for every component $Z$ the development constructs
\emph{rational point data} $P$
(\lref{GlobalTheory.T.nonempty_definiteRationalPointData_component}): an $F$-rational framed point
$y_0$ on $Z$ through which every point of $Z$ factors; for every finite place $w$ of $F$ an element
$a_w\in F^\times$ and matrices $B_{w}=(B_{w,\tau})_\tau$, $B_{w,\tau}\in\GL_{\mathrm{mult}(\tau)}(F)$,
describing the integral frames at $w$, with $a_w=1$ and $B_w=1$ for almost all $w$; and for every
$\sigma\colon F\to\C$ a real $N_\sigma>0$ and complex matrices $X_\sigma=(X_{\sigma,\tau})_\tau$
describing the metrics at $\sigma$. For a number field $K$, reals $(r_\sigma)_\sigma$ and
$(\varrho_w)_w$ (almost all zero) put
\[
\widehat{\deg}(r,\varrho)=\frac1{[K:\Q]}\Bigl(\sum_w\varrho_w\log\Nm w-\sum_\sigma\log r_\sigma\Bigr),
\]
the normalised Arakelov degree of a hermitian line with a generator of valuations $\varrho_w$ and norms
$r_\sigma$, and for matrices $\widehat{\deg}_{\det}(B,X)=\widehat{\deg}(|\det X_\sigma|,v_w(\det B_w))$. The
point data define the degrees of the Lie algebra and of the theta line,
\[
d^{\mathrm{Lie}}_P(\tau)=\widehat{\deg}_{\det}\bigl((B_{w,\tau})_w,(X_{\sigma,\tau})_\sigma\bigr),\qquad
d^{\theta}_P=\widehat{\deg}\bigl((N_\sigma^{-1/2})_\sigma,(v_w(a_w))_w\bigr).
\]

\begin{definition}\label{pf:lift1:flat}
For $\tau\in\Psi$ and a component $Z$, a \emph{flat section} on $Z$ for $\tau$ is a section of the
automorphic line of the flat weight attached to $\tau$ whose archimedean norms are constant on $Z$ and
whose valuations are prescribed by the flat defect of $\tau$; each flat section $s$ has a degree
$\deg s\in\R$. For $c_1\in\R$, $\liftoneFl(c_1)$ means: for every component $Z$ and every
$\tau\in\Psi$ flat sections on $Z$ for $\tau$ exist, and all of them have degree $-c_1$. A deep cover $Y$
has the \emph{flat degree property} if $\liftoneFl(-\liftoneflat(\liftonelogOm))$ holds for $Y$; the
\emph{flat degree condition} is the cover condition defined by this property.
\end{definition}

\begin{definition}\label{pf:lift1:fibreat}
For a component $Z$, a weight $G$ and $s\in\R$, $\liftoneFib(Z,G,s)$ means: there are a finite
nonempty set $\iota$, a lattice $L\subseteq F^\iota$, a family $Y=(Y_\sigma)$ of invertible complex
$\iota\times\iota$ matrices and an injective $F$-linear map $\varphi\colon F^\iota\to\{\text{rules}\}$
with image $\liftonecusp_Z(G)$, such that $s\le\frac12\mu_Y(L)$; $\varphi(x)$ is integral at every
finite place for $x\in L$; and $\|\varphi(x)\|_{z}\le\|x\|_{Y,\sigma}$ for every $\sigma$, $x$ and every
frame $z$ of $Z$ at $\sigma$, where $\|\cdot\|_z$ is the pointwise (homogeneous) norm of forms of weight
$G$ at $z$.
\end{definition}

\begin{proposition}[Fibres at CM points]\label{pf:lift1:fibre}
Assume $\liftoneDef(U,\xi)$, $(a,n)$ dominant and $\liftoneFl(c_1)$. Then
$\liftoneFib(Z,\mathrm{wt}(a,n),\liftonedefB(a,n,c_1))$ holds for every component $Z$.
\formallabel{GlobalTheory.DefFibre.fibreAt_defCase}
\end{proposition}

\begin{proof}
Fix $Z$ and rational point data $P$. The development proves:
\begin{enumerate}[label=(\roman*)]
\item\label{pf:lift1:fib1} there is an injective $F$-linear map
  $\varphi\colon F^{I}\to\{\text{rules}\}$, $I=\prod_\tau\mathrm{Idx}(a_\tau)$ the index set of the
  standard basis of $\bigotimes_\tau\nabla(a_\tau)$, with image $\liftonecusp_Z(\mathrm{wt}(a,n))$, such
  that the value of $\varphi(x)$ at a framed point obtained from $y_0$ by a change of frame with
  matrices $(Y_\tau)$ and scalar $c$ is $c^n$ times the pairing of $x$ with the matrix coefficients of
  $(Y_\tau)$ in $\bigotimes_\tau\nabla(a_\tau)$
  (\lref{GlobalTheory.T.exists_injective_dualWeylTensor_cuspForms_rationalPoint});
\item\label{pf:lift1:fib2} let $\rho_a(g)$ be the matrix of $g^{\mathsf T}$ acting on $\nabla(a)$ in its
  standard basis and $O_a$ the change of basis to an orthonormal basis; for the matrices
  $A_w=a_w^n\bigotimes_\tau\rho_{a_\tau}(B_{w,\tau})$ and
  $Y_\sigma=N_\sigma^{-n/2}\bigotimes_\tau O_{a_\tau}\rho_{a_\tau}(X_{\sigma,\tau})$, the module
  $L=\{x\in F^I:A_wx\text{ is integral at }w\text{ for all }w\}$ is a lattice, $\varphi(x)$ is integral
  at every finite place for $x\in L$, and $\|\varphi(x)\|_z\le\|x\|_{Y,\sigma}$ for every frame $z$ of $Z$
  at $\sigma$ (\lref{GlobalTheory.T.integral_coordinates_forms_and_fibre_norm_bound},
  \lref{GlobalTheory.T.fibre_action_preserves_adelic_matrices},
  \lref{GlobalTheory.T.isLattice_latticeOfLocalMatrixFamily});
\item\label{pf:lift1:fib3} the slope of this fibre lattice satisfies
  \[
  \widehat{\deg}\bigl((N_\sigma^{-n/2}),(v_w(a_w^n))\bigr)+\sum_{t}\Bigl(\frac{\mathrm{size}(a_t)}{\mathrm{mult}(t)}
  \widehat{\deg}_{\det}\bigl((B_{w,t}),(X_{\sigma,t})\bigr)-\frac12\log\dim\nabla(a_t)\Bigr)\le\frac12\mu_Y(L),
  \]
  the sum over all embeddings $t$ of $E$ (\lref{GlobalTheory.T.slope_dualWeylFibreLattice_ge});
\item\label{pf:lift1:fib4} $d^{\mathrm{Lie}}_P(\tau)=-\deg s$ for every $\tau\in\Psi$ and every flat
  section $s$ on $Z$ for $\tau$, and $\mathrm{mult}(t)=0$ for $t\notin\Psi$
  (\lref{GlobalTheory.T.degree_dualDetLie_eq_neg_flatSection_degree});
\item\label{pf:lift1:fib5} $\frac12\sum_{\tau\in\Psi}d^{\mathrm{Lie}}_P(\tau)+|\Psi|r_+\bigl(\frac14\log2\pi
  -\frac12\log\liftonerd(E)\bigr)\le d^{\theta}_P$
  (\lref{GlobalTheory.T.thetaLine_degree_ge_half_detLie_definitePoint}).
\end{enumerate}
By \ref{pf:lift1:fib4} and $\liftoneFl(c_1)$, $d^{\mathrm{Lie}}_P(\tau)=c_1$ for $\tau\in\Psi$.
For $t\notin\Psi$ the multiplicity is $0$, so $\nabla(a_t)$ is trivial and the $t$-term in
\ref{pf:lift1:fib3} vanishes (with the convention $x/0=0$ of the development); for $\tau\in\Psi$ the
$\tau$-term is
$\frac{\mathrm{size}(a_\tau)}{r_+}c_1-\frac12\log\dim\nabla(a_\tau)$. The first term of
\ref{pf:lift1:fib3} equals $n\,d^{\theta}_P$, and by \ref{pf:lift1:fib5} and $n\ge0$,
$n\,d^{\theta}_P\ge n|\Psi|\frac{c_1}2+n|\Psi|r_+(\frac14\log2\pi-\frac12\log\liftonerd(E))$.
Adding up, the left side of \ref{pf:lift1:fib3} is at least $\liftonedefB(a,n,c_1)$
by \eqref{pf:lift1:eq:defB}. Thus $(I,L,Y,\varphi)$ witnesses
$\liftoneFib(Z,\mathrm{wt}(a,n),\liftonedefB(a,n,c_1))$.
\end{proof}

\subsubsection*{Eigenlattices under free actions}

\begin{lemma}\label{pf:lift1:orbit}
Let $K$ be a field, $V$ a $K$-vector space, $G$ a finite group acting freely on a finite set $C$,
$(e_Z)_{Z\in C}$ $K$-linear endomorphisms of $V$ with $e_Z^2=e_Z$, $e_Ze_{Z'}=0$ for $Z\ne Z'$ and
$\sum_Ze_Z=1$, $\rho$ a representation of $G$ on $V$ with $\rho(g)e_Z=e_{gZ}\rho(g)$, and
$\chi\colon G\to K^\times$ a character. Let $O\subseteq C$ be a set of representatives of the
$G$-orbits. Then:
\begin{enumerate}[label=(\alph*)]
\item every $Z\in C$ is $g\cdot o$ for a unique pair $(g,o)\in G\times O$;
\item if $\rho(g)F=\chi(g)F$ for all $g$, then $F=\sum_{o\in O}\sum_{g\in G}\chi(g)^{-1}\rho(g)e_oF$;
\item if $f\colon C\to V$ satisfies $e_of_o=f_o$ for all $o$, then
  $\Phi(f)=\sum_{o\in O}\sum_{g\in G}\chi(g)^{-1}\rho(g)f_o$ satisfies $\rho(h)\Phi(f)=\chi(h)\Phi(f)$ for
  all $h$, and $e_{g\cdot o}\Phi(f)=\chi(g)^{-1}\rho(g)f_o$ for all $g\in G$, $o\in O$.
\end{enumerate}
\formallabel{GlobalTheory.freeOrbit_projector_eigen}
\end{lemma}

\begin{proof}
(a) Existence holds by the choice of $O$; if $g\cdot o=g'\cdot o'$ then $o,o'$ lie in the same orbit, so
$o=o'$, and $g'^{-1}g$ fixes $o$, so $g=g'$ by freeness. (b) By (a),
$F=\sum_Ze_ZF=\sum_{o}\sum_ge_{go}F$, and
$e_{go}F=e_{go}\chi(g)^{-1}\rho(g)F=\chi(g)^{-1}\rho(g)e_oF$. (c) Substituting $g\mapsto h^{-1}g$ gives
$\rho(h)\Phi(f)=\sum_o\sum_g\chi(h^{-1}g)^{-1}\rho(g)f_o=\chi(h)\Phi(f)$. Further
$e_{go}\rho(g')f_{o'}=\rho(g')e_{g'^{-1}go}e_{o'}f_{o'}$, which is $0$ unless $g'^{-1}g\cdot o=o'$, i.e.\
unless $(g',o')=(g,o)$ by (a); hence $e_{go}\Phi(f)=\chi(g)^{-1}\rho(g)f_o$.
\end{proof}

\begin{lemma}\label{pf:lift1:sumslope}
Let $K$ be totally complex, $V$ a $K$-vector space with a norm family $\nu$ having a Gram frame, and
$M\subseteq V$ finitely generated. Let $O$ be a finite nonempty set and, for $o\in O$, $\iota_o$ a finite
nonempty set, $Y_o$ a family of invertible complex $\iota_o\times\iota_o$ matrices and
$L_o\subseteq K^{\iota_o}$ a lattice. Let $\Phi\colon\bigoplus_oK^{\iota_o}\to V$ be injective with
$\Phi(\bigoplus_oL_o)\subseteq M\subseteq\Phi(\bigoplus_oK^{\iota_o})$ and
$\nu_\sigma(\Phi x)^2\le\sum_o\|x_o\|_{Y_o,\sigma}^2$ for all $\sigma,x$. If $s\le\mu_{Y_o}(L_o)$ for all
$o$, then $s\le\mu_\nu(M)$.
\formallabel{GlobalTheory.SumSlope.le_slope_sq_of_sum}
\end{lemma}

\begin{proof}
Give $W=\bigoplus_oK^{\iota_o}$ the norms $N_\sigma(x)^2=\sum_o\|x_o\|^2_{Y_o,\sigma}$ (block diagonal
frame) and let $L=\bigoplus_oL_o$, a lattice in $W$. By \cref{pf:lift1:facts}(a),
$\liftonerk L\cdot\mu_N(L)=\sum_o\liftonerk L_o\,\mu_{Y_o}(L_o)\ge s\liftonerk L$. As $L$ spans $W$,
$\Phi(L)\subseteq M\subseteq\Phi(W)$ gives $\liftonerk M=\dim W=\liftonerk L$. Since
$\nu_\sigma(\Phi x)\le N_\sigma(x)$, \cref{pf:lift1:facts}(b) gives $\mu_N(L)\le\mu_\nu(M)$.
\end{proof}

\begin{lemma}\label{pf:lift1:sloeig}
Let $K$ be totally complex, $V$ a $K$-vector space with a norm family $\nu\ge0$ having a Gram frame,
and $G,C,(e_Z),\rho,\chi$ as in \cref{pf:lift1:orbit} with $C\ne\emptyset$. Let $M_0\subseteq V$ be a
finitely generated $\OO_K$-submodule with $\rho(g)M_0\subseteq M_0$ and
$M=\{v\in M_0:\rho(g)v=\chi(g)v\ \forall g\}$. For $o\in C$ let $\iota_o,Y_o,L_o$ be as in
\cref{pf:lift1:sumslope} and $\varphi_o\colon K^{\iota_o}\to V$ injective with
$\varphi_o(K^{\iota_o})=\{v:e_ov=v\}$ and $\varphi_o(L_o)\subseteq M_0$. Assume the domination property:
for all $\sigma$, $v\in V$ and $B\in\R$, if for every $Z\in C$ there are $g,o,x,u$ with $Z=g\cdot o$,
$|\sigma(u)|=1$, $e_Zv=u\,\rho(g)\varphi_o(x)$ and $\|x\|_{Y_o,\sigma}\le B$, then $\nu_\sigma(v)\le B$.
If $s\le\mu_{Y_o}(L_o)$ for all $o$, then $M\ne0$ and $s\le\mu_\nu(M)$.
\formallabel{GlobalTheory.DefPt.OrbitGeo.slope_eig}
\end{lemma}

\begin{proof}
Choose $O$ as in \cref{pf:lift1:orbit} and define $\Phi\colon\bigoplus_{o\in O}K^{\iota_o}\to V$ by
\[
\Phi(x)=\sum_{o\in O}\sum_g\chi(g)^{-1}\rho(g)\varphi_o(x_o).
\]
By \cref{pf:lift1:orbit}(c) with $f_o=\varphi_o(x_o)$, $\Phi(x)$ is a $\chi$-eigenvector and
$e_{go}\Phi(x)=\chi(g)^{-1}\rho(g)\varphi_o(x_o)$; with $g=1$, $e_o\Phi(x)=\varphi_o(x_o)$, so $\Phi$ is
injective. Since $\chi(g)^{|G|}=1$, the values of $\chi$ are roots of unity: they lie in $\OO_K$ and
have absolute value $1$ at every embedding. Hence $\Phi(\bigoplus_oL_o)\subseteq M$. Choosing
$o_0\in O$ and $0\ne x\in L_{o_0}$ gives $0\ne\Phi(x)\in M$, so $M\ne0$. If $v\in M$, then $e_ov$ lies in
the image of $\varphi_o$, say $e_ov=\varphi_o(x_o)$, and $v=\Phi(x)$ by \cref{pf:lift1:orbit}(b); so
$M\subseteq\Phi(\bigoplus_oK^{\iota_o})$. $M$ is finitely generated, being a submodule of $M_0$. Finally,
for $x$ and $\sigma$ let $B=(\sum_{o\in O}\|x_o\|^2_{Y_o,\sigma})^{1/2}$: every $Z$ is $g\cdot o$ with
$o\in O$, and $e_Z\Phi(x)=\chi(g)^{-1}\rho(g)\varphi_o(x_o)$ with $\|x_o\|_{Y_o,\sigma}\le B$; the domination
property gives $\nu_\sigma(\Phi x)\le B$. \Cref{pf:lift1:sumslope} concludes.
\end{proof}

A \emph{level action} of a finite group $G$ on $X^o$ is a map $\mathrm{lev}$ from $G$ to pairs of
$\OO_E$-linear endomorphisms of $\Lambda$ which acts on framed points by change of level structure;
it then acts on rules. It is \emph{free} if $g=1$ whenever $g$ maps a framed point $y$ to a framed
point $y'$ over the same underlying point. For a character $\chi\colon G\to F^\times$ and a weight $G'$
the \emph{eigenlattice} is
$\liftoneeig(G')=\{f\in\liftonelat(G'):g\cdot f=\chi(g)f\ \forall g\in G\}$.

\begin{proposition}[Eigenlattice slope]\label{pf:lift1:eiglat}
Assume $\liftoneDef(U,\xi)$, that $\mathrm{Comp}$ is finite and nonempty, that the dictionaries
take values in the domains of the uniformisations, that $\mathrm{lev}$ is a free level action of the
finite group $G$, that $(a,n)$ is dominant and the forms of weight $\mathrm{wt}(a,n)$ are good. If
$\liftoneFib(Z,\mathrm{wt}(a,n),s)$ holds for every component $Z$, then
$\liftoneeig(\mathrm{wt}(a,n))\ne0$ and $s\le\frac12\mu(\liftoneeig(\mathrm{wt}(a,n)))$.
\formallabel{GlobalTheory.DefPt.OrbitGeo.eigLat_slope}
\end{proposition}

\begin{proof}
Apply \cref{pf:lift1:sloeig} with $K=F$ (totally complex, since complex conjugation does not fix any
embedding of the coefficient field), $V$ the space of rules, $C=\mathrm{Comp}$, $e_Z$ multiplication by
the idempotent of $Z$, $\rho$ the action of $G$ on rules through $\mathrm{lev}$,
$M_0=\liftonelat(\mathrm{wt}(a,n))$, $\nu=\mathrm{nrm}(\mathrm{wt}(a,n))$, and $\iota_o,L_o,Y_o,\varphi_o$
from $\liftoneFib$ (with $2s$ in place of $s$). The hypotheses are checked from the definitions
and the following facts proved in the development: $G$ acts on $\mathrm{Comp}$ compatibly with the
action on points (\lref{GlobalTheory.DefPt.OrbitGeo.exists_mulAction}); the action is free, since a
component of a definite moduli space has a rational point through which all its points factor
(\lref{GlobalTheory.FlatLine.DefinitePoint.main}, which states this) and the level action is free;
$g\cdot e_Z=e_{gZ}$ (\lref{GlobalTheory.DefPt.OrbitGeo.actLev_idem}); $G$ preserves integral forms
(\lref{GlobalTheory.DefPt.OrbitGeo.actLev_int}); the eigenlattice is the set of $\chi$-eigenvectors
in $M_0$ (\lref{GlobalTheory.DefPt.OrbitGeo.mem_eigLat}); the image of $\varphi_o$ is
$\liftonecusp_o=\{v:e_ov=v\}$; and the domination property holds because the norm $\nu_\sigma(v)$ is
bounded by the supremum over the frames $z$ of the pointwise norms of $v$
(\lref{GlobalTheory.DefPt.OrbitGeo.petAt_le}), and the pointwise norm of
$e_Zv=u\,g\cdot\varphi_o(x)$ at a frame of $Z=g\cdot o$ equals that of $\varphi_o(x)$ at a frame of $o$
(\lref{GlobalTheory.DefPt.OrbitGeo.homN_actLev}), which is $\le\|x\|_{Y_o,\sigma}$ by
$\liftoneFib$.
\end{proof}

\begin{corollary}\label{pf:lift1:eigdef}
Under the hypotheses of \cref{pf:lift1:eiglat} without the $\liftoneFib$ assumption, if
$\liftoneFl(c_1)$ holds then $\liftoneeig(\mathrm{wt}(a,n))\ne0$ and
$\liftonedefB(a,n,c_1)\le\frac12\mu(\liftoneeig(\mathrm{wt}(a,n)))$.
\formallabel{GlobalTheory.DefPt.eigLat_slope_definite}
\end{corollary}

\begin{proof}
Combine \cref{pf:lift1:fibre} and \cref{pf:lift1:eiglat} with $s=\liftonedefB(a,n,c_1)$.
\end{proof}

\subsubsection*{CM points and the degree of flat lines}
Let $E$ be a CM field with a host whose prime is $d$. The host is a \emph{host field} if $E$ is abelian over
$\Q$ and every prime $w$ of $E$ dividing the different satisfies $\bar w\ne w$ and lies over a prime
$\ell$ with either $\ell\ne d$ and $-d$ a square modulo $\ell$, or $\ell=d$ and $w^3\nmid\ell$. A Dirichlet
character $\chi$ modulo $d$ is a \emph{Kronecker character} if $\chi(n)=\bigl(\frac{-d}n\bigr)$ (Jacobi
symbol) for all odd $n$; put
$\liftonelogOm_E(\chi)=\gamma+\log(2\pi/\liftonerd(E))-L'(1,\chi)/L(1,\chi)\in\C$.

\begin{definition}[CM objects]\label{pf:lift1:cmobj}
Let $K$ be a number field in which $E$ splits, $\rho\in\OO_K$ with $\rho^2=-d$, $\sigma_1\colon K\to\C$
and $p,q\in\N$. A \emph{CM object} of signature $(p,q)$ over $K$ consists of polarised
abelian schemes $A_+$, $A_-$ (the \emph{letters}) and $B$ over $\OO_K$ with actions of $\OO_E$ of types $(1,0)$, $(0,1)$
and $(p,q)$ with respect to $\rho$, and $\OO_E$-linear homomorphisms $\varphi_j\colon A_\pm\to B$ for $j$
in $J=\{1,\dots,p\}\sqcup\{1,\dots,q\}$ (with source $A_+$ for the first $p$ indices and $A_-$ for the
others), such that every complex point of $B$ (via $\sigma_1$) is a sum of images of complex points
under the $\varphi_j$, and the pull-back of the Weil pairing of $B$ by $\varphi_j$ is the Weil pairing of
the source composed with the action of some $h_j\in\OO_E$. \emph{Gram matrices} are
$g_+\in M_p(\OO_E)$ and $g_-\in M_q(\OO_E)$ such that the pull-back of the Weil pairing of $B$ by
$(\varphi_a,\varphi_b)$ is the Weil pairing of $A_\pm$ composed with the action of $g_\pm(a,b)$; they exist
for every CM object (\lref{GlobalTheory.T.exists_gramMatrices_cmUnitaryShimuraObject}).

A \emph{generic fibre} of $B$ is an identification of $B_K$ with a polarised abelian variety over $K$
together with a basis of its Lie algebra; it identifies $\Lie B_K$ with $K^g$ and $\Lie B$ with an
$\OO_K$-lattice. For $\tau\colon\OO_E\to\OO_K$ the \emph{eigenlattice} $\liftonelat_B(\tau)$ is the
sublattice of $\Lie B$ on which $\OO_E$ acts through $\tau$, and $\bar\tau=\tau\circ\bar{\ }$. There is a
unique hermitian family $h$ on $K^g$ induced by the polarisation
(\lref{GlobalTheory.CM.LieMetric.exists_unique_isHerm}); $\deg_h$ denotes the Arakelov degree of
lattices for $h$.
\end{definition}

\begin{theorem}[Lie eigenlattices of definite CM objects]\label{pf:lift1:cmdiff}
Let the host be a host field, $\chi$ a Kronecker character modulo $d$, $X$ a CM object of signature
$(p,q)$ over $K$ with $q=0$, a generic fibre of $B$, $h$ the hermitian family induced by the
polarisation, and $\tau\colon\OO_E\to\OO_K$ with $\tau(\mathrm{rt})=\rho$. Then
\[
\deg_h\liftonelat_B(\tau)-\deg_h\liftonelat_B(\bar\tau)=\frac{p-q}2\,\Re\liftonelogOm_E(\chi).
\]
\formallabel{GlobalTheory.T.cmObject_definite_lieEigenLattice_degree_difference}
\formallabel{GlobalTheory.FlatCore.CMDiff0.cmDiff_definite}
\end{theorem}

\begin{proof}
As $q=0$, $J=\{1,\dots,p\}$ and all $\varphi_j$ have source $A_+$. Choose generic fibres of $A_\pm$, and
let $M_j\colon\Lie A_{+,K}\to\Lie B_K$ be the differential of $\varphi_j$ in the chosen coordinates
(\lref{GlobalTheory.CM.DiffCoord.existsUnique_isDiffOf}, existence and uniqueness of differentials).
Put $\liftonelat'(\tau)=\sum_jM_j(\liftonelat_{A_+}(\tau))$, the lattice spanned by the images of the
$\tau$-eigenlattices of the letters, and let $\liftonelat'(\bar\tau)=0$ (the empty sum over the indices of
$A_-$).

\emph{Ranks.} Since $\rho^2=-d\ne0$, $\bar\tau(\mathrm{rt})=-\rho\ne\rho$. The eigenspace decomposition of
$\Lie B$ for an action of type $(p,0)$ (\lref{GlobalTheory.CM.EigDec.eigen_decomposition}) gives
$\liftonerk\liftonelat_B(\tau)=p$ and $\liftonerk\liftonelat_B(\bar\tau)=0$. The images
$M_j(\liftonelat_{A_+}(\tau))$ lie in $\liftonelat_B(\tau)$, have rank $1$ each and are independent
(\lref{GlobalTheory.CM.Indep.letters_independent}), so $\liftonelat'(\tau)\subseteq\liftonelat_B(\tau)$
has rank $p$.

\emph{Indices.} For finitely generated lattices $N\subseteq L$ of the same rank, the index ideal
$\mathrm{idx}(L,N)\subseteq\OO_K$ is nonzero and
$\deg_hL=\deg_hN+\log\Nm\mathrm{idx}(L,N)/[K:\Q]$
(\lref{GlobalTheory.T.lattice_degree_eq_sublattice_degree_add_log_index}); in rank $0$ the index
ideal is $\OO_K$. Let $g_\pm$ be Gram matrices. The development proves that there is an ideal $\mathfrak n$
of $\OO_E$ with $\mathfrak n\bar{\mathfrak n}=(\det g_+)$ and
$\mathrm{idx}(\liftonelat_B(\tau),\liftonelat'(\tau))=\tau(\mathfrak n)\OO_K$
(\lref{GlobalTheory.T.cmDefinite_lieIndex_ideal_eq_norm_ideal}); this is the arithmetic core of the
statement.

\emph{Degrees of the images.} The development computes $\det g_\pm\ne0$ and
\begin{align*}
\deg_h\liftonelat'(\tau)&=\frac p2\Re\liftonelogOm_E(\chi)-\frac{\log|\Nm_{E/\Q}\det g_+|}{2[E:\Q]},\\
\deg_h\liftonelat'(\bar\tau)&=\frac q2\Re\liftonelogOm_E(\chi)-\frac{\log|\Nm_{E/\Q}\det g_-|}{2[E:\Q]}
\end{align*}
(\lref{GlobalTheory.T.degree_images_eigenLines_cmUnitaryShimuraObject}; the constant
$\liftonelogOm_E(\chi)$ enters through the degrees of the eigenlines of the letters). Here $q=0$ and
$\det g_-=1$, so $\deg_h\liftonelat'(\bar\tau)=0$ and hence $\deg_h\liftonelat_B(\bar\tau)=0$.

\emph{Norms.} Let $n=[K:\Q]/[E:\Q]$. Then $\Nm(\tau(\mathfrak n)\OO_K)=\Nm(\mathfrak n)^n$
(the norm of an extended ideal), and $\Nm\bar{\mathfrak n}=\Nm\mathfrak n$,
so $\Nm(\mathfrak n)^2=|\Nm_{E/\Q}\det g_+|$. Hence
\[
\frac{\log\Nm\mathrm{idx}(\liftonelat_B(\tau),\liftonelat'(\tau))}{[K:\Q]}=\frac{n\log\Nm\mathfrak n}{n[E:\Q]}
=\frac{\log|\Nm_{E/\Q}\det g_+|}{2[E:\Q]},
\]
and $\deg_h\liftonelat_B(\tau)=\frac p2\Re\liftonelogOm_E(\chi)$. Subtracting
$\deg_h\liftonelat_B(\bar\tau)=0$ gives the claim.
\end{proof}

Let now $D_w$ be a datum, $X^o$ a moduli space for $D_w$ over a coefficient field $F$ and $U$
uniformisations satisfying the structural laws. For $\tau\in\Psi$, the \emph{flat weight} of $\tau$ is the
parallel weight equal to $-1$ on the entries of $u_\tau$ and $0$ elsewhere.

\begin{proposition}[Degree of flat lines]\label{pf:lift1:framedeg}
Let the host be a host field, $\chi$ a Kronecker character modulo $d$, and $r_-=0$. Let $Z$ be a
component, $\tau\in\Psi$, $s$ a section of the flat weight of $\tau$ on $Z$ with constant archimedean
norms $c_\sigma(s)$, and $\varrho$ a real function on the finite places of $F$ such that every valuation
of $Z$ at a place $w$ takes the value $\varrho_w$ at $s$. Then $\varrho$ has finite support and
\[
\widehat{\deg}\bigl((c_\sigma(s))_\sigma,\varrho\bigr)=\frac{r_+-r_-}2\,\Re\liftonelogOm_E(\chi).
\]
\formallabel{GlobalTheory.DefCoreLie.frameDeg_def}
\end{proposition}

\begin{proof}
There are a number field $K$ containing $F$, a $K$-rational framed point $y$ on $Z$, a CM object of
signature $(r_+,r_-)$ over $K$, with $\rho=\tau'(\mathrm{rt})$ (image in $K$) for every $\tau'\in\Psi$, and
an abelian variety over $K$ mapping to the universal abelian variety at $y$ and to $B_K$
(\lref{GlobalTheory.T.exists_cm_goodReduction_point_on_component}: every component contains a
point whose universal abelian variety is related in this way to the abelian scheme $B$ of a CM object
over $\OO_K$). Choose a generic fibre of $B$
(\lref{GlobalTheory.CM.GenFib.nonempty_of_baseChange}) and let $h$ be the hermitian family induced by
the polarisation. Let $t\colon\OO_E\to\OO_K$ be $\tau$ followed by $F\subseteq K$. The development proves that
$\varrho$ has finite support and that the degree of the flat line is
\[
\widehat{\deg}\bigl((c_\sigma(s))_\sigma,\varrho\bigr)=\deg_h\liftonelat_B(t)-\deg_h\liftonelat_B(\bar t)
\]
(\lref{GlobalTheory.T.flatLineDegree_eq_eigenLattice_degree_sub}). As $t(\mathrm{rt})=\rho$ and
$r_-=0$, \cref{pf:lift1:cmdiff} gives the value $\frac{r_+-r_-}2\Re\liftonelogOm_E(\chi)$.
\end{proof}

\begin{remark}
In the formal statement of \cref{pf:lift1:framedeg} the conclusion of \cref{pf:lift1:cmdiff} (for all
CM objects) is a hypothesis; it is discharged by \cref{pf:lift1:cmdiff} where the proposition is used.
\end{remark}

\subsubsection*{The core datum at the top index}

Let $p$ satisfy $\mathrm{PointConditions}'\wedge\mathrm{RecTied}$ and $i+1=m$, and assume
$\mathsf{GramAt}$, $\liftoneNw(p,i,PP)$ and $\liftoneNh(PP,A,M)$. Let $F_T$ be the coefficient
field of the core datum and $X^o_m$ its moduli space at level $m$. At such points $F_T$ is
algebraically closed in the function fields of the components of $X^o_m$
(\lref{GlobalTheory.CoreDesc.fnAlgClosed_core_alt}); assuming this, the definite case
$\liftoneDef(U,\xi_R)$ holds for the core datum at level $m$ (with $\xi_R=\xi_R(\mathcal D)$), the
set of components is finite and nonempty, and the dictionaries take values in the domains of the
uniformisations (\lref{GlobalTheory.DefCore.defCase_core_alt}).

\begin{proposition}\label{pf:lift1:eigcore}
In this situation, if $F_T$ is algebraically closed in the function fields of the components of
$X^o_m$, there are a finite group $G$, a free level action $\mathrm{lev}$ of $G$ on $X^o_m$ and a
character $\chi\colon G\to F_T^\times$ with $\mu(\liftoneeig(G_\kappa))\le\liftonemuc$.
\formallabel{GlobalTheory.DefEig.exists_finite_group_lev_char_slope}
\end{proposition}

\begin{proof}
$G$ is the finite group $P$ of the level structure of the face of $p$, acting through the level maps of
the orbit data of $p$; this is a level action (\lref{GlobalTheory.DefEig.actKeeps_alt}), and it is
free at the top index (\lref{GlobalTheory.T.freePt_last_index_of_fnAlgClosed}). The lifting map from
cusp forms of weight $G_\kappa$ to the space of lifted forms intertwines the action of $P$ on forms with
its action $\rho$ on lifted forms, and there is a character $\chi$ of $P$ with values in $F_T^\times$
whose image in $K^\times$ is $\chi_{\mathrm{FW}}\psi^{-1}$, where $\chi_{\mathrm{FW}}$ is the character of
the space of lifted forms and $\psi$ the character of the face
(\lref{GlobalTheory.DefEig.exists_chi_algebraMap_eq_FW_chi_mul_psi_inv}). The lattice
$\liftonelat_{\mathrm{type}}$ is the intersection of $\liftonelat(G_\kappa)$ with the span of the forms
whose lifts have the type $(P,\psi^{-1})$, i.e.\ transform under $\rho(u)$ by
$\chi_{\mathrm{FW}}(u)\psi^{-1}(u)$. by the intertwining property and the injectivity of the lifting
map, a form has this property if and only if it is a $\chi$-eigenvector of $P$. Hence
$\liftonelat_{\mathrm{type}}=\liftoneeig(G_\kappa)$, and as the norms defining $\liftonemuc$ are those
of $\mathrm{nrm}(G_\kappa)$, $\mu(\liftoneeig(G_\kappa))=\liftonemuc$.
\end{proof}

\begin{proposition}[Flat sections at the top index]\label{pf:lift1:flatcore}
In the same situation, $\liftoneFl(-\liftoneflat(\liftonelogOm))$ holds for the core datum at
level $m$: flat sections exist on every component for every $\tau\in\Psi$ and all have degree
$\liftoneflat(\liftonelogOm)=\frac r2\liftonelogOm$.
\formallabel{GlobalTheory.DefFlat.flatExistsAt_core_alt}
\end{proposition}

\begin{proof}
Flat sections exist in the definite case
(\lref{GlobalTheory.FlatLine.Definite.nonempty_flatSec}). Let $s$ be a flat section on $Z$ for
$\tau$ and $P$ rational point data for $Z$. By
\lref{GlobalTheory.T.degree_dualDetLie_eq_neg_flatSection_degree}, $\deg s=-d^{\mathrm{Lie}}_P(\tau)$,
and by \cref{pf:lift1:deglie} below $d^{\mathrm{Lie}}_P(\tau)=-\liftoneflat(\liftonelogOm)$.
\end{proof}

\begin{proposition}[Degree of the Lie algebra at CM points]\label{pf:lift1:deglie}
Under the hypotheses of \cref{pf:lift1:eigcore}, for every component $Z$ of $X^o_m$, rational point
data $P$ for $Z$ and $\tau\in\Psi$,
\[
d^{\mathrm{Lie}}_P(\tau)=-\liftoneflat(\liftonelogOm)=-\frac r2\liftonelogOm.
\]
\formallabel{GlobalTheory.DefCoreLie.degLie_core_alt}
\end{proposition}

\begin{proof}
By \lref{GlobalTheory.DefCore.defCase_core_alt} the definite case holds, and the uniformisations of
the core datum satisfy the structural laws (\lref{GlobalTheory.LawsCore.laws_core}). By
\lref{GlobalTheory.T.logPeriod_eq_kroneckerCharacter_logPeriod_of_host} the host is a host field and
there is a Kronecker character $\chi$ modulo $d$ such that $\liftonelogOm(p)=\Re\liftonelogOm_E(\chi)$. By
\cref{pf:lift1:framedeg} (with \cref{pf:lift1:cmdiff}), every flat section $s$ on a component $Z$ for
$\tau\in\Psi$, with its constant archimedean norms and its valuations, spans a hermitian line of degree
$\frac{r_+-r_-}2\Re\liftonelogOm_E(\chi)$. Since $r_-=0$ and
$d^{\mathrm{Lie}}_P(\tau)=-\deg s$ for a flat section $s$
(\lref{GlobalTheory.T.degree_dualDetLie_eq_neg_flatSection_degree}), we get
$d^{\mathrm{Lie}}_P(\tau)=-\frac r2\liftonelogOm=-\liftoneflat(\liftonelogOm)$.
\end{proof}

\begin{remark}
\Cref{pf:lift1:deglie} is the point where the Chowla--Selberg constant $\liftonelogOm$ enters the core
slope: the degree of the Lie line at the CM points of the definite moduli space is $-\frac r2\Re
\liftonelogOm_E(\chi)$, and by the Chowla--Selberg formula \cite{ChowlaSelberg1967}, in the form of
Colmez \cite{Colmez1993}, $\liftonelogOm$ is an explicit affine function of the Faltings height of the
elliptic curves with complex multiplication by the integers of $\Q(\sqrt{-d})$
(\cref{rem:covolume:omega}). The development uses $\liftonelogOm$ only through its defining formula.
\end{remark}

\begin{proposition}[Defect bound]\label{pf:lift1:defect}
Assume (Hdg). There are reals $M$ and $C_0\ge0$ such that for every point $p$ with
$\mathrm{PointConditions}'\wedge\mathrm{RecTied}$, $i+1=m$, $\mathsf{GramAt}(p,i)$,
$\liftoneNw(p,i,PP)$ and $\liftoneNh(PP,A,M)$, if $F_T$ is algebraically closed in the
function fields of the components of $X^o_m$, then
\[
2\,\liftonedefB\bigl(a_\kappa,n_\kappa,-\liftoneflat(\liftonelogOm)\bigr)-C_0U\le\liftonemuc.
\]
\formallabel{GlobalTheory.CovTop.DefCore.defect_bound_of_connected_alt}
\end{proposition}

\begin{proof}
Take $M=C_0=0$. By \cref{pf:lift1:eigcore} there are $G,\mathrm{lev},\chi$ with
$\mu(\liftoneeig(G_\kappa))\le\liftonemuc$; the definite case and the other hypotheses of
\cref{pf:lift1:eiglat} on the core datum hold by \lref{GlobalTheory.DefCore.defCase_core_alt}; and
$\liftoneFl(-\liftoneflat(\liftonelogOm))$ holds by \cref{pf:lift1:flatcore}. The forms of
weight $G_\kappa$ are good by $\mathsf{GramAt}$ and $(a_\kappa,n_\kappa)$ is dominant by
$\liftoneNw$. \Cref{pf:lift1:eigdef} with $c_1=-\liftoneflat(\liftonelogOm)$ gives
$2\,\liftonedefB(a_\kappa,n_\kappa,c_1)\le\mu(\liftoneeig(G_\kappa))\le\liftonemuc$.
\end{proof}

\begin{corollary}\label{pf:lift1:defectlift}
Assume (Hdg). Then the lift condition defined by
\[
i+1=m\ \wedge\ \mathsf{GramAt}(p,i)\ \Longrightarrow\
2\,\liftonedefB\bigl(a_\kappa,n_\kappa,-\liftoneflat(\liftonelogOm)\bigr)-C\,U\le\liftonemuc
\]
holds.
\formallabel{GlobalTheory.CovTop.DefCore.lifts_two_mul_defBound_sub_unitE_le_muCore}
\end{corollary}

\begin{proof}
Combine (\cref{pf:lift1:combine}) the lift conditions defined by $\liftoneNw$ and by
$\liftoneNh(\cdot,\cdot,M)$ for the constant $M$ of \cref{pf:lift1:defect}. At a point where they
hold, $U\ge0$. If $F_T$ is algebraically closed in the function fields of the components of $X^o_m$,
\cref{pf:lift1:defect} applies with its constant $C_0$. Otherwise the statement holds vacuously, as
this closedness always holds at the points considered
(\lref{GlobalTheory.CoreDesc.fnAlgClosed_core_alt}). The proof of \cref{pf:lift1:defect} gives
$C_0=0$, so the constant of the lift condition is $0$.
\end{proof}

\begin{proof}[Proof of \cref{pf:lift1:lowdef-thm}]
By \cref{pf:lift1:combine} we may work at a point $p$ with $i+1=m$ where $\liftoneAnc(C_A)$,
$\mathsf{GramAt}$, $\liftoneNw$, the crude error bound of \cref{pf:lift1:crude} (constant $C_3=35$)
and the inequality of \cref{pf:lift1:defectlift} (constant $C_0$) hold. By $\liftoneNw$, $r_-=0$
and $r_+=r\ge1$, and $n_\kappa$ is even. \Cref{pf:lift1:twodef} with $k_c=k-1$ and
$\ell=\liftonelogOm$, combined with \cref{pf:lift1:crude}, gives
\[
-\sum_{\tau\in\Psi}r_+\mathrm{anc}(\tau)\liftonelogOm-C_3U\le2\,\liftonedefB\bigl(a_\kappa,n_\kappa,-\liftoneflat(\liftonelogOm)\bigr),
\]
and $\liftoneAnc(C_A)$ with $r_-=0$ gives
$-\liftonemainT\le-\sum_\tau r_+\mathrm{anc}(\tau)\liftonelogOm+C_AU$. With \cref{pf:lift1:defectlift},
$-\liftonemainT-(C_A+C_3+C_0)U\le\liftonemuc$, which is $\mathsf{LowAt}(C_A+C_3+C_0)$.
\end{proof}

% ---------------------------------------------------------------------------
\subsection{The height bound}\label{pf:lift1:height}

Throughout this subsection $p$ is a point, $i+1=j<m$, $Y$ is a good deep cover of level $j$ with
coefficient field $F=F_Y$ and level $l_Y$, and $u_1$ is an admissible weight; all forms, components and
invariants are those of $Y$. We write $|u|_1=\sum_{\tau,a}|u_{\tau,a}|$.

\subsubsection*{The predicates}
\begin{definition}[Height predicate]\label{pf:lift1:heightat}
$\liftoneH(PP,C,u_1)$ means: there are reals $c_s,c_c,X\ge0$ and a function
$\mathrm{sl}\colon\mathrm{Comp}\to\R$ such that
\begin{enumerate}[label=(H\arabic*)]
\item\label{pf:lift1:H1} $0\le\liftonehgt(u_1,Z)+c_s$ for every component $Z$;
\item\label{pf:lift1:H2} $(1-\rho(u_1))\max(0,\mathrm{sl}(Z)+c_s)\le\liftonehgt(u_1,Z)+c_s+c_c$ for every $Z$;
\item\label{pf:lift1:H3} $-X\le\frac1{\#\mathrm{Comp}}\sum_Z\mathrm{sl}(Z)$;
\item\label{pf:lift1:H4} $TX\le CU$, $Tc_c\le CU$ and
  $c_s\bigl(6\liftonedfl+10+T(\zeta(u_1)+2\rho(u_1))\bigr)\le CU$.
\end{enumerate}
The \emph{height condition} is the cover condition defined by requiring $\liftoneH(PP,C,u_1)$ for
every admissible weight $u_1$.
\end{definition}

Recall the mean-slope bound $\mathrm{MeanAll}(s,C)$ and the quantities $\mathrm{WtGap}$,
$\liftonebnd(u,g)$, $\liftoneerr_s(u,g)$ and $\ell(\mathcal D)=\log\liftonerd(E)+|\frac{L'}{L}(1,\chi_d)|$ of
\cref{def:covolume:deepstr}. The volume slope there is defined as follows. For a family $\mathcal V$ of
valuations at the finite places of $F$, a component $Z$, a weight $G$ and $t\in\R$, let
$\liftonevol_{\mathcal V}(Z,G,t)\in[0,\infty]$ be the normalised growth (the development's
counting volume, with multiplicity $8$) of the numbers of forms of weight $\tau G$ on $Z$, $\tau\to\infty$,
which are integral for $\mathcal V$ and have pointwise norm at most $e^{\tau t}$ at all frames of $Z$
at the good embeddings. For a weight $u$ with $\mathrm{top}(Z,u)>0$ and $s\in\R$,
\begin{equation}\label{pf:lift1:eq:volslope}
s\le\mathrm{volSlope}_{\mathcal V}(Z,u)\iff\mathrm{top}(Z,u)\,(s+t)\le\liftonevol_{\mathcal V}(Z,\mathrm{wt}(u),t)
\text{ for all }t\in\R
\end{equation}
(where a negative left side is read as $0$). In $\mathrm{MeanAll}$ the family is
$\mathcal V_B$, $B=\liftonebnd(u,g)$: the \emph{height valuations} at the places over
the primes of $P_B=\{q\text{ prime}:q\le B\text{ or }q\mid2|d_E|M'\}$ ($M'$ the level
modulus of $Y$) and the standard valuations elsewhere. Similarly
$\liftonevol^{\mathrm B}(Z,G,t)$ is the counting volume of the forms $f$ of weight $\tau G$ on $Z$ with
$v(f)\ge\tau\,\mathrm{depth}_v(Z,G)$ for every standard valuation $v$ (where $\mathrm{depth}_v(Z,G)$ is the
minimum of $v$ on $\liftonelat_Z(G)$) and $|f(z)|^2\le e^{2\tau t}\mathrm{Berg}_Z(G)(z)^\tau$ at all frames
$z$, $\mathrm{Berg}_Z(G)$ being the Bergman kernel of $\liftonecusp_Z(G)$.

\begin{definition}[Inputs of the height bound]\label{pf:lift1:inputs}
Let $\Lambda\ge1$, $c\ge0$ and $u,u'$ weights.
\begin{enumerate}[label=(\alph*)]
\item A family $w_1,\dots,w_n\in\liftonecusp_Z(\mathrm{wt}(u))$ is a \emph{short generating family} of
  size $\Lambda$ if it generates $\liftonelat_Z(\mathrm{wt}(u))$ as a $\Z$-module and
  $\mathrm{nrm}_{Z,\sigma}(w_l)\le\Lambda$ for all $l,\sigma$.
\item $\mathrm{Mult}(Z,u,u',c)$ means: no finite place is a base place of the forms of weight
  $\mathrm{wt}(u)$ on $Z$ (if some form on $Z$ has finite valuation $v$, some element of
  $\liftonelat_Z(\mathrm{wt}(u))$ does), and there are a rule $\theta$, reals $e_w\ge0$ at finitely many
  finite places and reals $t_\sigma$ (invariant under complex conjugation) such that: $\theta$ is a section
  of exponent $8(u-u')$, nonzero at the generic point of $Z$, with $\theta^kf$ a form of weight
  $8k\,u$ on $Z$ for every form $f$ of weight $8k\,u'$ on $Z$; $8\,\mathrm{depth}_v\le v(\theta)+8e_w$ for
  every finite place $w$ and every standard valuation $v$ at $w$ which is finite on
  $\liftonelat_Z(\mathrm{wt}(u))$; $|\theta(z)|^2\le e^{16t_\sigma}\mathrm{Berg}_Z(\mathrm{wt}(u))(z)^8$ at
  every frame $z$ of $Z$ at every good embedding $\sigma$; and
  $\mathrm{cost}(e,t)=\frac1{[F:\Q]}(\sum_we_w\log\Nm w+\sum_\sigma t_\sigma)\le c$.
\item The \emph{height comparison property} $\liftoneHC(Y)$ means: for every prime $q$, every finite place $w\mid q$ of $F$ and every
  standard valuation $v$ at $w$ there is a height valuation $v'$ at $w$ with
  $v'(f)\le v(f)+v_w(q)\,\eta_q\,|u|_1$ for every weight $u$ and every section $f$ of exponent
  $u$, where $\eta_q=\frac1{q-1}+40\frac{v_q(d_E)}{[E:\Q]}$. Put
  $\liftonehtc(B)=\sum_{q\in P_B}\eta_q\log q$.
\end{enumerate}
The \emph{pointwise inputs} at $(p,i,Y,u_1)$ are:
\begin{itemize}
\item $\liftoneS(PP,C,u_1)$: there is $\Lambda\ge1$ such that every component $Z$ has a short
  generating family of size $\Lambda$ for $u_1$, and
  $\log\Lambda\,(6\liftonedfl+10+T(\zeta(u_1)+2\rho(u_1)))\le CU$;
\item $\liftoneG(u_1)$: for every $Z$, the forms of weight $\mathrm{wt}(u_1)$ on $Z$ are good,
  their ratios generate the function field of $Z$ at its generic point (\emph{generic generation}), and
  $(1-\rho(u_1))\liftonedg(u_1,Z)\le\delta'(Z,u_1)$;
\item $\liftoneM(PP,C,u_1)$: there is $c_c\ge0$ with $Tc_c\le CU$ and
  $\mathrm{Mult}(Z,u_1,u_1',c_c)$ for every $Z$;
\item $\liftoneB(PP,C,u_1)$: the level modulus $M'$ has at most $15$ prime factors, and there is
  $g$ with $\mathrm{WtGap}(u_1',g)$, $0\le\liftoneerr:=\liftoneerr_{\ell(\mathcal D)}(u_1',g)$,
  $T\liftoneerr\le CU$, $0\le\liftonehtc(B)$ and $T\liftonehtc(B)|u_1'|_1\le CU$, where
  $B=\liftonebnd(u_1',g)$;
\item $\liftoneP(u_1)$: there is $\delta>0$ with
  $\delta\le\liftonedg(u_1,Z)\le\delta/(1-\rho(u_1))$ for all $Z$.
\end{itemize}
The \emph{short}, \emph{geometric}, \emph{multiplicity} and \emph{budget conditions} are the cover
conditions defined by requiring these predicates for every admissible $u_1$;
the cover degree condition (\cref{tab:covolume:conditions}) is the one defined by $\liftoneP$.
\end{definition}

\subsubsection*{The pointwise argument}

\begin{lemma}\label{pf:lift1:htnonneg}
Assume (Hdg). If the forms of weight $\mathrm{wt}(u_1)$ on $Z$ are good, $\liftonedg(u_1,Z)>0$ and $Z$
has a short generating family of size $\Lambda$ for $u_1$, then $0\le\liftonehgt(u_1,Z)+\log\Lambda$.
\formallabel{GlobalTheory.CovHt.HtTopG.height_add_log_nonneg}
\end{lemma}

\begin{proof}
$F$ is totally complex (complex conjugation fixes no embedding of the coefficient field) and the norms
$\mathrm{nrm}_{Z,\sigma}$ are invariant under $\sigma\mapsto\bar\sigma$. The development proves, for a
hermitian lattice $\liftonelat$ with conjugation invariant norms having a Gram frame, spanning its
space and generated over $\Z$ by vectors of norm $\le\Lambda$, and a linear map to a field whose image
has positive degree, that the height of the image is at least $-\log\Lambda$
(\lref{GlobalTheory.ImageHeightNonneg.imageHeight_add_log_nonneg}). Apply it to
$\liftonelat_Z(\mathrm{wt}(u_1))$ and the generic point map of $Z$. Hypothesis (Hdg) is not used.
\end{proof}

\begin{lemma}[Slope versus height]\label{pf:lift1:slopeht}
Assume (V) and (Hdg). Let $Y$ be good with $\liftoneHC(Y)$, and let $Z$ be a component with a short
generating family of size $\Lambda$ for $u_1$ such that the forms of weight $\mathrm{wt}(u_1)$ on $Z$ are
good, generic generation holds, $(1-\rho(u_1))\liftonedg(u_1,Z)\le\delta'(Z,u_1)$, $\rho(u_1)<1$ and
$\liftonedg(u_1,Z)>0$. Let $\mathrm{Mult}(Z,u_1,u_1',c_c)$ with $c_c\ge0$, $B\in\R$ with
$h_c:=\liftonehtc(B)|u_1'|_1\ge0$, and $c_s\ge\log\Lambda$ with $\liftonehgt(u_1,Z)+c_s\ge0$. If
$s\le\mathrm{volSlope}_{\mathcal V_B}(Z,u_1')$, then
\[
(1-\rho(u_1))\max(0,s+c_s)\le\liftonehgt(u_1,Z)+c_s+c_c+h_c.
\]
\formallabel{GlobalTheory.CovHt.HtTopG.slope_le_height}
\end{lemma}

\begin{proof}
Write $\rho=\rho(u_1)$, $\delta=\liftonedg(u_1,Z)$ and $h=\liftonehgt(u_1,Z)$; it suffices to prove the
inequality with $\varepsilon>0$ added on the right. Since $\mathrm{top}(Z,u_1')=\delta'(Z,u_1)\ge
(1-\rho)\delta>0$, \eqref{pf:lift1:eq:volslope} with $t=c_s$ gives
$\delta'(Z,u_1)(s+c_s)\le\liftonevol_{\mathcal V_B}(Z,\mathrm{wt}(u_1'),c_s)$. Two estimates of the
development bound the right side:
\begin{enumerate}[label=(\roman*)]
\item (multiplication by $\theta$) if $\mathrm{Mult}(Z,u,u',c)$ and $\liftoneHC(Y)$, then for all
  $t$ and $\varepsilon>0$,
  $\liftonevol_{\mathcal V_B}(Z,\mathrm{wt}(u'),t)\le\liftonevol^{\mathrm B}(Z,\mathrm{wt}(u),t+c+\liftonehtc(B)|u'|_1+\varepsilon)$
  (\lref{GlobalTheory.CovHtCmp.compVol_le_volL}): a small form $f$ of weight $8k\,u'$ gives the form
  $\theta^kf$ of weight $8k\,u$; the archimedean and finite bounds for $\theta$ cost $c$, and the
  comparison of height and standard valuations at the primes of $P_B$ costs
  $\liftonehtc(B)|u'|_1$;
\item (volume versus height) if the forms of weight $\mathrm{wt}(u)$ on $Z$ are good, generic generation
  holds, no finite place is a base place, $\liftonedg(u,Z)>0$ and $Z$ has a short generating family of
  size $\Lambda\le e^t$, then
  $\liftonevol^{\mathrm B}(Z,\mathrm{wt}(u),t)\le\liftonedg(u,Z)\,(\liftonehgt(u,Z)+t)$
  (\lref{SelfIntersectionVolumeEstimates.VolUpG.volume_line_bundle_le}; its proof constructs a nef
  hermitian model of the image of $Z$,
  \lref{SelfIntersectionVolumeEstimates.ImageNefG.exists_nef_model_of_image}, and bounds the
  counting volume by its arithmetic self-intersection,
  \lref{GlobalTheory.T.countingVolume_le_selfIntersection_nefHermitianBundle}).
\end{enumerate}
Applying (i) with $(u,u',c,t)=(u_1,u_1',c_c,c_s)$ and (ii) with $u=u_1$ and
$t=c_s+c_c+h_c+\varepsilon\ge\log\Lambda$ (the absence of base places is part of $\mathrm{Mult}$), we get
$\delta'(Z,u_1)(s+c_s)\le\delta(h+c_s+c_c+h_c+\varepsilon)$, the right side being $\ge0$. If $s+c_s\le0$
the claim holds because its right side is $\ge0$. Otherwise
$(1-\rho)\delta(s+c_s)\le\delta'(Z,u_1)(s+c_s)\le\delta(h+c_s+c_c+h_c+\varepsilon)$, and we divide by
$\delta>0$.
\end{proof}

\begin{proposition}\label{pf:lift1:heightpt}
Assume (V) and (Hdg). Let $Y$ be good (for some parameters), with $\liftoneHC(Y)$ and nonempty
$\mathrm{Comp}$; let $C_0,C_s,C_m,C_b\ge0$, $U\ge0$, $\mathrm{MeanAll}(\ell(\mathcal D),C_0)$ on $Y$ (with
$x=\xi_R(\mathcal D)$), $\rho(u_1)<1$, and assume $\liftoneP(u_1)$, $\liftoneS(PP,C_s,u_1)$,
$\liftoneG(u_1)$, $\liftoneM(PP,C_m,u_1)$ and $\liftoneB(PP,C_b,u_1)$. Then
$\liftoneH(PP,C_s+C_m+C_b+C_0C_b,u_1)$ holds.
\formallabel{GlobalTheory.CovHt.HtTopG.heightAt_eq_of_cover}
\end{proposition}

\begin{proof}
Let $\Lambda$, $c_c$ and $g$ be given by $\liftoneS$, $\liftoneM$ and $\liftoneB$, and
$B=\liftonebnd(u_1',g)$. $\mathrm{MeanAll}$ applied to $(u_1',g)$ gives $\mathrm{sl}$ with
$\mathrm{sl}(Z)\le\mathrm{volSlope}_{\mathcal V_B}(Z,u_1')$ and
$\frac1{\#\mathrm{Comp}}\sum_Z\mathrm{sl}(Z)\ge-C_0\liftoneerr$. By $\liftoneP$,
$\liftonedg(u_1,Z)>0$ for all $Z$; as $\mathrm{Comp}\ne\emptyset$, $\Lambda\ge1$. Take
\[
c_s=\log\Lambda\ge0,\qquad c'_c=c_c+\liftonehtc(B)|u_1'|_1\ge0,\qquad X=C_0\liftoneerr\ge0 .
\]
\ref{pf:lift1:H1} is \cref{pf:lift1:htnonneg}; \ref{pf:lift1:H2} (with $c'_c$) is \cref{pf:lift1:slopeht},
whose hypotheses are supplied by $\liftoneG$, $\liftoneM$ and \ref{pf:lift1:H1};
\ref{pf:lift1:H3} is the mean bound. For \ref{pf:lift1:H4}: $TX=C_0T\liftoneerr\le C_0C_bU$;
$Tc'_c\le C_mU+C_bU$; and $c_s(6\liftonedfl+\dots)\le C_sU$. Each is at most
$(C_s+C_m+C_b+C_0C_b)U$ since all constants and $U$ are $\ge0$.
\end{proof}

\subsubsection*{The inputs}

\begin{lemma}[Short generators]\label{pf:lift1:short}
Assume (Hdg), the deep goodness condition, the size condition and the flat degree
condition. Then the short condition holds.
\formallabel{GlobalTheory.CovHt.ShortPf.cover_height_short_section_alt}
\end{lemma}

\begin{proof}
Let $C\ge0$ be the constant of \lref{GlobalTheory.T.short_generator_norm_comparison}, which states:
for a core datum with the structural laws, a level $j$ with $r_+,r_-\ge1$ and $2<[E:\Q]/2$, a good
deep cover $Y$ (parameters $e,o$), and $L_b\ge1$ with $\log M\le30L_b$,
$\log[E:\Q]+\log|d_E|/[E:\Q]\le3L_b$, $\log(fr)\le L_b$, $o\le391(frL_b)^2$, if $\liftoneFl(c_1)$ holds on $Y$ with
$|c_1|\le9rL_b$, then for every weight $u$ which is non-increasing in each $\tau$, has gap
$u_{\tau,r_+-1}-u_{\tau,r_+}\ge8fr+10$ and $\log|u|_1\le10L_b$, every component has a short generating
family for $u$ of size $\Lambda$ with $\log\Lambda\le C((frL_b)^3|u|_1+(frL_b)^5)$. We apply it with
$L_b=L$ and $u=u_1$ at points where $\liftoneNh(PP,A,1)$, $\liftoneNw$ and the flat degree
condition hold: the level of a good deep cover satisfies $\log M\le30L$ (its modulus is
$g_0g_1\cdot15^{e}$ with $e=\lfloor3L/\log3\rfloor$), the host relations give the size bounds on $E$
and $fr$, the flat degree satisfies $|\liftoneflat(\liftonelogOm)|\le9rL$, and the development proves
(\lref{GlobalTheory.CovHt.ShortPf.weight_rows}) that $u_1$ is non-increasing with the required gap,
$\log|u_1|_1\le10L$, $r_+,r_-\ge1$, and
\[
\bigl((frL)^3|u_1|_1+(frL)^5\bigr)\bigl(6\liftonedfl+10+T(\zeta(u_1)+2\rho(u_1))\bigr)\le693\,U .
\]
Taking $\Lambda=\exp(C((frL)^3|u_1|_1+(frL)^5))$ gives $\liftoneS(PP,693C,u_1)$. The deep goodness
and size conditions are not used.
\end{proof}

\begin{lemma}[Geometric bound]\label{pf:lift1:geom}
Assume (X), (GM), (Q), (Hdg), (Sieg), (Tor) and (Lie). Let $p$ satisfy $\mathsf{SizeAt}(PP,C)$,
$\liftoneNh(PP,A,M)$ and $\liftoneNw(p,i,PP)$ with $j<m$, let $Y$ be a good deep cover and $u_1$
admissible. Then $\liftoneG(u_1)$ holds.
\formallabel{GlobalTheory.CovHt.GeomBound.geometric_height_bound_at_place}
\end{lemma}

\begin{proof}
By $\liftoneNw$, $r_+=k+j-1\ge1$ and $r_-=m-j\ge1$; by $\liftoneNh$, $[E:\Q]/2=f>2$. The
successive differences of $u_1$ are at least $(\liftonedfl+3)^2+2r$
(\lref{GlobalTheory.CovHt.GeomBound.quo_gaps}, which states this bound for admissible weights); in
particular $u_1$ is strictly decreasing in each $\tau$ and the gap at $r_+-1$ is $\ge1$. Then the forms
of weight $\mathrm{wt}(u_1)$ are good on every component
(\lref{GlobalTheory.T.covering_domain_lattice_regularity}, which states that the forms of weights
with non-increasing entries and a positive gap at $r_+-1$ are good on good deep covers), generic generation
holds on every component (\lref{GlobalTheory.T.birOn_wtOf_of_consecutive_gap_bound}, which states
this for weights with large consecutive gaps), and $\liftonedg(u_1,Z)\le\mathrm{top}(Z,u_1)$ for strictly
decreasing $u_1$ by the proportionality principle for the toroidal compactification
(\lref{GlobalTheory.T.proportionality_principle_compactification} and
\lref{GlobalTheory.DegTop.compDeg_le}, which deduce this inequality from Mumford's
proportionality). As $\liftonevdm(u_1)\ne0$, $\delta'(Z,u_1)=(1-\rho(u_1))\mathrm{top}(Z,u_1)$
(\lref{GlobalTheory.CompHt.delPrime_eq}, the identity of \cref{pf:lift1:setting}), and
$1-\rho(u_1)>0$ by the size condition; multiplying gives
$(1-\rho(u_1))\liftonedg(u_1,Z)\le\delta'(Z,u_1)$.
\end{proof}

\begin{corollary}\label{pf:lift1:geomcond}
Assume (X), (GM), (Q), (Hdg), (Sieg), (Tor), (Lie), the deep goodness condition and the size
condition. Then the geometric condition holds.
\formallabel{GlobalTheory.CovHt.GeomBound.geom_of_prints_toroidal_deepGood_size}
\end{corollary}

\begin{proof}
Combine (\cref{pf:lift1:combine}) the size condition with the numerical conditions of
\cref{pf:lift1:num} and apply \cref{pf:lift1:geom}; the constant is $0$. The deep goodness condition is
not used.
\end{proof}

\begin{lemma}[Multiplicity]\label{pf:lift1:mult}
Assume the reference weight condition: the cover condition defined by: for admissible $u_1$ there are a
finite set $S$ of primes, reals $e_q\ge0$ ($e_q=0$ for $q\notin S$) and $s$ with
$\sum_{q\in S}e_q\log q\le s$, $Ts\le CU$, and $\epsilon_q(\mathrm{wt}(u_1))\le e_q$ for every
prime $q$. Assume the carrying section condition: the cover condition defined by: for admissible $u_1$
there is $c_c\ge0$ with $Tc_c\le CU$ such that every component $Z$ carries $\theta,e,t$ as in
$\mathrm{Mult}(Z,u_1,u_1',c_c)$, except that the finite bound is replaced by the lower bound
$v(\theta(y))\ge\frac8{2500}\phi/e_0-8e_w$ for the normalised valuation at every integral algebraic
framed point $y$ over $w$ with an invariant $\phi$, at which some form on $Z$ does not vanish. Then the
multiplicity condition holds.
\formallabel{GlobalTheory.CovMult.multC_of_refW_alt}
\end{lemma}

\begin{proof}
At a point where both inputs hold with constants $C_1,C_2$, let $S,e,s$ and $c_c$ be as given; then
$s\ge0$. For each $Z$, the development derives from the reference exponent bounds and the lower bound
for the valuations of $\theta$ at integral points both the absence of base places and the finite bound
required in $\mathrm{Mult}$, for exponents $e'_w$ obtained by spreading the $e_q$ over the places above
$q$, with $\mathrm{cost}(e',t)\le\mathrm{cost}(e,t)+s$ (\lref{GlobalTheory.CovMult.mult_of}, which uses
the unitarity of the uniformisations of good deep covers). Hence $\mathrm{Mult}(Z,u_1,u_1',c_c+s)$, and
$T(c_c+s)\le(C_1+C_2)U$.
\end{proof}

Here $\epsilon_q(G)=\epsilon_q(Y,G)$, the \emph{reference exponent} of $Y$ at $q$, is the supremum over
places $w\mid q$, integral algebraic framed points $y$ over $w$ (with values in complete discrete
valuation rings $T$) and admissible invariants $\phi$ of $y$, of the depth of the forms of weight $G$ at
$y$ minus $\phi/(2500\,v_T(q))$. The same definition applies to the core moduli space $X^o_j$ of $p$ at
level $j$ with its uniformisations; we write $\epsilon_q(X^o_j,G)$. The two inputs of
\cref{pf:lift1:mult} are proved in \cref{pf:lift1:refw,pf:lift1:carry} below.

\subsubsection*{The reference weight condition}
A \emph{prime budget} is a triple $(S,e,s)$ of a finite set $S$ of primes, a function
$e\colon\N\to\R_{\ge0}$ vanishing outside $S$ and a real $s$ with $\sum_{q\in S}e_q\log q\le s$. The
reference weight condition of \cref{pf:lift1:mult} is the cover condition defined by: for every
admissible $u_1$ there is a prime budget $(S,e,s)$ with $Ts\le CU$ and
$\epsilon_q(Y,\mathrm{wt}(u_1))\le e_q$ for every prime $q$. It is reached in four steps: an oscillation
bound on the core moduli space, a bound for the reference exponents of the core moduli space for an
auxiliary weight (the theta weight), the passage to deep covers, and the passage from the theta weight
to $\mathrm{wt}(u_1)$.

\emph{Factors and theta weights.} For the datum $D_j$ let $\Phi$ be the set of \emph{factors}
$f=(\tau,l)$, $\tau\colon E\to\C$, $0\le l<\mathrm{mult}(\tau)$, with rank $\liftonerk f=l+1$, and put
$n_c=8\bigl(\lfloor\frac18\sum_{f\in\Phi}2\liftonerk f\rfloor+\#\Phi+1\bigr)$, a multiple of $8$ with
$\sum_f2\liftonerk f\le n_c$. A factor $(\tau,l)$ is \emph{big} if $\tau\in\Psi$ and $l+2\ge k$. For a
weight $\lambda$ and $m\in\Z$ let $\mathrm{wt}_m(\lambda)$ be the weight of forms with $n=2m$ and
exponents $\lambda_{\tau,l}-m$ at $(\tau,l)$ for $\tau\in\Psi$, $l<r_+$, and
$-\lambda_{\bar\tau,r-1-l}-m$ at $(\tau,l)$ for $\bar\tau\in\Psi$, $l<r_-$ (other exponents $0$); thus
$\mathrm{wt}(\lambda)=\mathrm{wt}_0(\lambda)$, and $\mathrm{wt}_m$ inverts the passage
$(a,n)\mapsto\lambda$ of \cref{pf:lift1:setting}. For a weight $u$ let
$\sigma_u(\tau)=\frac{n_c}2+1-u_{\tau,r_+-1}$, let $\mathrm{par}(\sigma_u)$ be the parallel weight equal
to $\sigma_u(\tau)$ on all entries of $u_\tau$, and let $a(u)$ be the exponents of
$\mathrm{wt}_{n_c/2}(u+\mathrm{par}(\sigma_u))$, negative values replaced by $0$. The \emph{theta
weight} of $u$ is $G(u)=\mathrm{wt}(a(u),n_c)$, and its \emph{gap exponents} are
$\gamma_u(\tau,l)=a(u)(\tau,l)-a(u)(\tau,l+1)$ (truncated subtraction, with $a(u)(\tau,\mathrm{mult}(\tau))=0$).

\emph{Depths and oscillation.} Let $q$ be a prime and $y$ an integral framed point of $X^o_j$ over a
complete discrete valuation ring $T$ with an embedding of $F$ into its fraction field. For $f\in\Phi$
and $k\in\N$ let $\nu_q(y,f,k)\in[-\infty,\infty]$ be the depth of the $k$-th power of the determinant of
the frame of $f$ in an integral chart of the formal group at $y$, minimised over the integral charts and
over the extensions of the valuation of $T$ normalised by $v(q)=1$; for $\gamma\colon\Phi\to\N$ put
$\nu_q(y,\gamma)=\sum_f\nu_q(y,f,\gamma(f))$. For $b\in\R$, the \emph{oscillation bound}
$\mathrm{Osc}_q(X^o_j,\gamma)\le b$ means: for every place $\lambda$ of $F$ above $q$ there is $A\in\N$ such
that for every such $T$ inducing $\lambda$ and every integral framed point $y$ over $T$,
$A\le\nu_q(y,\gamma)$, and $\nu_q(y,\gamma)\le A+b$ if $y$ is algebraic.

\begin{definition}[Intermediate conditions]\label{pf:lift1:refconds}
In the following, $u_1$ runs over the admissible weights and $(S,e,s)$ over the prime budgets with
$Ts\le CU$; $g_0,g_1$ are the auxiliary primes of the host.
\begin{enumerate}[label=(\alph*)]
\item The \emph{sum point condition} is the lift condition defined by: if $j<m$, then for every $u_1$
  there is $(S,e,s)$ with $\mathrm{Osc}_q(X^o_j,\gamma_{u_1})\le e_q$ for every prime $q\ne g_0,g_1$.
\item The \emph{reference core condition} is the lift condition defined by: if $j<m$, then for every
  $u_1$ there is $(S,e,s)$ with $\epsilon_q(X^o_j,G(u_1))\le e_q$ for every prime $q$.
\item The \emph{reference theta condition} is the cover condition defined by: for every $u_1$ there is
  $(S,e,s)$ with $\epsilon_q(Y,G(u_1))\le e_q$ for every prime $q$.
\end{enumerate}
\end{definition}

The sum point condition is deduced from two pointwise predicates. Let $S_0$ be the set of \emph{small
primes} of the host (the prime factors of an integer $Q_S\ge1$ of the host), $v_d(q)\ge0$ ($q\in S_0$)
its discriminant depths, and $q_H$ the second prime of the host besides $d$. For $C\ge0$:
\begin{itemize}
\item $\mathrm{Wt}(C,u_1)$ means: $r_++r_-+1=k+m$, $r_-\le m\le k$, $m\ge1$, $k\ge2$, $k\le r_++1$,
  $[E:\Q]=2f$, and there are $B\in\N$ and $S_g\in\R$ such that $u_1$ is non-increasing in each $\tau$,
  $u_{1,\tau,k-2}=0$, $u_{1,\tau,r_+-1}-u_{1,\tau,r_+}\ge n_c+1$, $|u_{1,\tau,l}|\le B$ and
  $\sum_{l<k-2}u_{1,\tau,l}\le S_g$ for all $\tau,l$, and $T(B+1)\le Cb$, $TS_g\le Cb$ and $T\ge1$;
\item $\mathrm{Num}(C)$ means: $L\ge1$, $b\ge1$, $\log b\le CL$, $\log\log x\ge1$,
  $fmbL\log\log x\le CU$, $\log\log k\le C\log\log x$, $\log r\le CL$, $\log_2Q_S\le CL$, $Tr^3\le Cb$, and
  $r\log q/(q-1)\le CmL$ for $q\in\{q_H,d\}$.
\end{itemize}
For reals $M_b,S_m$ and $M\in\N$ let $e=e_{M_b,S_m,M}\colon\N\to\R$ be the function
\begin{equation}\label{pf:lift1:eq:eof}
e(q)=\begin{cases}
0,&q>M,\\
fm\bigl(1+\tfrac{v_d(q)}2\bigr)M_b+S_m,&q\le M,\ q\in S_0,\\
\frac{frM_b+S_m}{q-1},&q\le M,\ q\in\{q_H,d\}\setminus S_0,\\
e_{\mathrm{sm}}(q)+1,&q\le M,\ q\notin S_0\cup\{q_H,d\},\ mq<k,\\
\frac{6fmM_b+S_m}{q-1}+1,&\text{otherwise},
\end{cases}
\end{equation}
where $e_{\mathrm{sm}}(q)=\bigl(fmM_b(4+\frac12\log_qM)+S_m\bigr)/(q-1)+f(m+1)r^2(\log_qM+1)$.

\begin{lemma}[Pointwise sum bound]\label{pf:lift1:sumat}
Let $j<m$, $C_1,C_2\ge0$ and $K=2C_1+C_2+1$, and let $u_1$ be a weight with $\mathrm{Wt}(C_1,u_1)$ and
$\mathrm{Num}(C_2)$. Then there is a prime budget $(S,e,s)$ with
$Ts\le40\,(4C_1K^2+C_2K+2C_1)\,C_2\,U$ and $\mathrm{Osc}_q(X^o_j,\gamma_{u_1})\le e_q$ for every prime
$q\ne g_0,g_1$.
\formallabel{GlobalTheory.RefSum.sumAt_of_sizes}
\end{lemma}

\begin{proof}
Let $B,S_g$ be as in $\mathrm{Wt}(C_1,u_1)$, $\gamma=\gamma_{u_1}$, and
\[
M_b=B+1,\qquad S_m=f(S_g+2r_-B),\qquad M=2B+1 .
\]
\emph{Exponents.} The development proves from the conditions on $u_1$ in $\mathrm{Wt}$ that $\gamma$
satisfies the \emph{exponent bound} for $(M_b,S_m,M)$: for every $\tau$ the sum of $\gamma(f)$ over the
big factors $f=(\tau,l)$ is at most $M_b$, $\sum_{f\text{ not big}}\liftonerk f\,\gamma(f)\le S_m$, and
$\gamma(f)\le M$ for all $f$ (\lref{GlobalTheory.T.thetaExponents_bigSmallFactorBounds}).

\emph{The host.} One has $2\in S_0$, $\sum_{q\in S_0}(1+\frac{v_d(q)}2)\log q\le\frac32\log_2Q_S$, and
every prime outside $S_0\cup\{q_H,d\}$ is unramified in $E$
(\lref{UpperTheory.KR.small_vd_sum}); and at every $q\in S_0$ the places of $E$ above $q$ have the
splitting type and discriminant depth $v_d(q)$ recorded by the host
(\lref{GlobalTheory.CharDepth.twoPlaces_host}).

\emph{Oscillation.} For a datum with $r_++r_-+1=k+m$, $r_-\le m\le k$, $m\ge1$, $k\ge2$, $[E:\Q]=2f$,
exponents $\gamma$ with the exponent bound for $(M_b,S_m,M)$ and the host data just described, the
development proves $\mathrm{Osc}_q(X^o_j,\gamma)\le e_{M_b,S_m,M}(q)$ for every prime $q$
(\lref{GlobalTheory.T.depthOscillation_le_bound_every_prime}). This is the main estimate behind
the reference weight condition; it bounds, place by place, how much the depths of the theta exponents
can vary over the integral points.

\emph{The prime sum.} Put $\ell_2=KL$ and $\ell_1=K\log\log x$; then $\ell_1,\ell_2\ge1$. The development
proves (\lref{GlobalTheory.T.weightedPrimeLogSum_oscillationBound_le}): if $M_b,S_m\ge0$,
$r+1=k+m$, $1\le m\le k$, $\log M\le\ell_2$, $\log r\le\ell_2$,
$\sum_{q\in S_0}(1+\frac{v_d(q)}2)\log q\le\frac32\ell_2$, $\log\log k\le\ell_1$ and
$r\log q/(q-1)\le m\ell_2$ for $q\in\{q_H,d\}$, then the set of primes $q\le M$, the function equal to $e$
at primes and to $0$ elsewhere, and
\[
s=40\bigl((fmM_b+S_m)\ell_2\ell_1+fr^3\ell_2+M\bigr)
\]
form a prime budget. The hypotheses hold: $\log(2B+1)\le\log(2C_1b)\le2C_1+C_2L\le\ell_2$, since
$2B+1\le2T(B+1)\le2C_1b$; the others follow from $\mathrm{Num}(C_2)$, the host bound above and $C_2\le K$.

\emph{The budget.} Using $T(B+1)\le C_1b$, $TS_g\le C_1b$, $r_-\le m$, $Tr^3\le C_2b$, $f,m,L,\log\log x\ge1$
and $b\le fmbL\log\log x\le C_2U$,
\begin{align*}
TfmM_b\,\ell_2\ell_1&\le C_1K^2\,fmbL\log\log x\le C_1C_2K^2U,\\
TS_m\,\ell_2\ell_1&\le3C_1K^2\,fmbL\log\log x\le3C_1C_2K^2U,\\
Tfr^3\ell_2&\le C_2K\,fbL\le C_2^2KU,\qquad TM\le2T(B+1)\le2C_1b\le2C_1C_2U,
\end{align*}
so $Ts\le40(4C_1K^2+C_2K+2C_1)C_2U$. Together with the oscillation estimate (which holds at every prime,
in particular at the primes $q\ne g_0,g_1$) this proves the lemma.
\end{proof}

\begin{theorem}[Sum points]\label{pf:lift1:sumpt}
The size condition implies the sum point condition.
\formallabel{GlobalTheory.T.sum_pt_of_size}
\formallabel{GlobalTheory.RefSum.sumPt_of_size_alt}
\end{theorem}

\begin{proof}
The size condition implies the \emph{weight point condition}, the lift condition defined by: if $j<m$,
then $\mathrm{Wt}(C,u_1)$ for every admissible $u_1$
(\lref{GlobalTheory.T.weight_pt_of_size}, which derives the bounds $B,S_g$ from the construction of
the admissible weights and the size relations), and the \emph{numerical point condition}, the lift
condition defined by: if $j<m$, then $\mathrm{Num}(C)$
(\lref{GlobalTheory.T.numerical_point_statement_of_size}, which derives these inequalities from
the size relations). By \cref{pf:lift1:combine} with
$\Phi(C_1,C_2)=40(4C_1K^2+C_2K+2C_1)C_2$, \cref{pf:lift1:sumat} gives the sum point condition.
\end{proof}

\begin{theorem}[Reference exponents of the core]\label{pf:lift1:refcore}
The size condition and the sum point condition imply the reference core condition.
\formallabel{GlobalTheory.T.reference_core_of_size_and_sum_point}
\formallabel{GlobalTheory.RefSum.refCore_of_size_sumPt}
\end{theorem}

\begin{proof}
We combine (\cref{pf:lift1:combine}) the sum point condition (constant $C_1$) with two lift conditions
proved from the size condition: for $j<m$ and admissible $u_1$,
\begin{enumerate}[label=(\roman*)]
\item $r_+,r_-\ge1$, $r_++r_-=r$, $u_1$ is non-increasing with $u_{1,\tau,r_+-1}-u_{1,\tau,r_+}\ge n_c+1$,
  $|u_{1,\tau,l}|\le2b$, $(a(u_1),n_c)$ is dominant, $a(u_1)$ is positive at the last entries, and
  $G(u_1)=\mathrm{wt}_{n_c/2}(u_1+\mathrm{par}(\sigma_{u_1}))$
  (\lref{GlobalTheory.T.quotient_weight_gap_bounds'}, which states these properties);
\item the jet numerical condition of the development holds for $\mathcal D$, the dimension attached to
  $a(u_1)$ and the constant $\frac1{2500}$; $2r\le n_c$; and $T(\log g_0+\log g_1)\le C_3U$
  (\lref{GlobalTheory.T.jetNum_nCol_lenT_bounds_of_quo}, which states these bounds).
\end{enumerate}
Two estimates of the development apply to the core datum at level $j$, which satisfies the structural
laws, and to exponents $(a,n)$ which are dominant, positive at the last entries, with
$\sum_f2\liftonerk f\le n$, $2r\le n$ and the jet numerical condition: at the auxiliary primes,
$\epsilon_{g_k}(X^o_j,\mathrm{wt}(a,n))\le1$
(\lref{GlobalTheory.T.supLatticeReferenceExponent_thetaWeight_level_le_one}); and at every other
prime $q$, $\mathrm{Osc}_q(X^o_j,\gamma)\le b$ for the gap exponents $\gamma$ of $a$ implies
$\epsilon_q(X^o_j,\mathrm{wt}(a,n))\le b$
(\lref{GlobalTheory.T.baseLocusExponent_le_oscillation_nonTypePrime}). The second is the passage
from the oscillation of the depths at integral points to the reference exponents.

Now let $p,i,u_1$ be given, with $j<m$, and let $(S,e,s)$ be given by the sum point condition. Put
$S'=S\cup\{g_0,g_1\}$, $e'_q=e_q+1$ for $q\in\{g_0,g_1\}$ and $e'_q=e_q$ otherwise, and
$s'=s+\log g_0+\log g_1$. Then $(S',e',s')$ is a prime budget (the added primes contribute at most
$\log g_0+\log g_1$) and $Ts'\le(C_1+C_3)U$ by (ii). By (i) and (ii) the estimates apply to
$(a(u_1),n_c)$, whose gap exponents are $\gamma_{u_1}$. For $q\in\{g_0,g_1\}$,
$\epsilon_q(X^o_j,G(u_1))\le1\le e'_q$ since $e_q\ge0$; for other primes,
$\epsilon_q(X^o_j,G(u_1))\le e_q=e'_q$ by the oscillation bound of the sum point condition.
\end{proof}

\begin{lemma}[From the core to deep covers]\label{pf:lift1:reftheta}
The size condition and the reference core condition imply the reference theta condition.
\formallabel{GlobalTheory.T.refTheta_of_size_refCore}
\formallabel{GlobalTheory.RefSum.refTheta_of_refCore_alt}
\end{lemma}

\begin{proof}
For the core datum at level $j<m$ with its structural laws, a good deep cover $Y$ at level $j$ and
dominant exponents $(a,n)$, $\epsilon_q(Y,\mathrm{wt}(a,n))\le\epsilon_q(X^o_j,\mathrm{wt}(a,n))$ for
every prime $q$ (\lref{GlobalTheory.RefSum.refExp_vp_cover_le}, which compares the two suprema).
Combine the reference core condition with the lift condition (i) in the proof of \cref{pf:lift1:refcore},
which gives the dominance of $(a(u_1),n_c)$: the prime budget of the reference core condition at
$(p,i,u_1)$ serves for every good deep cover $Y$, with the same constant.
\end{proof}

\begin{lemma}[From the theta weight to the weight]\label{pf:lift1:reftow}
Hypothesis (Hdg), the size condition, the flat degree condition and the reference theta condition
imply the reference weight condition.
\formallabel{GlobalTheory.T.refW_of_prints_size_flatDeg_refTheta}
\formallabel{GlobalTheory.RefSum.ToWeight.refinedSum_weight_of_refinedTheta_alt}
\end{lemma}

\begin{proof}
Let $B_Y$ be the set of prime factors of $2d\,|d_E|\,M'$, $M'$ the level modulus of $Y$. We combine the
reference theta condition (constant $C_1$), the flat degree condition, the lift condition (i) in the
proof of \cref{pf:lift1:refcore}, the cover condition $\sum_{q\in B_Y}\log q\le C_2\,frL$
(\lref{GlobalTheory.BadLog.logSum_badPrimes_le_alt}, a bound for the logarithm of the level), and
the lift condition $T(frL)^{12}\le U$, $L\ge1$, $\log\liftonerd(E)\le3L$, $n_c\le20fr^2$, $f,r\ge1$
(\lref{GlobalTheory.T.lenT_zBase_unitE_nCol_bounds}, which states these size bounds).

The development proves, under (Hdg), for a good deep cover $Y$ with $r_+,r_-\ge1$ on which flat
sections exist, every weight $u$, every parallel weight $\mathrm{par}(w_0)$, $m\in\N$ and every prime
$q$:
\begin{equation}\label{pf:lift1:eq:refshift}
\epsilon_q(Y,\mathrm{wt}(u))\le\epsilon_q\bigl(Y,\mathrm{wt}_m(u+\mathrm{par}(w_0))\bigr)
+\Bigl[q\in B_Y\Bigr]\Bigl(m\,fr\,\frac{v_q(d_E)}{[E:\Q]}+1\Bigr)
\end{equation}
(\lref{GlobalTheory.T.reference_exponent_valuation_bound}). At $(p,i,Y,u_1)$, let $(S,e,s)$ be
given by the reference theta condition, $a=\frac{n_c}2fr$, and
\[
S'=S\cup B_Y,\quad e'_q=e_q+[q\in B_Y]\Bigl(a\frac{v_q(d_E)}{[E:\Q]}+1\Bigr),\quad
s'=s+a\log\liftonerd(E)+\sum_{q\in B_Y}\log q .
\]
Since $\sum_q v_q(d_E)\log q=\log|d_E|$ and $\log\liftonerd(E)=\log|d_E|/[E:\Q]$, $(S',e',s')$ is a prime
budget. As $\frac{n_c}2\le10fr^2$, $\log\liftonerd(E)\le3L$ and $frL\ge1$,
\[
a\log\liftonerd(E)+\sum_{q\in B_Y}\log q\le30f^2r^3L+C_2frL\le(30+C_2)(frL)^{12},
\]
so $Ts'\le(C_1+30+C_2)U$. Finally \eqref{pf:lift1:eq:refshift} with $u=u_1$, $w_0=\sigma_{u_1}$,
$m=n_c/2$ (flat sections exist by the flat degree condition, and $r_+,r_-\ge1$ by (i)) and
$G(u_1)=\mathrm{wt}_{n_c/2}(u_1+\mathrm{par}(\sigma_{u_1}))$ give
$\epsilon_q(Y,\mathrm{wt}(u_1))\le e_q+[q\in B_Y](a\,v_q(d_E)/[E:\Q]+1)=e'_q$.
\end{proof}

\begin{theorem}[Reference weights]\label{pf:lift1:refw}
Hypothesis (Hdg) and the deep goodness, size and flat degree conditions imply the reference weight
condition.
\formallabel{GlobalTheory.RefSum.refinedSum_weight_of_size_alt}
\end{theorem}

\begin{proof}
By \cref{pf:lift1:sumpt} the size condition gives the sum point condition; by \cref{pf:lift1:refcore}
the reference core condition; by \cref{pf:lift1:reftheta} the reference theta condition; and by
\cref{pf:lift1:reftow}, with (Hdg) and the flat degree condition, the reference weight condition. The
deep goodness condition is not used.
\end{proof}

\subsubsection*{The carrying section condition}
Let $\Lambda$ be the hermitian $\OO_E$-module of the datum $D_j$. A \emph{theta function of level $5$}
is a rule $f$ whose value at every framed point is the theta value with characteristics given by a
family $x=(x_1,\dots,x_r)$ of classes in $\Lambda/5^4\Lambda$ which generate it as an $\OO_E$-module; $f^2$
is a section of the weight $\mathrm{wt}(0,2r)$, and $|f^2|(z)$ denotes its pointwise norm at a frame $z$.
For a frame $z$ at a good embedding let $\lambda_1(z)$ be the first minimum of the hermitian lattice
of $z$, and put $\phi(z)=2f/(\liftonerd(E)^2\lambda_1(z)^2)$. Let $P_{16}=P_{16}(E,r,g_0g_1,m)$ be the
explicit constant of the theta constant estimates and $\vartheta(R_0)$ the explicit threshold attached
to $(E,g_0g_1,R_0)$ there; $M'$ is the level modulus of $Y$.

\begin{theorem}[Theta functions and the Bergman kernel]\label{pf:lift1:thetaberg}
Assume (Hdg), (H), and the deep goodness, size and flat degree conditions. Then the cover condition
defined by the following predicate holds: for every admissible $u_1$ there is $\beta\ge0$ with
$T\beta\le CU$ such that for every component $Z$, every theta function $f$ of level $5$, every good
embedding $\sigma$ and every frame $z$ of $Z$ at $\sigma$,
\[
|f^2|(z)\le e^\beta\qquad\text{and}\qquad|f^2|(z)^3\le e^\beta\,\mathrm{Berg}_Z(\mathrm{wt}(u_1))(z).
\]
\formallabel{GlobalTheory.CovHt.KeySection.theta_function_le_bergman_kernel_alt}
\end{theorem}

\begin{proof}
We combine three cover conditions, the first proved from the size condition and the other two from
all the hypotheses:
\begin{enumerate}[label=(\alph*)]
\item $r_-\le m$, $P_{16}\ge0$, $TP_{16}\le C_1U$, and for every integer $1\le D_F\le b^2$, with
  $R_0=\sqrt{90(fr+D_F)}$, the parameters of the theta constant estimate exist for $(r,D_F,g_0g_1,R_0,m)$,
  and every $s\ge0$ with $M's<b$ satisfies $\sqrt s<\vartheta(R_0)$
  (\lref{GlobalTheory.T.eta_bound_length_unit_alt}, which states these numerical facts);
\item the \emph{thin part}: at frames $z$ of $Z$ with $M'\lambda_1(z)^2<b$,
  $\mathrm{Berg}_Z(\mathrm{wt}(u_1))(z)\ge\exp(-\phi(z)/625)$
  (\lref{GlobalTheory.T.bergman_kernel_thin_part_estimate_alt});
\item the \emph{thick part}: there is $B\ge0$ with $TB\le C_3U$ such that
  $\mathrm{Berg}_Z(\mathrm{wt}(u_1))(z)\ge e^{-B}$ at frames $z$ of $Z$ with $M'\lambda_1(z)^2\ge b$
  (\lref{GlobalTheory.T.exp_neg_B_le_bergOn_of_la1sq_bound}; its proof constructs peak sections by
  H\"ormander's $L^2$ method, with a Poincar\'e-type $L^2$ inequality on a fundamental set).
\end{enumerate}
We also use the theta constant estimate
(\lref{GlobalTheory.T.normSq_thetaConstant_spanning_classes_le_exp}): at a good embedding with
unitary uniformisation, for a theta function $f$ of odd level $R>1$ attached to generating classes
modulo $R^4$, if the parameters of (a) exist for $(r,D_F,q,R_0,m)$ and $r_+\le m$ or $r_-\le m$, then
$|f^2|(z)\le\exp\bigl((P_{16}(E,r,q,m,R)-\phi_\vartheta(z))/R^4\bigr)$ at every frame $z$, where
$\phi_\vartheta(z)=\phi(z)$ if $\lambda_1(z)<\vartheta(R_0)$ and $\phi_\vartheta(z)=0$ otherwise. We apply it
with $R=5$ ($R^4=625$), $q=g_0g_1$, $D_F=1$; the uniformisations of $Y$ at good embeddings are unitary.

Put $\beta=P_{16}+B$, so $\beta\ge0$ and $T\beta\le(C_1+C_3)U$. Let $a=|f^2|(z)$ and
$G=\mathrm{Berg}_Z(\mathrm{wt}(u_1))(z)$. As $\phi_\vartheta\ge0$, $a\le e^{P_{16}/625}\le e^\beta$. For the
second inequality it suffices to find reals $s,t$ with $a\le e^s$, $3s\le P_{16}+B+t$ and $e^t\le G$, for
then $a^3\le e^{3s}\le e^{\beta}e^t\le e^\beta G$. If $M'\lambda_1(z)^2<b$, then
$\lambda_1(z)<\vartheta(R_0)$ by (a), so $\phi_\vartheta(z)=\phi(z)$; take $s=(P_{16}-\phi(z))/625$ and
$t=-\phi(z)/625$, using (b) and $3s\le P_{16}+t$ (as $P_{16},\phi(z)\ge0$). Otherwise take
$s=P_{16}/625$ and $t=-B$, using (c).
\end{proof}

\begin{theorem}[Carrying sections]\label{pf:lift1:carry}
Hypothesis (Hdg), (H) and the deep goodness, size and flat degree conditions imply the carrying section
condition.
\formallabel{GlobalTheory.CovHt.KeySection.cover_height_section_bound_alt}
\end{theorem}

\begin{proof}
We combine \cref{pf:lift1:thetaberg} (constant $C_1$) with the flat degree, size, host numerical
($M=1$) and weight numerical conditions. Let $(p,i,Y,u_1)$ be given and $\beta$ as in
\cref{pf:lift1:thetaberg}. The uniformisations of $Y$ satisfy the structural laws
(\lref{GlobalTheory.LawsDeep.laws_deep}); the level modulus $M'=g_0g_1\cdot15^{e}$ of $Y$ is divisible
by $5^4$ since $e\ge4$ (\lref{GlobalTheory.ThetaSteep.four_le_expL}); the size relations give
$T\cdot PP\le U$ and $(frL)^{12}\le PP$; and $|\liftonelogOm|\le3L$. Put $c_c=\beta/2+10PP\ge0$; then
$Tc_c\le(C_1+10)U$.

Let $Z$ be a component. By the flat degree property, for every $\tau\in\Psi$ there is a flat section on
$Z$; its degree $\liftoneflat(\liftonelogOm)$ is the degree of its frame part plus
$r_-(\log d-\log\liftonerd(E))$, so the frame part has degree $\frac{r_+-r_-}2\liftonelogOm$. The product over
$\tau\in\Psi$ of the $(12r)$-th powers of these frames is a section $S$ of the parallel weight $-12r$ with
constant norms and
$\deg S=12r\,f\,\frac{r_+-r_-}2\liftonelogOm$, as degrees of sections with constant norms are additive
The development then constructs the carrying section
(\lref{GlobalTheory.CovHt.KeySection.exists_section}): if $5^4\mid M'$ and every theta function $f$
of level $5$ satisfies $|f^2|(z)^3\le e^\beta\mathrm{Berg}_Z(\mathrm{wt}(u_1))(z)$ at the frames of $Z$,
then there are $\theta,e,t$ with the properties required in the carrying section condition and
\[
\mathrm{cost}(e,t)=\frac\beta2+\frac{3r\cdot2fr\log(2\pi)}8-\frac{\deg S}8 ;
\]
here $\theta=(f^2)^{12}KS$ for a theta function $f$ of level $5$ which does not vanish identically on
$Z$ and a constant section $K$ of parallel weight with constant norms, so that $\theta$ has exponent
$24r\,h=8(u_1-u_1')$, and the archimedean bound follows from the bound on $|f^2|^3$. Finally, with
$\log(2\pi)\le7$ and $|\frac{r_+-r_-}2\liftonelogOm|\le\frac r2\cdot3L$,
\[
\frac{3r\cdot2fr\log(2\pi)}8-\frac{12rf\frac{r_+-r_-}2\liftonelogOm}8
\le\frac{21}4fr^2+\frac94fr^2L\le10(frL)^{12}\le10PP
\]
(as $f,r,L\ge1$), so $\mathrm{cost}(e,t)\le c_c$.
\end{proof}

\begin{lemma}[Height comparison]\label{pf:lift1:htcmp}
Assume (Hdg). For every CM field $E$ which is Galois over $\Q$, with a host, a datum, a level and a
moduli space $X^o$ over a coefficient field $F$, the height comparison property holds for $X^o$.
\formallabel{GlobalTheory.CovHt.cover_height_comparison}
\end{lemma}

\begin{proof}
This is \lref{GlobalTheory.CovHt.FrameCmp.htCmp}, which proves the height comparison property for
every such moduli space directly from the definitions of the two families of valuations. Hypothesis
(Hdg) is not used.
\end{proof}

\begin{lemma}[Error budget]\label{pf:lift1:budget}
The host numerical conditions (for every $M$), the weight numerical condition and the size condition
imply the budget condition; the proof gives the constant $83400$.
\formallabel{GlobalTheory.CovHt.BudgetPf.budget_alt}
\end{lemma}

\begin{proof}
By \cref{pf:lift1:combine} it suffices to prove $\liftoneB(PP,83400,u_1)$ at points where
$\liftoneNh(PP,A,1)$ and $\liftoneNw$ hold. This is a computation with the sizes
(\lref{GlobalTheory.CovHt.BudgetPf.budgetAt}, which carries it out): the level modulus of a good
deep cover is $g_0g_1\cdot15^{e}$, which has at most four prime factors; for the gap $g$ supplied by the
construction of $u_1$ one has $\mathrm{WtGap}(u_1',g)$ and $\log\liftonebnd(u_1',g)\le L+22$;
$\ell(\mathcal D)$ and $\log\liftonerd(E)$ are $O(L)$; $\liftonehtc(B)$ is bounded using that the primes
in $P_B$ are at most $B$ or divide $2|d_E|M'$; and the resulting bounds for $T\liftoneerr$ and
$T\liftonehtc(B)|u_1'|_1$ are at most $83400\,U$ by the weight relations. The size condition is not
used.
\end{proof}

\subsubsection*{The theorem}

\begin{proposition}[Heights from mean slopes]\label{pf:lift1:meanall}
Assume (V) and (Hdg), the deep goodness, size, cover degree, short, geometric, multiplicity and budget
conditions, and the conclusion of $\sfname{DeepStR}$: there is $C_0\ge0$ such that for every global datum
$\mathcal D$, every core datum $T$ of $\mathcal D$ with the structural laws and $m^2\le k$, every
$1\le j<m$ and every deep cover $Y$ of $T$ at level $j$ which is good for some parameters and whose level
modulus has at most $15$ prime factors, $\mathrm{MeanAll}(\ell(\mathcal D),C_0)$ holds on $Y$. Then the
height condition holds.
\formallabel{GlobalTheory.CovHt.HtTopG.height_of_meanAll_bound}
\end{proposition}

\begin{proof}
Let $C_1,\dots,C_7$ be the constants of the seven conditions, in the order listed; we show that the
height condition holds with constant $C_4+C_6+C_7+C_0C_7$, and by \cref{pf:lift1:combine} it suffices
to do this pointwise. Let $p,i$ with $i+1<m$, a good deep cover $Y$ and an admissible $u_1$ be given.
Deep goodness gives that $\mathrm{Comp}$ is nonempty; the size condition gives $U\ge b\ge1$ and
$\rho(u_1)<1$; and $Y$ is good with parameters $(e,o)$. The budget condition gives that the level modulus
of $Y$ has at most $15$ prime factors, so the hypothesis applies to the core datum $T$ of $p$ (which
satisfies the structural laws, and $m^2\le k$ by the choice of parameters of the family), the level
$j=i+1\in[1,m)$ and $Y$, and gives $\mathrm{MeanAll}(\ell(\mathcal D),C_0)$ on $Y$. The height comparison
property holds for $Y$ by \cref{pf:lift1:htcmp}. \Cref{pf:lift1:heightpt} with
$(C_s,C_m,C_b)=(C_4,C_6,C_7)$ gives $\liftoneH(PP,C_4+C_6+C_7+C_0C_7,u_1)$.
\end{proof}

\begin{theorem}[Height bound]\label{pf:lift1:height-thm}
Assume (V), (X), (GM), (Q), (Hdg), (H) and $\sfname{DeepStR}$, and the deep goodness, size, cover
degree and flat degree conditions. Then the height condition holds.
\formallabel{SelfIntersectionVolumeEstimates.height_lower_bound_general_alt}
\end{theorem}

\begin{proof}
The short condition holds by \cref{pf:lift1:short}; the geometric condition by \cref{pf:lift1:geomcond};
the multiplicity condition by \cref{pf:lift1:mult}, \cref{pf:lift1:refw} and \cref{pf:lift1:carry}; and the budget
condition by \cref{pf:lift1:budget} and the numerical conditions of \cref{pf:lift1:num}. As
$\sfname{DeepStR}$ is the hypothesis on mean slopes in \cref{pf:lift1:meanall}, that proposition gives the
height condition.
\end{proof}

% ---------------------------------------------------------------------------
\subsection{The Hilbert--Samuel lower bound}\label{pf:lift1:hs}

We keep the notation of \cref{pf:lift1:height}; in particular $j=i+1<m$, $Y$ is a good deep cover of
level $j$ and $u_1$ is admissible, so that $u_1$ is the weight of the parameter $P=(u_W,c_0,k-2,T)$ and
$u_M=\mathrm{wt}_T+w$ with $w$ parallel. Put $z_L=frL$.

\begin{definition}\label{pf:lift1:hsat}
$\liftoneR(PP,u_1)$ means: for every component $Z$,
$B_T\liftonedg(u_1,Z)\le\liftonerkC(u_M,Z)$ and $(1-\zeta(u_1))\liftonerkC(u_M,Z)\le B_T\liftonedg(u_1,Z)$.
$\liftoneHS(PP,C,u_1)$ means: there are $A>0$ and $R\ge0$ with $A\log\max(A,1)+R\le CU$ such that for
every component $Z$
\[
B_T\liftonedg(u_1,Z)\Bigl(T\liftonehgt(u_1,Z)-3(\liftonedfl+1)|\liftonehgt(u_1,Z)|-A\log\bigl(2+|\liftonehgt(u_1,Z)|\bigr)-R\Bigr)
\le\liftonedegC(u_M,Z).
\]
The \emph{share condition} and the \emph{cover height sum condition} are the cover conditions defined
by requiring these predicates for every admissible $u_1$.
\end{definition}

\begin{theorem}[Hilbert--Samuel lower bound]\label{pf:lift1:hs-thm}
Assume (X), (Hdg), (GM),
(Q), (Sieg), (Tor) and (Lie), and the deep goodness, size, flat degree, cover degree and share conditions. Then the cover
height sum condition holds.
\formallabel{GlobalTheory.T.coverHeightSum_of_flatDeg_coverDegree_share}
\formallabel{GlobalTheory.CovTop.ShortRoad.hilbertSamuel_of_component_estimates_alt}
\end{theorem}

The proof rests on a Hilbert--Samuel estimate on each component, in the following form.

\begin{definition}[Hilbert--Samuel claim]\label{pf:lift1:claim}
Let $C\in\R$ and $T,D\in\N$, and let $n$, $j_0$, $\Lambda$, $\beta$, $\Lambda_D$, $\Lambda_T$, $J$,
$\delta$, $h$, $d$ and $\deg$ be reals. Put $B=\binom{T-1}D$, $Z_s=1-B\delta/d$,
$\ell_\rho=\log\Lambda+C\log\bigl(nDj_0\delta(2+|h|+\log\Lambda)\bigr)$ and
$g=\frac12(D+1)(\log n+1)+1$. The claim
\[
\mathrm{HS}_C(T,D;\,n,j_0,\Lambda,\beta,\Lambda_D,\Lambda_T,J;\,\delta,h,d,\deg)
\]
is the statement: $0\le Z_s\le1$ and
\begin{multline*}
(1-Z_s)\Bigl(Th-3(D+1)|h|-J-CD^2\log\bigl(e^{\ell_\rho}T\bigr)\\
-(T+D+1)\log\beta-(T+D^2)(g+2)\Bigr)\\
-Z_s\Bigl(T\log\bigl(2e^{\ell_\rho}\sqrt\beta\bigr)+\log\Lambda_D+\log\Lambda_T\Bigr)\le\frac{\deg}d .
\end{multline*}
\end{definition}

In the applications $\delta=\liftonedg(u_1,Z)$, $h=\liftonehgt(u_1,Z)$, $d=\liftonerkC(\mathrm{wt}_T,Z)$ and
$\deg=\liftonedegC(\mathrm{wt}_T,Z)$; thus the claim bounds the normalised slope $\deg/d$ of the forms of
weight $\mathrm{wt}_T$ on $Z$ from below by $T$ times the height of $Z$ for $u_1$, up to explicit errors,
the fraction $Z_s$ of the rank not accounted for by $B\delta$ being estimated separately.

\subsubsection*{The abstract Hilbert--Samuel estimate}
The claim is deduced from an estimate for abstract data, which we now describe. Let $F$ be a number
field. For integers $B,N$, a polynomial in the variables $X_{i,k}$ ($i<B$, $k<N$) is
\emph{multihomogeneous} of degree $a\in\N^B$ if it is homogeneous of degree $a_i$ in each block
$(X_{i,k})_k$. Given a matrix $w\in F^{N\times n}$ whose columns span $F^N$ and invertible complex
matrices $A_\sigma$ ($\sigma\colon F\to\C$), defining the frame norms $\|u\|_{A_\sigma}$ on $\C^N$, every
nonzero multihomogeneous $G$ has a height $\mathrm{ht}(G)=\mathrm{ht}_{w,A}(G)\in\R$ (the development's
height of multiforms for these hermitian data); it is additive on products and vanishes on nonzero
constants (\lref{GlobalTheory.HSShort.htMul}), and if the frame norms of the columns of $w$ are at
most $\Lambda\ge1$, then $\mathrm{ht}(G)\ge-c(n,\Lambda)\sum_ia_i$ with
$c(n,\Lambda)=\log(n\Lambda)+\frac12(\log n+1)$ (\lref{GlobalTheory.T.multiformHeight_lowerBound}).

\begin{definition}[Hilbert--Samuel data]\label{pf:lift1:hsdata}
A \emph{Hilbert--Samuel datum} $X$ over $F$ consists of: a smooth projective geometrically integral
$F$-variety $Y$ of dimension $D$ with a base point free linear system given by sections
$b_1,\dots,b_N$ of a line bundle $L$, a proper closed boundary subset with open complement $Y^o$, and a
section of a power of $L$; hermitian data $(w,A)$ as above and a real $\Lambda$; the same kind of data for
a second line bundle $L_2$; the canonical bundle of $Y$ and a line bundle $L_A$ with a section of
$L_A\otimes(L_2\otimes L^{D})^{-1}$ not vanishing on $Y^o$; an integer $T\ge D+1$; for every $j\in\N$ an
$F$-space $V_j$ of sections of degree $j$, containing the global sections for $j\le T$ and equal to them for
$j=0$, with a full $\OO_F$-lattice $M_j$ and a conjugation invariant hermitian family; and reals
$\beta,\Lambda_D,\Lambda_T\ge1$. Put $\delta_X=(L^D)$ (top self-intersection), $d_X=\liftonerk M_T$ and
$\deg_X=\deg M_T$. A real $h$ is a \emph{height} of $X$ if there are $\delta\ge1$ and a Chow form $f$ of
degree $\delta$ of the image of $Y$ under $(b_1,\dots,b_N)$ such that $(D+1)\delta h$ is the height of $f$
normalised with $\delta$. A \emph{linear section} of $Y$ by $s\in E^{D\times N}$ ($E\supseteq F$) is the
common zero set of the sections $\sum_ks_{i,k}b_k$, $i<D$; it is \emph{good} if it is finite \'etale and
contained in $Y^o$.

For $\Lambda_0\in\R$ the \emph{geometric hypothesis} $\mathrm{GH}(X,\Lambda_0)$ means: $F$ is totally
complex; the linear system is birational onto its image; $L_A$ and $L_2$ are nef and big; multiplication
by the elements of the lattice spanned by the columns of $w$ maps $M_j$ into $M_{j+1}$ and multiplies the
norms by at most $\sqrt\beta$ times the frame norm ($j<T$); the frame norms of the columns of $w$ are at
most $\Lambda$; $\Lambda,\Lambda_0\ge1$; and $M_0$, $M_D$, $M_T$ contain bases of $V_0$, $V_D$, $V_T$ whose
elements have norms at most $\Lambda_0$, $\Lambda_D$, $\Lambda_T$. The \emph{geometric bound} is the
statement: for every Hilbert--Samuel datum $X$ over a number field, every $\Lambda_0$ with
$\mathrm{GH}(X,\Lambda_0)$ and every height $h$ of $X$,
\[
\mathrm{HS}_{26}\bigl(T,D;\,n,1,\Lambda,\beta,\Lambda_D,\Lambda_T,\log\Lambda_0;\,\delta_X,h,d_X,\deg_X\bigr).
\]
\end{definition}

The proof of the geometric bound passes through hermitian systems, which isolate the linear algebra.

\begin{definition}[Hermitian systems]\label{pf:lift1:hsys}
A \emph{hermitian system} $S$ over $F$ consists of an $F$-vector space $A$ with commuting endomorphisms
$m(u)$ depending linearly on $u\in F^N$, hermitian data $(w,A_\sigma)$ with $w\in F^{N\times n}$, integers
$1\le D$, $D+1\le T$, subspaces $V_j\subseteq A$ with full lattices $M_j$ ($j\le T$) and hermitian
families, and reals $\Lambda,\beta,\Lambda_0,\Lambda_D,\Lambda_T\ge1$, such that: the columns of $w$ span
$F^N$ and have frame norms at most $\Lambda$; $m(w_k)M_j\subseteq M_{j+1}$ for $j<T$ and all columns
$w_k$; multiplication from $V_j$ to $V_{j+1}$ increases the norms (of tensors $\sum_{k,i}u_kc_i\,m(e_k)f_i$)
by at most $\sqrt\beta\|u\|_{A_\sigma}$; and $M_0,M_D,M_T$ contain bases of $V_0,V_D,V_T$ with norms at most
$\Lambda_0,\Lambda_D,\Lambda_T$. For $s\in F^{D\times N}$ and $\alpha\in\N^D$ put
$s^\alpha=\prod_im(s_i)^{\alpha_i}$. For $\delta\ge1$ and $f$ as below, $S$ is \emph{flat} at $s$ if, over
some extension $E$ of $F$, there are points $y_1,\dots,y_\delta\in E^N\setminus0$ with
$f(s,l)=c\prod_p\langle l,y_p\rangle$ ($c\ne0$), and a decreasing filtration $(U_a)_{a\le T-D}$ of
$A\otimes E$ with injective symbol maps from $U_a/U_{a+1}$ to the functions on pairs $(\alpha,p)$ of a
monomial of degree $a$ in $D$ variables and an index $p\le\delta$, under which $s^\alpha$ followed by
multiplication by linear forms $l_1,\dots,l_{T-a}$ acts on symbols by $\prod_q\langle l_q,y_p\rangle$, and
whose deepest step $U_{T-D}$ is the span of the $s^\alpha V_D$, $|\alpha|=T-D$ (the
\emph{Hilbert--Samuel filtration} at $s$). With $c=c(n,\Lambda)$ and $Z_s=1-\binom{T-1}D\delta/d$, the
\emph{slope claim} $\mathrm{Cl}(\delta,h,d,\deg)$ is: $0\le Z_s\le1$ and
\[
(1-Z_s)\Bigl(Th-D(D+1)c-\log\Lambda_0-\frac T2\log\beta\Bigr)
-Z_s\Bigl(Tc+\frac T2\log\beta+\log\Lambda_D+\log\Lambda_T\Bigr)\le\frac{\deg}d .
\]
\end{definition}

\begin{theorem}[Slope of the top lattice]\label{pf:lift1:slope}
Let $S$ be a hermitian system, $\delta\ge1$, $f$ an irreducible polynomial in $X_{i,k}$ ($i\le D$, $k<N$),
multihomogeneous of degree $\delta$ in each of the $D+1$ blocks, and $Q\ne0$ a polynomial in $X_{i,k}$
($i<D$), such that $S$ is flat at every $s\in F^{D\times N}$ with $Q(s)\ne0$. Then $d=\dim V_T>0$ and
$\mathrm{Cl}\bigl(\delta,\mathrm{ht}(f)/((D+1)\delta),d,\deg M_T\bigr)$ holds.
\formallabel{GlobalTheory.T.slopeLowerBound_topLattice_offHypersurface}
\formallabel{GlobalTheory.HSShort.SlopeTop.slopeBound}
\end{theorem}

\begin{proof}
Let $b$ be the basis of $V_T$ in $M_T$ with norms at most $\Lambda_T$, $d=\dim V_T$, and extend its
coordinate functions to linear forms $c_i$ on $A$. Call $s$ \emph{good} if $Q(s)\ne0$; good $s$ exist since
$Q\ne0$ and $F$ is infinite. For good $s$ the \emph{deep rank}
$\dim\sum_{|\alpha|=T-D}s^\alpha V_D$ is at most $d$
(\lref{GlobalTheory.T.finrank_deepestFiltrationStep_eq_rational}, which identifies it with the
dimension of the deepest step of the filtration); let $y\le d$ be its maximum over good $s$, attained at
$s_0$.

\emph{Columns.} A \emph{column} is a pair of an element $x\in V_j$ ($j\in\{0,D,T\}$) and an assignment of
the $t=T-j$ remaining multiplications to blocks $0,\dots,D$; it defines the polynomial vector
$(c_i(\text{the product of }x\text{ with the corresponding linear forms}))_i$, and for a family of $d$ columns
let $\Delta\in F[X_{i,k}:i\le D]$ be the determinant of these vectors. Using the flatness at $s_0$, the
development chooses $d$ columns with short vectors ($x\in M_j$ with norm at most $\Lambda_j$) and
$\Delta\ne0$, of which $\binom{T-1}D\delta$ have $j=0$ (indexed by the pairs of a monomial of degree
$a<T-D$ in $D$ variables and an index $p\le\delta$), $y$ have $j=D$ and the others $j=T$
(\lref{GlobalTheory.T.exists_short_columns_det_ne_zero_hilbertSamuel}).

\emph{Divisibility.} Put $e=\sum_{a<T-D}\binom{a+D-1}{D-1}(T-a)$. At every good $s$, the specialisation of
$\Delta$ at $s$ (as a polynomial in the last block) is a constant times $(\prod_p\langle l,y_p\rangle)^e$
(\lref{GlobalTheory.T.det_rational_sections_eq_pow_product_linearForms}, which uses that the deep
rank at $s$ is at most $y$), hence a constant times $f(s,l)^e$. As $f$ is irreducible and multihomogeneous
and this holds off the hypersurface $Q=0$, $f^e$ divides $\Delta$
(\lref{GlobalTheory.HSShortDiv.pow_dvd_of_specialisations}). Write $\Delta=f^eG$. $\Delta$ is
multihomogeneous of a degree $a$ whose total is $\sum_kt_k$
(\lref{GlobalTheory.HSShort.det_isMultiHomogeneous}), so $G$ is multihomogeneous of a degree $a'$ with
$a=e\delta\mathbf 1+a'$ (a factor of a nonzero multihomogeneous polynomial is multihomogeneous).

\emph{Heights.} By additivity, $\mathrm{ht}(\Delta)=e\,\mathrm{ht}(f)+\mathrm{ht}(G)$, and
$\mathrm{ht}(f)\ge-c(D+1)\delta$, $\mathrm{ht}(G)\ge-c\sum_ia'_i$. The coefficients of the column vectors
lie in $M_T$ (\lref{GlobalTheory.T.coeff_linearSubst_column_mem_integralLattice}), and on the
products of unit spheres of the frame norms the $k$-th column has norm at most
$B_k=\sqrt\beta^{\,t_k}\Lambda_{j_k}$ (\lref{GlobalTheory.T.hermFamComplexTensorNorm_column_unitSpheres_le}
and the norm bound of $S$). A lower bound of the degree of a lattice by the height of a determinant of
polynomial vectors (\lref{GlobalTheory.T.lattice_degree_ge_height_polynomialVectors_det}) gives
\[
\mathrm{ht}(\Delta)-\sum_k\log B_k\le\deg M_T .
\]
\emph{Counting.} The columns with $j=0$, $D$, $T$ have $t=T$, $T-D$, $0$, so
$e(D+1)\delta+\sum_ia'_i=\binom{T-1}D\delta T+y(T-D)$, and
$(D+1)e=\binom{T-1}D(T+D(D+1))$ (elementary binomial identities, from
$\#\{\alpha:|\alpha|=a\}=\binom{a+D-1}{D-1}$ and $\sum_{a<T-D}\binom{a+D-1}{D-1}=\binom{T-1}D$). Further
$\binom{T-1}D\delta+y\le d$ and
$\sum_k\log B_k=\binom{T-1}D\delta\log(\sqrt\beta^T\Lambda_0)+y\log(\sqrt\beta^{T-D}\Lambda_D)
+(d-\binom{T-1}D\delta-y)\log\Lambda_T$. An elementary computation
(\lref{GlobalTheory.HSShortEnd.claimS_of_numbers}, which derives the slope claim from these relations,
the two height lower bounds and the degree bound) gives $d>0$ and
$\mathrm{Cl}(\delta,\mathrm{ht}(f)/((D+1)\delta),d,\deg M_T)$.
\end{proof}

\begin{theorem}[Bertini over number fields]\label{pf:lift1:bertini}
Assume (Q). Let the data $Y$, $(b_k)$, $(w,A)$ of a Hilbert--Samuel datum over $F$ be given (with the
boundary and the section of a power of $L$), with birational linear system, and let $f$ be a Chow form of
degree $\delta$ of the image. Then $\delta=(L^D)$, and there is a nonzero polynomial $Q$ in $X_{i,k}$
($i<D$, $k<N$) such that for every number field $E\supseteq F$ and every $s\in E^{D\times N}$ with
$Q(s)\ne0$: the linear section of $Y$ by $s$ is good, has exactly $\delta$ geometric points, the linear
system is injective on them, and $f(s,\cdot)$ is a nonzero constant times the product of the linear forms
$l\mapsto\sum_kl_kb_k(x)$ over these points $x$.
\formallabel{GlobalTheory.T.bertiniNumberField_of_hilbertSamuel₂}
\formallabel{GlobalTheory.HSShort.bertiniNumberField_of_hilbertSamuel}
\end{theorem}

\begin{proof}
We prove the statement for every field $E\supseteq F$; number fields are a special case. Under (Q), the image of $Y$ has degree $(L^D)$ and dimension $D$
(\lref{GlobalTheory.ImageDegree.degree_image_eq_topSelf}); its ideal is prime and homogeneous, and the
degree of a Chow form of a prime homogeneous ideal is its Hilbert degree
(\lref{GlobalTheory.ChowExists.degree_eq_hilbertDegree}); hence $\delta=(L^D)$. There
are forms $g_1,g_2$ outside the ideal of the image such that $g_1$ vanishes on the boundary and the
linear system has finite fibres off $g_2=0$
(\lref{GlobalTheory.T.exists_form_outside_imageIdeal_vanishing_on_boundary},
\lref{GlobalTheory.T.exists_form_finite_fibres_birational_linearSystem}); an elimination form
$Q_1\ne0$ such that, if $Q_1(s)\ne0$, no point of the image in the zero set of $g_1g_2$ lies on the linear
space defined by $s$ (\lref{GlobalTheory.T.exists_eliminationForm_generalLinearForms_cone}); and a
discriminant $Q_2\ne0$ of pencils of lines such that $f(s,\cdot)$ has no repeated linear factor if
$Q_2(s)\ne0$ (\lref{GlobalTheory.PencilDisc.exists_poly_noSquare}, a Bertini-type argument with the
discriminant of a generic pencil). Put $Q=Q_1Q_2$. If $Q(s)\ne0$, the linear section is finite and
contained in $Y^o$ and in the locus of finite fibres
(\lref{GlobalTheory.T.linearSection_finite_inOpen_off_badForm}), and since $f(s,\cdot)$ is
squarefree the section is \'etale with the stated properties
(\lref{GlobalTheory.T.linearSection_etale_of_squarefree_chowForm}).
\end{proof}

\begin{theorem}[Explicit arithmetic Hilbert--Samuel bound]\label{pf:lift1:geobound}
Hypothesis (Q), the conclusion of \cref{pf:lift1:bertini}, the height comparison stated in the proof
below, and the conclusion of \cref{pf:lift1:slope} imply the geometric bound.
\formallabel{GlobalTheory.T.explicit_arithmetic_hilbertSamuel_bound}
\formallabel{GlobalTheory.HSShortTop.geoBound_of_abstract}
\end{theorem}

\begin{proof}
Let $X$ be a Hilbert--Samuel datum with $\mathrm{GH}(X,\Lambda_0)$ and $h$ a height of $X$, defined by a Chow
form $f$ of degree $\delta$. By \cref{pf:lift1:bertini}, $\delta=\delta_X$ and there is $Q\ne0$ such that the
linear sections by $s$ with $Q(s)\ne0$ (over number fields) are good and simple. Let $S$ be the hermitian
system with $A$ the space of all sections, $m$ the multiplication by the linear combinations of the
$b_k$, the data $V_j,M_j$, norms and constants of $X$, and $\Lambda_0$; it satisfies the hypotheses of
\cref{pf:lift1:hsys}: each of them is a condition of $\mathrm{GH}(X,\Lambda_0)$ or of the definition of a
Hilbert--Samuel datum (and $N\ge1$, as the linear system is base point free). At every $s$ with $Q(s)\ne0$, $S$ admits the
Hilbert--Samuel filtration, i.e.\ is flat at $s$
(\lref{GlobalTheory.T.hilbertSamuel_filtration_at_good_sections}, proved from (Q), the geometric hypothesis
and the simplicity of the linear sections given by \cref{pf:lift1:bertini}). The Chow form $f$ is irreducible, since it is
irreducible after an extension of scalars (by definition of Chow forms) and the units of polynomial rings
are the nonzero constants, and it is multihomogeneous of degree $\delta$.
\Cref{pf:lift1:slope} gives $d_X>0$ and $\mathrm{Cl}(\delta,h',d_X,\deg_X)$ with
$h'=\mathrm{ht}(f)/((D+1)\delta)$, as $\dim V_T=d_X$ and the degree of $M_T$ is $\deg_X$ (the
inclusion of $V_T$ into the space of all sections is isometric, and degrees are invariant under isometries).

\emph{Height comparison.} For the height $h$ defined by $f$, $h-\frac12(\log n+1)\le h'$
(\lref{GlobalTheory.HeightTieProof.heightTie}, which shows
$h'=h-\frac12\sum_{k=1}^{N-1}\frac1k$ by the additivity of heights).

\emph{From the slope claim to the Hilbert--Samuel claim.} Let $\ell_\rho=\log\Lambda+26\log(nD\delta(2+|h|+\log\Lambda))$ and
$g=\frac12(D+1)(\log n+1)+1$. As $nD\delta(2+|h|+\log\Lambda)\ge2n$, $0\le c\le\ell_\rho$. Then
\begin{multline*}
Th-3(D+1)|h|-\log\Lambda_0-26D^2(\ell_\rho+\log T)-(T+D+1)\log\beta-(T+D^2)(g+2)\\
\le Th'-D(D+1)c-\log\Lambda_0-\tfrac T2\log\beta,
\end{multline*}
using $Th-\frac T2(\log n+1)\le Th'$, $D(D+1)c\le2D^2\ell_\rho$, $\frac T2\log\beta\le(T+D+1)\log\beta$
and $\frac T2(\log n+1)\le(T+D^2)(g+2)$; and
$Tc+\frac T2\log\beta+\log\Lambda_D+\log\Lambda_T\le T\log(2e^{\ell_\rho}\sqrt\beta)+\log\Lambda_D+\log\Lambda_T$.
Multiplying by $1-Z_s\ge0$ and $Z_s\ge0$ and inserting in $\mathrm{Cl}$ gives
$\mathrm{HS}_{26}(T,D;n,1,\Lambda,\beta,\Lambda_D,\Lambda_T,\log\Lambda_0;\delta_X,h,d_X,\deg_X)$.
\end{proof}

\begin{corollary}\label{pf:lift1:geobound-cor}
Hypothesis (Q) implies the geometric bound.
\formallabel{GlobalTheory.T.arithmeticHilbertSamuelBound}
\formallabel{GlobalTheory.HSShort.geoBound_of_quoted'}
\end{corollary}

\begin{proof}
\Cref{pf:lift1:bertini} holds under (Q); the height comparison holds unconditionally
(\lref{GlobalTheory.HeightTieProof.heightTie} with \lref{GlobalTheory.HSShort.htMul}); and
\cref{pf:lift1:slope} holds unconditionally. Apply \cref{pf:lift1:geobound}.
\end{proof}

\begin{theorem}[Hilbert--Samuel claim at a component]\label{pf:lift1:hscomp}
Assume (Q). Let $X^o$ be a unitary moduli space over $F$ with a model, a flag cover of the model, a
component $Z$, uniformisations, a weight parameter $P$ and the auxiliary choices of the development (a
modification of the flag cover adapted to $Z$ and $P$, a canonical identification, and the readiness
data), cusp frames $f_1$ (of size $n$) and $f_2$ on $Z$ for $u_1$ and $u_2$, a section $\theta$ and
constants $(\Lambda,\beta,\Lambda_D,\Lambda_T)$. Let $X_Z$ be the associated Hilbert--Samuel datum
(the variety being the closure of $Z$ in the modification, $L$ the line bundle of $u_1$ with the frame
$f_1$, and $V_j$ the forms of the weights $\mathrm{wt}_j$), and let $\Lambda_0$ satisfy
$\mathrm{GH}(X_Z,\Lambda_0)$. Then
\[
\mathrm{HS}_{26}\bigl(T,\liftonedfl;\,n,1,\Lambda,\beta,\Lambda_D,\Lambda_T,\log\Lambda_0;\,
\liftonedg(u_1,Z),\liftonehgt(u_1,Z),\dim\liftonecusp_Z(\mathrm{wt}_T),\liftonedegC(\mathrm{wt}_T,Z)\bigr).
\]
\formallabel{GlobalTheory.T.hilbertSamuel_component_short_basis_bound}
\formallabel{GlobalTheory.HSShort.hsClaimS_component_geo}
\end{theorem}

\begin{proof}
The datum $X_Z$ presents $Z$: its dimension is $\liftonedfl$ and its linear system is the map defined by
the cusp forms of weight $u_1$ (\lref{GlobalTheory.CompHS.compPresents_data}, which states this). Under
(Q), the degree $\delta_{X_Z}$ equals the degree $\liftonedg(u_1,Z)$ of the image
(\lref{GlobalTheory.HSCore.ImageDegree.imageDeg_eq_del}), and, $F$ being totally complex,
$\liftonehgt(u_1,Z)$ is a height of $X_Z$ (\lref{GlobalTheory.T.height_birationalImage_eq_tuple_height},
which identifies the height of the image computed from the lattice and norms of the forms with the height
of a Chow form). \Cref{pf:lift1:geobound-cor}, applied to $X_Z$, gives the claim with $\delta_{X_Z}$, $h=\liftonehgt(u_1,Z)$,
$d_{X_Z}$ and $\deg_{X_Z}$. Finally $\deg_{X_Z}=\deg M_T$ is the degree $\liftonedegC(\mathrm{wt}_T,Z)$ of the
lattice of cusp forms of weight $\mathrm{wt}_T$ on $Z$ (\lref{GlobalTheory.T.degree_supLattice_component_eq}),
$d_{X_Z}=\dim\liftonecusp_Z(\mathrm{wt}_T)$, and the parameters $T$, $D$, $n$, $\Lambda$, $\beta$,
$\Lambda_D$, $\Lambda_T$ of $X_Z$ are those of the statement, by construction of $X_Z$.
\end{proof}

\subsubsection*{The estimate at a component}

\begin{proposition}[Hilbert--Samuel estimate at a component]\label{pf:lift1:hslevel}
Assume the part (Hdg') of (Hdg) consisting of seventeen of its statements (among them the existence of
toroidal compactifications, Mumford's proportionality principle and the extension of homomorphisms of
semi-abelian schemes), (GM), (Q), Nadel vanishing for non-klt divisors, principalisation, (Sieg),
toroidal charts for fans and (Lie). Let $Y$ be good, $2<[E:\Q]/2$, $Z$ a component and $P$ a weight
parameter satisfying the largeness condition of the development, with the forms of weights
$\mathrm{wt}_t$ ($t\le T$) good.
Let $PP,\kappa\ge0$; let $f_2$ and $f_1$ be cusp frames on $Z$ for $u_2$ and $u_1$ ($f_1$ of size $n$) and
$\Lambda,\beta\ge1$ such that the frame norms of the columns of $f_1$ are at most $\Lambda$ and the
Bergman kernel of $u_1$ on $Z$ is bounded by $\beta$; assume that the ratios of the forms of weights
$u_1$ and $u_2$ generate the function field of $Z$; that $\liftonelat_Z(\mathrm{wt}_0)$ is spanned by
elements of norm $\le\exp(\liftonedfl PP(\kappa/T+1))$; and that every $\liftonelat_Z(\mathrm{wt}_t)$, $t\le T$,
is spanned by elements of norm $\le e^{PP\kappa}$. Then $n\ge1$ and
\begin{multline*}
\mathrm{HS}_{26}\Bigl(T,\liftonedfl;\,n,1,\Lambda,\beta,e^{PP\kappa},e^{PP\kappa},\liftonedfl PP(\tfrac\kappa T+1);\\
\liftonedg(u_1,Z),\liftonehgt(u_1,Z),\liftonerkC(\mathrm{wt}_T,Z),\liftonedegC(\mathrm{wt}_T,Z)\Bigr).
\end{multline*}
\formallabel{GlobalTheory.HSShortRun.hilbertSamuel_estimate_at_level}
\end{proposition}

\begin{proof}
The development provides a toroidal flag model of the moduli space of $Y$, a modification of it and a
hermitian line bundle $\omega$ whose sections realise the forms of the weights $\mathrm{wt}_t$
(\lref{GlobalTheory.T.exists_weight_hodge_bundle_generated}); ratios of cusp forms generate the
function field at the generic point of the flag space
(\lref{GlobalTheory.T.cuspForm_ratios_generate_at_flagSpace_genericPoint}); the spanning families
give short families of sections (\lref{GlobalTheory.HSApp.shortFamily_of_spanning}); and the
geometric hypothesis $\mathrm{GH}(X_Z,\Lambda_0)$ of \cref{pf:lift1:hsdata} holds for the associated datum
$X_Z$ with $\Lambda_0=\exp(\liftonedfl PP(\kappa/T+1))$, $\Lambda_D=\Lambda_T=e^{PP\kappa}$
(\lref{GlobalTheory.T.hermitian_geometric_hypothesis_holds}). \Cref{pf:lift1:hscomp} then gives
the claim with $d$ the dimension of $\liftonecusp_Z(\mathrm{wt}_T)$, which is $\liftonerkC(\mathrm{wt}_T,Z)$ since
the forms of weight $\mathrm{wt}_T$ are good.
\end{proof}

\begin{proposition}[Degree parameters]\label{pf:lift1:degpar}
Under the hypotheses of \cref{pf:lift1:hs-thm} except the cover degree and share conditions, the cover
condition defined by the following predicate (which does not involve the constant) holds: for every component $Z$ there
are reals $n,\Lambda,\beta,\Lambda_D,\Lambda_T\ge1$ and $J$ with
$\log n,\log\beta,\log\liftonedg(u_1,Z)\le PP$, $\log\Lambda\le PP(\kappa/T+1)$,
$\log\Lambda_D,\log\Lambda_T\le PP\kappa$,
$J\le4(\liftonedfl PP(\kappa/T+1)+TPP)+2PP(\kappa/T)\liftonedfl$ and
\[
\mathrm{HS}_{26}\bigl(T,\liftonedfl;\,n,1,\Lambda,\beta,\Lambda_D,\Lambda_T,J;\,\liftonedg(u_1,Z),
\liftonehgt(u_1,Z),\liftonerkC(\mathrm{wt}_T,Z),\liftonedegC(\mathrm{wt}_T,Z)\bigr),
\]
where $\kappa=|\kappa_i|$ and $u_1$ is the weight of the parameter $P$.
\formallabel{GlobalTheory.HSShortRun.hsClaim_short_alt_alt}
\formallabel{GlobalTheory.T.exists_degree_parameters_bound_alt_alt}
\end{proposition}

\begin{proof}
By \cref{pf:lift1:combine} we combine the following cover conditions: cusp frames $f_1,f_2$ for $u_1,u_2$
with $\Lambda,\beta\ge1$, $\log n\le PP$, $\log\beta\le PP$, $\log\Lambda\le PP(\kappa/T+1)$, the frame norm
bound and the Bergman bound (\lref{GlobalTheory.T.hilbertSamuel_bergman_frame_bound_alt}); the
generation of function fields for $u_1$ and $u_2$
(\lref{GlobalTheory.T.birOn_weightPair_both_components_alt}); $\log\liftonedg(u_1,Z)\le PP$
(\lref{GlobalTheory.T.log_dg_u1_le_of_flatDeg}); the largeness condition for $P$
(\lref{GlobalTheory.HSApp.wtBig_run_alt}); the properties of the weights
(\lref{GlobalTheory.T.weight_parameter_strict_antitone_alt}) and sizes
(\lref{GlobalTheory.T.size_bounds_for_family_alt}); the short spanning family for
$\mathrm{wt}_0$ (\lref{GlobalTheory.T.short_vectors_span_cusp_lattice_alt}); and the short spanning
families for $\mathrm{wt}_t$, $t\le T$, together with $PP\kappa\ge0$
(\lref{GlobalTheory.T.short_path_weight_bound_alt}). The goodness of the weights $\mathrm{wt}_t$
follows from \lref{GlobalTheory.T.covering_domain_lattice_regularity}, as their entries are
non-increasing with the required gap. \Cref{pf:lift1:hslevel} gives the claim with
$\Lambda_D=\Lambda_T=e^{PP\kappa}$ and $J=\liftonedfl PP(\kappa/T+1)$, and this $J$ satisfies the stated
bound because the remaining terms are nonnegative.
\end{proof}

\begin{lemma}\label{pf:lift1:alog}
Let $1\le D\le T$ be integers and $PP,b,z$ reals with $PP\ge1$, $b\ge1$, $\log b\le PP$,
$T\,PP\le\sqrt b$ and $0\le z\le\frac12$. Then $A=26(26D^2+2zT)$ satisfies $A\log\max(A,1)\le6000\,b$.
\formallabel{GlobalTheory.hilbertSamuel_ALogA_le}
\end{lemma}

\begin{proof}
$A\ge1$ and $A\le702T^2$. From $T\le\sqrt b$, $\log T\le\frac12\log b\le\frac12PP$, so
$\log A\le\log702+2\log T\le7+PP\le8PP$. Hence $A\log A\le702T^2\cdot8PP\le5616(TPP)^2\le5616\,b$, using
$PP\le PP^2$ and $(TPP)^2\le b$.
\end{proof}

\begin{lemma}\label{pf:lift1:degnum}
Let $T,D\in\N$ and $PP,\kappa,b,z,\rho,n,\Lambda,\beta,\Lambda_D,\Lambda_T,J,\delta,h,d,\deg,\deg_M,c\in\R$ with
$1\le D$, $3D+5\le T$, $PP\ge1$, $b\ge1$, $\log b\le PP$, $TPP\le\sqrt b$, $z=7(D^2/T+D\rho)$, $\rho\ge0$,
$z\le\frac12$, $zPP\kappa\le b$, $\kappa\ge0$, $n,\Lambda,\beta,\Lambda_D,\Lambda_T,\delta\ge1$, $d>0$,
$\log n,\log\beta,\log\delta\le PP$, $\log\Lambda\le PP(\kappa/T+1)$, $\log\Lambda_D,\log\Lambda_T\le PP\kappa$,
$J\le4(DPP(\kappa/T+1)+TPP)+2PP(\kappa/T)D$, $(1-z)d\le\binom{T-1}D\delta$, the claim
$\mathrm{HS}_{26}(T,D;n,1,\Lambda,\beta,\Lambda_D,\Lambda_T,J;\delta,h,d,\deg)$, $c\ge0$ and
$\deg-dc\le\deg_M$. Then
\[
\binom{T-1}D\delta\Bigl(Th-3(D+1)|h|-26(26D^2+2zT)\log(2+|h|)-(4000b+2c)\Bigr)\le\deg_M .
\]
\formallabel{GlobalTheory.CovHSNum.deg_bound_of_claim_numbers}
\end{lemma}

\begin{proof}
Let $B=\binom{T-1}D$ and $l_0=27PP(q+1)+65PP$ with $q=\kappa/T$. First,
$\log\Lambda+26\log(nD\delta(1+\log\Lambda))\le l_0$, from the bounds on $\log n,\log\delta,\log\Lambda$,
$\log D\le\log T\le PP/2$ and $\log(1+\log\Lambda)\le\log\Lambda$. The development rearranges the claim
(\lref{GlobalTheory.T.hilbertSamuel_degree_ratio_lower_bound_rearranged}) into
\begin{multline*}
B\delta\Bigl(Th-3(D+1)|h|-26(26D^2+2zT)\log(2+|h|)-\bigl(J+26D^2(l_0+\log T)+(T+D+1)\log\beta\\
+(T+D^2)(g+2)+2z(T(\log2+l_0+\tfrac12\log\beta)+\log\Lambda_D+\log\Lambda_T)\bigr)\Bigr)\le\deg ;
\end{multline*}
this uses $2+|h|+\log\Lambda\le(1+\log\Lambda)(2+|h|)$, so that $\ell_\rho\le l_0+26\log(2+|h|)$;
$Z_s\le z$, which follows from $(1-z)d\le B\delta$; and $d(1-Z_s)=B\delta$ and $dZ_s\le2zB\delta$ (as
$z\le\frac12$). The size hypotheses bound the bracketed error by $4000b$: e.g.\ $(T+D+1)\log\beta\le2b$,
$g+2\le5DPP$ and $(T+D^2)(g+2)\le10b$ (using $D^2\le T/14$, which follows from $z\le\frac12$, and
$T\,D\,PP\le(TPP)^2\le b$), and the term with factor $2z$ is at most $200b$. Finally $d\le2B\delta$, so
$dc\le2B\delta c$, and $\deg_M\ge\deg-dc$ gives the claim.
\end{proof}

\begin{proof}[Proof of \cref{pf:lift1:hs-thm}]
Let $C_p\ge0$ be the constant of \lref{GlobalTheory.T.twisted_bundle_degree_lower_bound}, which
states: under the hypotheses of \lref{GlobalTheory.T.short_generator_norm_comparison} (see the proof of
\cref{pf:lift1:short}), for every weight $\lambda$ and every parallel weight $w$ with $|w(\tau)|\le3r$
and every component $Z$ on which the forms of weight $\mathrm{wt}(\lambda)$ are good,
$\liftonedegC(\lambda,Z)-\liftonerkC(\lambda,Z)\,C_pz_L^2\le\liftonedegC(\lambda+w,Z)$ (twisting by a parallel
weight changes degrees by at most $C_pz_L^2$ per unit of rank). We show that the cover height sum
condition holds with constant $10000+2C_p$, combining by \cref{pf:lift1:combine}
\cref{pf:lift1:degpar}, the size bounds (\lref{GlobalTheory.T.size_bounds_for_family_alt}), the weight
properties (\lref{GlobalTheory.T.weight_parameter_strict_antitone_alt}) and the size, flat degree,
cover degree and share conditions. Fix $p,i,Y$ and an admissible $u_1$; then $u_1$ is the weight of
$P$ and $u_M=\mathrm{wt}_T+w$ with $|w(\tau)|\le3r$. The size conditions give $b\ge1$, $b\le U$,
$3\liftonedfl+5\le T$, $0\le\rho(u_1)$, $\zeta(u_1)\le\frac12$, $\zeta(u_1)PP\kappa\le b$,
$z_L^{12}\le PP$, $TPP\le\sqrt b$ and $|\liftoneflat(\liftonelogOm)|\le9rL$ on $Y$; also $\liftonedfl\ge1$,
$z_L\ge1$ (as $f\ge8$, $r\ge100$), $PP\ge1$, $\log b=L\le z_L\le PP$, and $z_L^2\le z_L^{12}\le PP\le TPP\le\sqrt
b\le b$.

Put $A=26(26\liftonedfl^2+2\zeta(u_1)T)>0$ and $R=4000b+2C_pz_L^2\ge0$. By \cref{pf:lift1:alog},
$A\log\max(A,1)+R\le6000b+4000b+2C_pz_L^2\le(10000+2C_p)U$.

Let $Z$ be a component. The cover degree condition gives $\liftonedg(u_1,Z)>0$, hence
$\liftonedg(u_1,Z)\ge1$ (it is an integer). The rank of the forms of a weight does not change under
twisting by a parallel weight (\lref{GlobalTheory.T.twisted_flat_bundle_ranks}), so the share
condition gives $B_T\liftonedg(u_1,Z)\le\liftonerkC(\mathrm{wt}_T,Z)$ and
$(1-\zeta(u_1))\liftonerkC(\mathrm{wt}_T,Z)\le B_T\liftonedg(u_1,Z)$; as $B_T\ge1$,
$\liftonerkC(\mathrm{wt}_T,Z)>0$. The forms of weight $\mathrm{wt}_T$ on $Z$ are good
(\lref{GlobalTheory.T.covering_domain_lattice_regularity}), so the twisting bound gives
$\liftonedegC(\mathrm{wt}_T,Z)-\liftonerkC(\mathrm{wt}_T,Z)C_pz_L^2\le\liftonedegC(u_M,Z)$. Now
\cref{pf:lift1:degnum}, applied with $D=\liftonedfl$, $\kappa=|\kappa_i|$, $z=\zeta(u_1)$, $\rho=\rho(u_1)$,
the numbers of \cref{pf:lift1:degpar}, $\delta=\liftonedg(u_1,Z)$, $h=\liftonehgt(u_1,Z)$,
$d=\liftonerkC(\mathrm{wt}_T,Z)$, $\deg=\liftonedegC(\mathrm{wt}_T,Z)$, $\deg_M=\liftonedegC(u_M,Z)$ and
$c=C_pz_L^2$, gives the inequality of $\liftoneHS(PP,10000+2C_p,u_1)$ at $Z$.
\end{proof}

% ---------------------------------------------------------------------------
\subsection{A cusp-regular level determinant}\label{pf:lift1:cusp}

Let $\mathcal D$ be a global datum, $T$ a core datum of $\mathcal D$ satisfying the structural laws, $j$ a
level in range, $Y$ a deep cover of $T$ at level $j$ with moduli space $X^o_Y$ over the coefficient field
$F$ and level $l_Y=(M,\mathfrak M,A)$. The group $\mathcal A_{\mathrm{lev}}$ of level automorphisms of $D_j$ and
$l_Y$ acts on framed points of $X^o_Y$, $y\mapsto\gamma\cdot y$, and hence on rules. For a finite subgroup
$S\subseteq\mathcal A_{\mathrm{lev}}$ and a rule $u$, the \emph{level determinant} is the group determinant
\[
\Delta_S(u)=\det\bigl(u_{hg}\bigr)_{h,g\in S},\qquad u_\gamma(y)=u(\gamma^{-1}\cdot y),
\]
a rule which is a section of some integral exponent when $u$ is. The \emph{section ring} is the
$F$-algebra generated by the sections of all integral exponents. A \emph{type frame} $\mathrm{fr}$ for
$(D_j,l_Y)$ determines subgroups $S_{\mathrm{HR}}(\mathrm{fr})$ and $S_{\mathrm{GHR},k}(\mathrm{fr})$
($k\in\{0,1\}$) of $\mathcal A_{\mathrm{lev}}$, built from the unitary and general unitary level groups of the
host.

\begin{theorem}[Cusp-regular level determinant]\label{pf:lift1:cusp-thm}
Assume (Hdg), (GM), (Q), (Sieg), (Tor) and (Lie). Let $T$ satisfy the structural laws,
$r_+,r_-\ge1$ at level $j$, $2<[E:\Q]/2$, $j$ in range, $Y$ a deep cover which is good with parameters
$(e,o)$, $e\ge4$, and $g_k>r+1$ for the auxiliary primes $g_k$ of the host. Let $\mathrm{fr}$ be a type frame and
$S=S_{\mathrm{HR}}(\mathrm{fr})$ or $S=S_{\mathrm{GHR},k}(\mathrm{fr})$. Then there is a cusp form
$u\in\liftonecusp(\mathrm{wt}(c_0h))$ of weight $c_0=c_0(\mathcal D,j)$ times the Hodge weight such that
$\Delta_S(u)$ is a non-zero-divisor in the section ring: $\Delta_S(u)\,f=0$ implies $f=0$.
\formallabel{GlobalTheory.T.orbits_separated_by_regular_cusp}
\formallabel{GlobalTheory.OrbitSep.exists_cusp_regular_level_determinant}
\end{theorem}

The proof uses two algebraic lemmas and a geometric separation statement.

\begin{lemma}[Generic group determinants]\label{pf:lift1:groupdet}
Let $F\subseteq\C$ be a field of characteristic $0$, $V$ an $F$-vector space, $S$ a finite group, $\iota$
a finite set, and for $i\in\iota$ let $\mathrm{ev}_i\colon V\to\C^S$ be $F$-linear, $J_i$ a finite set with a
map $r_i\colon J_i\to S$ and $W_i\subseteq V$ a subspace, such that: $1\in r_i(J_i)$;
$\mathrm{ev}_i(v)(h)=0$ for $v\in W_i$ and $h\notin r_i(J_i)$; and if $a\in\C^{J_i}$ satisfies
$\sum_ja_j\mathrm{ev}_i(v)(r_i(j))=0$ for all $v\in W_i$, then $a=0$. Then there is $v\in V$ with
$\det\bigl(\mathrm{ev}_i(v)(hg)\bigr)_{h,g\in S}\ne0$ for every $i$.
\formallabel{GlobalTheory.GenericGroupDet.exists_vec_det_ne_zero}
\end{lemma}

\begin{proof}
Fix $i$. The last hypothesis says that the $\C$-span of the vectors
$(\mathrm{ev}_i(v)(r_i(j)))_j$, $v\in W_i$, is $\C^{J_i}$; in particular $r_i$ is injective (if
$r_i(j)=r_i(j')$ with $j\ne j'$, the vector $a=e_j-e_{j'}$ would be orthogonal to all of them). Hence there
are $w_{i,1},\dots,w_{i,n_i}\in W_i$ and $c_{i,l}\in\C$ with
$\sum_lc_{i,l}\mathrm{ev}_i(w_{i,l})(r_i(j))=[r_i(j)=1]$ for all $j$; since the
$\mathrm{ev}_i(w_{i,l})$ vanish outside $r_i(J_i)$ and $1\in r_i(J_i)$, the function
$\sum_lc_{i,l}\mathrm{ev}_i(w_{i,l})$ is the indicator $\delta_1$ of $1\in S$. The matrix
$(\delta_1(hg))_{h,g}$ is a permutation matrix (its $(h,g)$ entry is $1$ exactly when $g=h^{-1}$), so
its determinant is $\pm1$. Consider the polynomials
\[
P_i(X)=\det\Bigl(\sum_{i',l}X_{i',l}\,\mathrm{ev}_i(w_{i',l})(hg)\Bigr)_{h,g}\in\C[X_{i',l}].
\]
Evaluating $P_i$ at $X_{i,l}=c_{i,l}$ and $X_{i',l}=0$ for $i'\ne i$ gives $\det(\delta_1(hg))\ne0$, so
$P_i\ne0$, and $\prod_iP_i\ne0$. A nonzero polynomial does not vanish identically on $A^N$ for an infinite
set $A\subseteq\C$; taking $A=F$ (infinite, as $\mathrm{char}F=0$) gives $t_{i',l}\in F$ with
$\prod_iP_i(t)\ne0$. Then $v=\sum t_{i',l}w_{i',l}$ satisfies
$\det(\mathrm{ev}_i(v)(hg))=P_i(t)\ne0$ for all $i$, by linearity of $\mathrm{ev}_i$.
\end{proof}

\begin{lemma}[Non-zero-divisors in the section ring]\label{pf:lift1:nzd}
Let $X^o$ be a unitary moduli space over $F$ and $\Phi$ a section of an integral exponent $e$ such that
every component $Z$ contains a complex framed point $y$ with $\Phi(y)\ne0$. If $f$ is in the section
ring and $\Phi f=0$, then $f=0$.
\formallabel{GlobalTheory.SecRingNonZeroDiv.eq_zero_of_mul_eq_zero}
\end{lemma}

\begin{proof}
For a component $Z$ let $\Omega_Z$ be its function field and $\mathrm{gen}_Z$ the $F$-algebra
homomorphism from rules to $\Omega_Z$ given by evaluation at the generic frame of $Z$. A section of an
integral exponent vanishes at all complex framed points of $Z$ if and only if its image under
$\mathrm{gen}_Z$ vanishes (\lref{GlobalTheory.CompHt.genVal_eq_zero_iff}); hence
$\mathrm{gen}_Z(\Phi)\ne0$ for every $Z$. Write $f=\sum_df_d$ (finite sum) with $f_d$ a section of exponent
$d$. For $t\in(\Omega_Z^\times)^{\mathrm{Ix}}$ there is a framed point $y_t$ over $\Omega_Z$, obtained from
the generic frame by a Borel change of frame with diagonal $t$
(\lref{GlobalTheory.SecSeparation.exists_frBor}), and the value at $y_t$ of a section $g$ of exponent
$d$ is $\chi_d(t)\,\mathrm{gen}_Z(g)$ with $\chi_d(t)=\prod_\iota t_\iota^{d_\iota}$. Evaluating $\Phi f=0$ at
$y_t$ and dividing by the unit $\chi_e(t)\mathrm{gen}_Z(\Phi)$ gives
$\sum_d\mathrm{gen}_Z(f_d)\chi_d(t)=0$ for all $t$. The characters $\chi_d$ of the group
$(\Omega_Z^\times)^{\mathrm{Ix}}$ are pairwise distinct (evaluate at $t=2$ in one coordinate and $1$
elsewhere; $\Omega_Z$ has characteristic $0$), hence linearly independent over $\Omega_Z$ by Artin's
theorem on characters. So $\mathrm{gen}_Z(f_d)=0$ for all $d$ and $Z$, and a section of an integral
exponent with this property is $0$ (\lref{GlobalTheory.CompHt.secZ_eq_zero_of_genVal}). Hence $f=0$.
\end{proof}

\subsubsection*{Three geometric inputs}
Let $\mathrm{TF}$ be a toroidal flag datum for $X^o$ (\cref{pf:lift1:hodge}), with model $\overline X$, interior
$j\colon X^o_F\to\overline X$ and flag model $\pi$, and let $Z$ be a component. Let $X_Z$ be the closure of
$Z$ in $\overline X$, $X_Z^o=j^{-1}(X_Z)$ its interior part and $j_Z\colon X_Z^o\to X_Z$. A \emph{generic
eigenframe} is a frame of the Lie algebra of the semi-abelian scheme at the generic point of $X_Z$ adapted
to the eigenspaces of $\OO_E$; for a local eigenframe at a point $x$ let $d_\tau\in K(X_Z)^\times$ be the
determinant of the matrix expressing its $\tau$-part in the generic one. The development proves that
$X_Z$ is integral (\lref{GlobalTheory.BaseVar.isVariety_barComp}) and that there are a generic
eigenframe and, for every $a\colon\Hom(E,\C)\to\Z$, a Cartier divisor $D_a$ on $X_Z$ with local equations
$\prod_\tau d_\tau^{-a(\tau)}$ whose pull-back to the flag model has the class of the automorphic line
bundle with exponent $a(\tau)$ on the $\tau$-entries (\lref{GlobalTheory.T.exists_cartierDivisor_det_lieEigenparts_toroidal}).
For the Hodge weight $h$ the exponents are $0$ for $\tau\in\Psi$ and $1$ for $\tau\notin\Psi$; for the
weight $\kappa_{\mathrm{base}}$ equal to $r_-$ on the first $r_+$ entries and to $-r_+$ on the others they
are $r_-$ for $\tau\in\Psi$ and $r_+$ for $\tau\notin\Psi$.

An action $\rho'$ of a group $\Gamma$ on an integral scheme $X'$ and a line bundle $L$ on $X'$ admit a
\emph{rational linearisation} if there are a Cartier divisor $D'$ with $\OO(D')$ in the class of $L$ and a
map $c\colon\Gamma\to K(X')^\times$ with $c(\gamma\gamma')=\gamma'^*(c(\gamma))\,c(\gamma')$, such that for
every local equation $\eta$ of $D'$ at $\rho'(\gamma)(x)$ the function $\gamma^*(\eta)\,c(\gamma)^{-1}$ is a local
equation of $D'$ at $x$, and, if $\gamma$ fixes $x$, $\gamma^*(\eta)\,c(\gamma)^{-1}\eta^{-1}-1$ vanishes at $x$ (so
that $\gamma$ acts trivially on the fibre at its fixed points).

\begin{theorem}[Linearisation of the Hodge line bundle]\label{pf:lift1:hodgelin}
Assume (Hdg). Let $r\ge1$, let the $A$-part of the level be $1$ and let $F$ be algebraically closed in the
function fields of the components. Then there is a line bundle $\delta$ on $X_Z$ whose pull-back to the
flag model has the class of the automorphic line bundle of $h$, with the following property. Let $I$ be an
ideal of $\OO_E$ and $g$ a prime with $\Nm I=g$, $g\notin I^2$, $g>r+1$ and $\mathfrak M\subseteq I$; let
$\Gamma$ be a finite group with homomorphisms $\iota$ to the group of level automorphisms, such that
every $\iota(\gamma)$ is unipotent modulo $I$, and $\rho$ to $\Aut(X_Z^o)$, such that $\rho$ is compatible
with the structure map and acts on framed points by the change of level $\iota(\gamma)^{-1}$; and let $X'$
be a normal integral scheme with a morphism $\varphi\colon X'\to X_Z$ projective over $F$, an open immersion
$\ell\colon X_Z^o\to X'$ with $\varphi\ell=j_Z$, and an action $\rho'$ of $\Gamma$ on $X'$ with
$\rho'(\gamma)\ell=\ell\rho(\gamma)$. Then $\rho'$ and $\varphi^*\delta$ admit a rational linearisation.
\formallabel{GlobalTheory.T.hodgeLineBundle_rational_linear_equiv_class}
\formallabel{GlobalTheory.HodgeLin.hodgeLineBundle_rationally_linear}
\end{theorem}

\begin{proof}
Let $\delta=\OO(D_a)$ for the exponents $a$ of $h$; its pull-back to the flag model has the class of the line
bundle of $h$. Now let $I,g,\Gamma,\iota,\rho,X',\varphi,\ell,\rho'$ be as in the statement. The development
constructs an equivariant structure on the pull-back of the semi-abelian scheme to $X'$: a group scheme
$G'$ over $X'$ with an $\OO_E$-action and a homomorphism to the semi-abelian scheme, and morphisms
$m(\gamma)$ of $G'$ over $\rho'(\gamma)$ forming an equivariant linearisation
(\lref{GlobalTheory.T.equivariant_linearization_pullback_compat}; the unipotence of the $\iota(\gamma)$
modulo the neat prime $I$, the normality of $X'$ and the extension of homomorphisms of semi-abelian schemes
from (Hdg) are used here). The interior part of $X_Z$ is a nonempty open subset of
the integral scheme $X_Z$ (it contains the points of $Z$), so it contains the generic point, and $\varphi$, whose image contains
$\varphi(\ell(X_Z^o))=j_Z(X_Z^o)$, is dominant. Let $s_Z$ be the structure morphism of $X_Z$. The morphisms
$s_Z\varphi\rho'(\gamma)$ and $s_Z\varphi$ from $X'$ to $\mathrm{Spec}\,F$ agree on the dense open subset $\ell(X_Z^o)$,
where $\rho'=\rho$ is compatible with the structure map; as $X'$ is integral, they agree. Finally, the transformation rule of the determinants
$d_\tau$ under the equivariant structure gives the rational linearisation of $\varphi^*\OO(D_a)$
(\lref{GlobalTheory.T.lieDeterminant_linearisation_equivariant_pullback}, which proves this for
every dominant $\varphi$, action $\rho'$ compatible with the structure map and equivariant structure).
\end{proof}

\begin{theorem}[Log canonical bundle]\label{pf:lift1:kappa}
Let $r\ge1$, let $F$ be algebraically closed in the function fields of the components, and suppose that
$X_Z$ is smooth over $F$ of relative dimension $\dim D_w$, with canonical bundle $K$. Then there are a
line bundle $\kappa$ on $X_Z$ whose pull-back to the flag model has the class of the automorphic line
bundle of $\kappa_{\mathrm{base}}$, and a section $\theta$ of $\kappa\otimes K^{-1}$ whose nonvanishing locus
is exactly the interior $j_Z(X_Z^o)$.
\formallabel{GlobalTheory.T.exists_toroidalCanonicalBundle_weight_line}
\formallabel{GlobalTheory.LogCanon.exists_kap_theta}
\end{theorem}

\begin{proof}
Let $\kappa=\OO(D_a)$ for the exponents $a$ of $\kappa_{\mathrm{base}}$; its pull-back has the required
class. Let $D'$ be a Cartier divisor presenting $K$ (from a trivialisation of $K$ at the generic point). The
logarithmic Kodaira--Spencer isomorphism of the toroidal model gives a rational function $h_0$ on $X_Z$
such that at every point $x$, for local equations $g$ of $D_a$ and $g'$ of $D'$, the function $h_0g/g'$ is
regular at $x$, and a unit at $x$ if and only if $x$ lies in the interior
(\lref{GlobalTheory.T.exists_logKodairaSpencer_divisor_toroidalComponent}; this expresses that
$\kappa$ is isomorphic to $K$ twisted by an effective divisor whose support is the boundary). The functions $h_0g/g'$ are the local expressions of a
global section $\theta$ of $\OO(D_a)\otimes K^{-1}$ in the trivialisations given by $g$ and $g'$; it is regular
everywhere and nonvanishing exactly on the interior.
\end{proof}

\begin{theorem}[Separation in a fibre of a finite quotient]\label{pf:lift1:quotsep}
Assume Nadel vanishing for non-klt divisors, the numerics of nef divisors and principalisation. Let $F$
be a number field and $\sigma\colon F\to\C$. Let $X$ be a smooth projective geometrically connected
integral $F$-variety of dimension $n\ge1$ with canonical bundle $K$; $\varphi\colon X'\to X$ a morphism from an
integral scheme $X'$ projective over $F$; $q\colon X'\to W$ a finite surjective morphism to an integral
projective $F$-scheme, compatible with the structure maps; and $V$ a scheme with open immersions
$i\colon V\to X$, $\ell\colon V\to X'$, $\varphi\ell=i$. Let $\delta,N',B$ be line bundles on $X$ and $N$ one on
$W$ with $q^*[N]=\varphi^*[\delta]$, $\delta$ nef and big, $N'$ nef, $m>\binom{n+1}2$ and
$[B]=m[\delta]+[N']$. Let $(v_j)_{j\in J}$ be distinct complex points of $V$ over $\sigma$, all in one fibre of
$q\ell$, one of which lies over the generic point of $X$. If $a\in\C^J$ satisfies
$\sum_ja_j\,g(v_j)=0$ for every global section $g$ of $K\otimes B$ (values in local trivialisations at the
points $v_j$), then $a=0$.
\formallabel{GlobalTheory.T.sections_separate_quotient_generic_fibre_descends}
\formallabel{GlobalTheory.QuotSep.sections_separate_quotient_fibre}
\end{theorem}

\begin{proof}
\emph{Resolution.} Using principalisation, the development constructs a dense open $U\subseteq X$, a
modification $\mu\colon M\to X$ with $M$ smooth projective and $\mu$ an isomorphism over $U$, a smooth
projective geometrically connected integral $W'$ of dimension $n$ with a projective surjective morphism
$f\colon M\to W'$, a line bundle $N_W$ on $W'$ with $f^*[N_W]=\mu^*[\delta]$, and complex points $z_j$ of $M$
over $U$ with $\mu(z_j)=i(v_j)$, all lying over the generic point of $W'$
(\lref{GlobalTheory.QuotSetup.exists_resolved_modification}). The $z_j$ are distinct since the $v_j$ are, and $M$ is
geometrically connected as $X$ is and $\mu$ is birational.

\emph{Positivity.} By the numerics of nef divisors, $\mu^*\delta$ and $\mu^*N'$ are nef; $\mu^*\delta$ is big
since $\mu$ is an isomorphism over a dense open set; $N_W$ is nef because its pull-back $f^*N_W$, whose
class is $\mu^*[\delta]$, is nef and $f$ is surjective and projective
(\lref{GlobalTheory.T.isNefLineBundle_of_pullback_surjective_projective}); and $f^*N_W$ is big. Hence
$N_W$ is ample at the generic point of $W'$, and $f$ is \'etale over an open neighbourhood $U_W$ of the
generic point (\lref{GlobalTheory.QuotSep.generic_data}).

\emph{Sections.} Let $K_M$ be the canonical bundle of $M$. There is a map
$\Psi\colon H^0(M,K_M\otimes\mu^*B)\to H^0(X,K\otimes B)$ such that at every complex point $z$ of $M$ over $U$
the values satisfy $\Psi(g)(\mu(z))=\lambda_z\,g(z)$ with $\lambda_z\in\C^\times$ depending only on $z$ and
the trivialisations (\lref{GlobalTheory.T.exists_push_sections_canonical_modification}: canonical
forms push forward along a modification). Moreover
$[\mu^*B]=m\,f^*[N_W]+[\mu^*N']$.

\emph{Separation on $M$.} The development proves, using Nadel vanishing and the numerics of nef
divisors: for $M$ and $W'$ smooth projective geometrically connected of dimension $n\ge1$, $f$ surjective
and \'etale over $U_W$, $N_W$ nef, $N'_M$ nef, $m>\binom{n+1}2$ and $[B_M]=m\,f^*[N_W]+[N'_M]$, the global
sections of $K_M\otimes B_M$ separate any finite set of distinct complex points of $M$ over $\sigma$ lying in
one fibre of $f$ over a point of $U_W$ at which $N_W$ is ample
(\lref{GlobalTheory.QuotSep.separate_numberField'}). Apply this to the $z_j$ and
$b_j=a_j\lambda_{z_j}$: for every $g$,
$\sum_jb_jg(z_j)=\sum_ja_j\Psi(g)(v_j)=0$ by hypothesis. Hence $a_j\lambda_{z_j}=0$, so $a=0$.
\end{proof}

\begin{proposition}[Separation of orbits by cusp forms]\label{pf:lift1:sep}
Assume the inputs of \cref{pf:lift1:cusp-thm}. Let $Y$ be good, $2<[E:\Q]/2$, $r_+,r_-\ge1$, and
$S\subseteq\mathcal A_{\mathrm{lev}}$ a subgroup, and suppose there are an ideal $I$ of $\OO_E$ and a prime
$g>r+1$ with $\Nm I=g$, $g\notin I^2$ and $\mathfrak M\subseteq I$, such that every $\gamma\in S$ is
unipotent modulo $I$ ($(\gamma-1)^m\Lambda\subseteq I\Lambda$ for some $m$). Let
$s\ge r+\binom{\dim D_j+1}2+1$, where $\dim D_j=\frac{[E:\Q]}2r_+r_-$. Then every component $Z$ contains a
complex framed point $y$ such that the following holds: if $(a_h)$ are complex numbers indexed by the
$h\in S$ with $h^{-1}\cdot y\in Z$, and $\sum_ha_h\,u(h^{-1}\cdot y)=0$ for every
$u\in\liftonecusp_Z(\mathrm{wt}(s\,h))$, then $a=0$.
\formallabel{GlobalTheory.OrbitSep.separate_orbits_by_cusp}
\end{proposition}

\begin{proof}
The steps are the following; each is a theorem of the development, proved above or cited.
\begin{enumerate}[label=(\arabic*)]
\item There is a toroidal flag datum $\mathrm{TF}$ for $X^o_Y$ (a smooth projective toroidal model with
  normal crossings boundary and a flag model over it,
  \lref{GlobalTheory.T.torFlag_nonempty_of_one_le_dimFlag}); since $F$ is algebraically closed in
  the function field of $Z$, the closure $X_Z$ of $Z$ in it is integral, smooth, projective and
  geometrically connected (\lref{GlobalTheory.T.flagSpace_component_smooth_projective_connected},
  \lref{GlobalTheory.T.compactificationComponent_integral_smooth_projective}).
\item Choose a complex framed point $y$ over the generic point of $Z$
  (\lref{GlobalTheory.T.exists_complexFramedPoint_over_componentGenericPoint}).
\item The stabiliser $\Gamma\subseteq S$ of $Z$ acts on the open part of $X_Z$ by change of level
  (\lref{GlobalTheory.T.componentStabilizer_acts_by_levelChange}), and there are a normal
  projective $\Gamma$-equivariant compactification $X'$ with a morphism $X'\to X_Z$ and a finite surjective
  quotient map $q\colon X'\to W$ to a projective $W$
  (\lref{GlobalTheory.T.equivariant_normal_projective_compactification}).
\item The Hodge line bundle $\delta$ on $X_Z$ of \cref{pf:lift1:hodgelin} (whose pull-back to the flag
  space has the class of the line bundle of the Hodge weight) admits, after pull-back to $X'$, a rational
  linearisation for the action of $\Gamma$; here the unipotence modulo $I$ is used. By descent along the
  finite quotient there is a line bundle $N$ on $W$ with $q^*N\cong\varphi^*\delta$, $\varphi\colon X'\to X_Z$
  (\lref{GlobalTheory.T.exists_descent_linearised_lineBundle_finiteGroup_quotient}).
\item Let $K$ be the canonical bundle of $X_Z$. By \cref{pf:lift1:kappa} there are a line bundle $\kappa$ on
  $X_Z$, whose pull-back to the flag space has the class of the line bundle of $\kappa_{\mathrm{base}}$, and
  a section $\theta$ of $\kappa\otimes K^{-1}$ whose nonvanishing locus is the interior. Write $s=m_0+t$ with
  $m_0=\binom{\dim D_j+1}2+1$, and let $N'$ be a line bundle of class $t\delta-\kappa$ and $B$ one of class
  $m_0\delta+N'$; then $(\kappa\otimes K^{-1})\otimes(K\otimes B)$ has class $s\delta$. The bundles
  $\delta$ and $N'$ are nef, by descent of nefness along the surjective projective map from the flag
  space and the nefness of the line bundles of weights with nonnegative gaps
  (\lref{GlobalTheory.T.isNefLineBundle_of_pullback_surjective_projective},
  \lref{GlobalTheory.T.base_weight_line_bundle_is_nef}); $\delta$ is big
  (\lref{GlobalTheory.T.hodgeLineBundle_bigness_on_shimuraVariety}, whose proof uses
  \cref{pf:lift1:hodge-thm}).
\item Global sections of a line bundle of class $s\delta$ whose nonvanishing locus lies in the interior
  give cusp forms of weight $s\,h$ on $Z$, whose values at interior points are those of the sections
  up to nonzero constants
  (\lref{GlobalTheory.T.exists_cuspRule_of_boundaryVanishing_componentSection}). For every section
  $\sigma$ of $K\otimes B$, $\theta\sigma$ is such a section.
\item Effective separation (\cref{pf:lift1:quotsep}): as $X_Z$ is smooth projective geometrically
  connected of dimension $n=\dim D_j\ge1$, $q$ is finite surjective, $\delta$ is nef and big with
  $q^*N=\varphi^*\delta$, $N'$ is nef and $B=m_0\delta+N'$ with $m_0>\binom{n+1}2$, the sections of $K\otimes B$
  separate any finite set of distinct complex points of the interior lying in one fibre of $q$, one of
  which lies over the generic point.
\item The points $h^{-1}\cdot y$, $h\in S$ with $h^{-1}\cdot y\in Z$, are distinct: if two coincide, the
  corresponding element of $S$ is unipotent modulo $I$ and fixes a complex point, hence is $1$
  (\lref{GlobalTheory.T.levelAutomorphism_unipotent_fixes_complexPoint_eq_one}); they lie in one
  fibre of $q$ since they differ by elements of $\Gamma$; and $y$ lies over the generic point.
\end{enumerate}
Now let $a$ be as in the statement. For every global section $\sigma$ of $K\otimes B$, $\theta\sigma$ gives
a cusp form $u$ of weight $s\,h$ by (6), and the hypothesis for $u$ becomes
$\sum_ha_h\lambda_h\theta(h^{-1}\cdot y)\,\sigma(h^{-1}\cdot y)=0$ with constants $\lambda_h\ne0$. By (7) and
(8), $a_h\lambda_h\theta(h^{-1}\cdot y)=0$ for all $h$, and $\theta(h^{-1}\cdot y)\ne0$ as these points lie
in the interior. Hence $a=0$.
\end{proof}

\begin{proof}[Proof of \cref{pf:lift1:cusp-thm}]
The $A$-part of the level of $Y$ is $1$, $r_++r_-=r$, and $M=g_0g_1\cdot15^e$. Since $g_k>r+1$, the
development finds an ideal $I$ and a prime $g$ as in \cref{pf:lift1:sep} with every element of $S$
unipotent modulo $I$ (\lref{GlobalTheory.OrbitSep.exists_neatPrime_unipAt}). By definition
$c_0\ge r+\binom{f r_+r_-+1}2+2$, and $\dim D_j=fr_+r_-$, so \cref{pf:lift1:sep} applies with
$s=c_0$ and gives points $y(Z)$, one for each of the finitely many components. Apply
\cref{pf:lift1:groupdet} with $V=\liftonecusp(G)$, $G=\mathrm{wt}(c_0h)$, $\iota=\mathrm{Comp}$,
$\mathrm{ev}_Z(u)(h)=u(h^{-1}\cdot y(Z))$, $J_Z=\{h\in S:h^{-1}\cdot y(Z)\in Z\}$ with the inclusion into
$S$, and $W_Z$ the forms supported on $Z$: $1\in J_Z$; a form supported on $Z$ vanishes at points off
$Z$; and the independence hypothesis is \cref{pf:lift1:sep}. We obtain $u\in\liftonecusp(G)$ with
$\det(u((hg)^{-1}\cdot y(Z)))_{h,g}\ne0$ for every $Z$. Evaluation at $y(Z)$ is a ring homomorphism on
rules and the entries of the group matrix evaluate to $u((hg)^{-1}\cdot y(Z))$
(\lref{GlobalTheory.LevelKeeps.evR_transl}), so $\Delta_S(u)(y(Z))\ne0$ for every $Z$. Since
$\Delta_S(u)$ is a section of an integral exponent
(\lref{GlobalTheory.LevDesc.levDet_mem_secZ}), \cref{pf:lift1:nzd} concludes. The hypothesis $e\ge4$ is not used.
\end{proof}

% ---------------------------------------------------------------------------
\subsection{The Hodge determinant dominates}\label{pf:lift1:hodge}

Let $X^o$ be a unitary moduli space over a coefficient field $F$ (for a datum $D_w$ and a level $l$), and
$X^o_F$ its scheme over $F$. A \emph{model} of $X^o$ is a normal, flat, projective $F$-scheme $\overline X$
with an open immersion $j\colon X^o_F\to\overline X$ (the image is the \emph{interior}) and a
semi-abelian scheme over $\overline X$ extending the universal abelian scheme; $\det\Omega$ denotes the
dual of the determinant of the Lie algebra of this semi-abelian scheme (the Hodge determinant). A
\emph{toroidal flag datum} $\mathrm{TF}$ for $X^o$ is a model $\overline X$ which is smooth of relative
dimension $\dim D_w$ over $F$, has a boundary with simple normal crossings and a logarithmic
Kodaira--Spencer isomorphism, together with a flag model $\pi\colon\mathrm{Fl}\to\overline X$; on
$\mathrm{Fl}$ every weight $\lambda$ defines an automorphic line bundle $\mathcal L_\lambda$. Let $\mu_0$
be the weight equal to $1$ on the first $r_+$ entries and to $-1$ on the others. Line bundles are
compared through their classes in the (\v{C}ech) Picard group.

We use the following inputs of \cref{tab:covolume:inputs} and of the Hodge package: (Ext) extension of
homomorphisms of semi-abelian schemes; (Sieg) the existence of toroidal compactifications of Siegel
moduli spaces, here only through the polyhedral cone decompositions it provides; (Ch) toroidal charts
for fans: for an admissible invariant fan without self-intersection there is a smooth toroidal model with
a boundary divisor $\mathrm{BD}$ whose canonical section does not vanish on the interior, and $k\ge1$
such that $(\det\Omega)^k\otimes\OO(-\mathrm{BD})$ is ample.

\begin{theorem}[The Hodge determinant dominates]\label{pf:lift1:hodge-thm}
Assume (Ext), (Sieg) and (Ch), and let $\mathrm{TF}$ be a toroidal flag datum for $X^o$, with model
$\overline X$. There is a line bundle $\omega$ on $\overline X$ such that $\pi^*\omega$ has the class of
$\mathcal L_{\mu_0}$, and such that for every point $w$ of $X^o_F$ and every line bundle $A$ on
$\overline X$ there are $m\ge1$ and a global section $s$ of $\omega^{\otimes m}\otimes A^{-1}$ with
$s(j(w))\ne0$.
\formallabel{GlobalTheory.HodgeDominates.exists_lb_dominates}
\end{theorem}

\begin{lemma}\label{pf:lift1:common}
Let $N$ be a quasi-compact quasi-separated scheme, $p\colon N\to X$ and $p'\colon N\to X'$ morphisms such
that $p^\sharp\colon\Gamma(V,\OO_X)\to\Gamma(p^{-1}V,\OO_N)$ is bijective for every open $V\subseteq X$,
and $n_0\in N$ such that for every open $W\ni p(n_0)$ there is an open $W'\ni p'(n_0)$ with
$p'^{-1}(W')\subseteq p^{-1}(W)$. Let $\omega$ be a line bundle on $X$, and $\omega',P,B$ line bundles on
$X'$ with $P$ ample, $p^*\omega\cong p'^*\omega'$, and $[P]+[B]=k[\omega']$ for some $k\ge1$. If $B$ has a
section $b$ with $b(p'(n_0))\ne0$, then for every line bundle $A$ on $X$ there are $m\ge1$ and a section
$s$ of $\omega^{\otimes m}\otimes A^{-1}$ with $s(p(n_0))\ne0$.
\formallabel{GlobalTheory.HodgeDominates.exists_sec_of_common}
\end{lemma}

\begin{proof}
Let $W\ni p(n_0)$ be an open set on which $A$ is trivial, and $W'$ as in the hypothesis. Since $P$ is ample,
there are $n\ge1$ and a section $t$ of $P^{\otimes n}$ such that $X'_t=\{t\ne0\}$ is an affine open
neighbourhood of $p'(n_0)$ contained in $W'$ \cite[4.5.2]{Grothendieck1961ega2}. Then $t'=t\otimes b^{\otimes n}$ is a
section of $T^{\otimes n}$, $T=P\otimes B$, with $\{t'\ne0\}=X'_t\cap X'_b$. Its pull-back $p'^*t'$ is a
section of $p'^*T^{\otimes n}$ whose nonvanishing locus is contained in $p'^{-1}(W')\subseteq p^{-1}(W)$, on
which $p^*A$ is trivial. As $N$ is quasi-compact and quasi-separated, a section of a line bundle over
the nonvanishing locus of $p'^*t'$ extends to $N$ after multiplication by a power of $p'^*t'$
\cite[II, Lemma 5.14]{Hartshorne1977}; applied to a trivialisation of $p^*A^{-1}$ this gives $r\ge1$ such
that every line bundle $\mathcal D$ on $N$ with $[\mathcal D]=r[p'^*T^{\otimes n}]-[p^*A]$ has a section
$s'$ with $s'(n)\ne0$ wherever $p'^*t'(n)\ne0$. Take $\mathcal D=p^*(\omega^{\otimes rnk}\otimes A^{-1})$: its
class is $rnk[p^*\omega]-[p^*A]=rnk[p'^*\omega']-[p^*A]=r[p'^*T^{\otimes n}]-[p^*A]$. Since $p^\sharp$ is
bijective on sections over all opens, pulling back sections of line bundles along $p$ is bijective, so
$s'=p^*s$ for a section $s$ of $\omega^{\otimes rnk}\otimes A^{-1}$. As $t(p'(n_0))\ne0$ and
$b(p'(n_0))\ne0$, $s'(n_0)\ne0$, i.e.\ $s(p(n_0))\ne0$. Take $m=rnk$.
\end{proof}

\begin{lemma}[Common modification]\label{pf:lift1:commonmod}
Assume (Ext). For two models $\overline X,\overline X'$ of $X^o$ over $F$ whose Hodge determinants are line
bundles, there are a quasi-compact quasi-separated scheme $N$, a morphism $u\colon X^o_F\to N$ and
morphisms $p\colon N\to\overline X$, $p'\colon N\to\overline X'$ with $p\circ u=j$, $p'\circ u=j'$,
$p^\sharp$ bijective on sections over every open, $p(n)=j(w)$ whenever $p'(n)=j'(w)$, and
$p^*\det\Omega\cong p'^*\det\Omega'$.
\formallabel{GlobalTheory.HodgeDominates.exists_common_modification}
\end{lemma}

\begin{proof}
Let $N$ be the normalisation of $\overline X\times_F\overline X'$ in $X^o_F$ with respect to $(j,j')$
(the relative normalisation); $u\colon X^o_F\to N$ is an open immersion, and $p,p'$ are the compositions
of $N\to\overline X\times_F\overline X'$ (a finite morphism) with the two projections. Then $p\circ u=j$
and $p'\circ u=j'$, and the compatibility of points over the interior follows. $N$ is normal and
noetherian, $p$ is proper and birational onto the normal scheme $\overline X$, and $p^\sharp$ is
bijective on sections over every open (\lref{GlobalTheory.HodgeDominates.bijective_app}, a form of
Zariski's main theorem: $p_*\OO_N=\OO_{\overline X}$). The pull-backs to $N$ of the two semi-abelian
schemes are identified over the dense open $u(X^o_F)$, both being the universal abelian scheme there; by
(Ext) this identification extends to an isomorphism over the normal scheme $N$. Taking determinants of
Lie algebras and duals gives $p^*\det\Omega\cong p'^*\det\Omega'$.
\end{proof}

\begin{proposition}\label{pf:lift1:secdet}
Assume (Ext), (Sieg) and (Ch), $r=r_++r_-\ge1$, and let $\overline X$ be a model of $X^o$ whose Hodge
determinant $\det\Omega$ is a line bundle. For every point $w$ of $X^o_F$ and every line bundle $A$ on
$\overline X$ there are $m\ge1$ and a section of $(\det\Omega)^{\otimes m}\otimes A^{-1}$ not vanishing at
$j(w)$.
\formallabel{GlobalTheory.HodgeDominates.exists_sec_detOm}
\end{proposition}

\begin{proof}
(Sieg) provides a polyhedral cone decomposition for the Siegel space of genus
$g=\max(1,[E:\Q]\min(r_+,r_-))$. From it, the development constructs a fan $\varphi$ for the cones of
hermitian forms over $\OO_E$ of rank at most $\min(r_+,r_-)$ which is polyhedral, admissible (invariant
under the twists by the host, positive and homogeneous) and without self-intersection
(\cref{pf:lift1:fan}). By (Ch) there are a smooth toroidal model $\overline X'$ of $X^o$, a boundary
divisor $\mathrm{BD}$ whose canonical section $b$ is nonzero at the points of the interior, and $k\ge1$
such that $P=(\det\Omega')^{\otimes k}\otimes\OO(-\mathrm{BD})$ is ample; as $\overline X'$ is smooth,
$\det\Omega'$ is a line bundle (\lref{GlobalTheory.FanTheta.isLB_detOm_of_smooth}, which states that
the Hodge determinant of a model smooth over $F$ is a line bundle when $r\ge1$). Then
$[P]+[\OO(\mathrm{BD})]=k[\det\Omega']$. Let $N,u,p,p'$ be given by \cref{pf:lift1:commonmod} and
$n_0=u(w)$, so $p(n_0)=j(w)$ and $p'(n_0)=j'(w)$. For an open $W\ni j(w)$ of $\overline X$, the set
$W'=j'(j^{-1}(W))$ is an open neighbourhood of $j'(w)$ with $p'^{-1}(W')\subseteq p^{-1}(W)$, by the
compatibility of points over the interior. \Cref{pf:lift1:common} with $\omega=\det\Omega$,
$\omega'=\det\Omega'$, $B=\OO(\mathrm{BD})$ and $b$ the canonical section gives the claim.
\end{proof}

\subsubsection*{Cone decompositions}
Let $E$ be a CM field with a host. For a finitely generated projective $\OO_E$-module $X$, a
\emph{balanced form} is a symmetric $\Z$-bilinear form $\beta$ on $X$ with $\beta(ay,z)=\beta(y,\bar az)$;
they form a lattice, and the positive semidefinite real balanced forms form a cone. Cone decompositions
are formalised by \emph{invariant functions} $\varphi$: functions of forms on the modules $\OO_E^n$,
invariant under pull-back by surjections; $\varphi$ evaluates on real and integral balanced forms. A family
$\gamma=(\gamma_i)$ of integral positive semidefinite balanced forms is a \emph{cell} of $\varphi$ if
$\varphi(\sum_ic_i\gamma_i)=\sum_ic_i\varphi(\gamma_i)$ for $c_i\ge0$ and the cell is strict (there are $N>0$ and
an integral linear form $\ell$ with $N\varphi\le\ell$ on the integral positive semidefinite balanced forms,
with equality exactly on the rational nonnegative combinations of the $\gamma_i$); a \emph{cone} is a cell
whose members are part of a $\Z$-basis of the balanced forms. For $q\in\N$, $\varphi$ is a \emph{polyhedral
fan} (up to rank $q$) if on every $X$ of rank at most $q$ it is positively homogeneous, superadditive,
integral and positive on the nonzero integral positive semidefinite balanced forms, and finitely many
cells cover the cone up to automorphisms of $X$; it is a \emph{smooth fan} if moreover the faces of cones
are cones, every positive semidefinite form is a nonnegative combination of a cone, and there are
finitely many cones up to automorphisms. It has \emph{no self-intersection} if no two members of a cone are
equivalent under an automorphism of $X$, and it is \emph{twist invariant} if it is invariant under the ring
automorphisms of $\OO_E$ fixing the element $w_0$ of the host. $\varphi'$ is \emph{subordinate} to $\varphi$ if
every invariant function which is linear on the cells of $\varphi$ is linear on the cells of $\varphi'$.
\emph{Domination} means: for integral linear forms $\ell,m$ on the balanced forms of $X$ with $\ell>0$ on the
nonzero integral positive semidefinite balanced forms, $m\le C\ell$ on them for some $C\in\N$.

\begin{theorem}[Invariant smooth refinements]\label{pf:lift1:smoothfan}
Let $q\in\N$ and let $\psi$ be a twist invariant polyhedral fan (up to rank $q$), and assume domination
for every $X$. Then there is a twist invariant smooth fan $\varphi'$ (up to rank $q$) without
self-intersection such that $\psi$ is linear on every cone of $\varphi'$ on every $X$ of rank at most $q$.
\formallabel{GlobalTheory.T.exists_invariant_smoothProjective_refinement_polyFan}
\formallabel{GlobalTheory.PolFan.exists_smooth_refinement}
\end{theorem}

\begin{proof}
First, there are twist invariant invariant functions $\varphi$ and $\mathrm{lab}$ such that $\varphi$ is a
polyhedral fan subordinate to $\psi$ and $(\varphi,\mathrm{lab})$ is \emph{simplicial and labelled}:
$\mathrm{lab}$ is invariant under positive integral dilation, and every cell of $\varphi$ is contained in the
cone spanned by a linearly independent cell of $\varphi$ whose members have pairwise distinct labels
(\lref{GlobalTheory.T.exists_invariant_simplicial_labelled_fan_refinement}; obtained by an invariant
barycentric subdivision, using domination). Second, from such a pair the development constructs a twist invariant
smooth fan $\varphi'$ without self-intersection and subordinate to $\varphi$
(\lref{GlobalTheory.T.exists_invariant_smooth_refinement_polyFan}; an invariant resolution, in which star
subdivisions lower the maximal multiplicity of the cells). Subordination is transitive, so $\varphi'$ is
subordinate to $\psi$. Since $\psi$ is linear on its own cells, it is then linear on the cells, in
particular on the cones, of $\varphi'$.
\end{proof}

\begin{proposition}[Fine invariant fans]\label{pf:lift1:fan}
Let $E$ be a CM field with a host, $g,q\in\N$ with $[E:\Q]q\le g$, and $\varphi_s$ a polyhedral cone
decomposition for positive quadratic forms in $g$ variables (as provided by (Sieg); cone decompositions
are formalised by invariant functions on forms). Then there is a cone decomposition $\varphi$ for the
positive semidefinite hermitian forms over $\OO_E$ of rank at most $q$ which is polyhedral, has no
self-intersection, is invariant under the twists of the host, positive and homogeneous, and such that the
restriction of $\varphi_s$ to hermitian forms is linear on every cone of $\varphi$; given a finite set of
integers $N\ge1$, $\varphi$ can moreover be chosen fine at each such $N$.
\formallabel{GlobalTheory.PolFan.exists_fine_invariant_fan}
\end{proposition}

\begin{proof}
The development first replaces $\varphi_s$ by a base function adapted to the hermitian forms over the
CM field, proves that it is polyhedral using reduction theory for balanced forms
(\lref{GlobalTheory.T.hasPolyhedralCover_balancedForms_CMField}, a polyhedral cover of the cone of
balanced hermitian forms) and finiteness of rational polyhedral cells up to automorphisms
(\lref{GlobalTheory.T.exists_finite_rational_polyhedral_cells_upto_aut}), and that linear forms on
balanced forms are dominated (\lref{GlobalTheory.T.dominatesLinearFormsOnBalancedForms_cm}). It then
takes an invariant smooth refinement without self-intersection (\cref{pf:lift1:smoothfan}), extends it homogeneously and
positively (\lref{GlobalTheory.T.exists_homogeneous_positive_polFan_extension}), and checks linearity
and fineness on its cones
(\lref{GlobalTheory.T.linearOn_isFine_of_linearOn_polarisationBaseFunction}). These steps are the combinatorics of
admissible cone decompositions for toroidal compactifications \cite{AshMumfordRapoportTai1975,FaltingsChai1990};
the refinement step is proved in \cref{pf:lift1:smoothfan}, the others are cited.
\end{proof}

\begin{proof}[Proof of \cref{pf:lift1:hodge-thm}]
If $r=0$, the datum has dimension $0$, so $\overline X$ is \'etale over $F$; being proper, it is finite
over $F$, hence affine. Take $\omega=\OO$: the weight $\mu_0$ has no entries, so $\mathcal L_{\mu_0}$ is
trivial and $\pi^*\omega$ has its class. Given $w$ and $A$, take $m=1$: on the affine (in particular
quasi-affine) scheme $\overline X$ the line bundle $A^{-1}$ has a section not vanishing at $j(w)$.

If $r\ge1$, $\overline X$ is smooth over $F$, so $\omega=\det\Omega$ is a line bundle
(\lref{GlobalTheory.FanTheta.isLB_detOm_of_smooth}). Its pull-back to the flag model is the Hodge
determinant of the flag model, and the latter is isomorphic to $\mathcal L_{\mu_0}$
(\lref{GlobalTheory.FlagHodge.lbIso_detOm_Llam}, which states this isomorphism). The domination
statement is \cref{pf:lift1:secdet}.
\end{proof}

